\documentclass[a4paper]{book}
\usepackage{bm}
\usepackage{tikz}
\usepackage{eucal}
\usepackage{tikz-cd}
\usepackage{fancyhdr}
\usepackage{mathrsfs}
\usepackage{graphicx}
\usepackage{enumerate}
\usepackage{mathtools}
\usepackage{circledsteps}
\usepackage{amsmath,amsthm,amssymb,amsfonts,stmaryrd}
\usepackage{hyperref}
\usepackage[textwidth=2cm,textsize=small,colorinlistoftodos]{todonotes}

\usepackage[nameinlink,capitalise,noabbrev]{cleveref}

%----------Removing chapter title from headers---------------------

\fancypagestyle{plain}{
    \lhead{}
    \fancyhead[R]{\thepage}
    \fancyhead[L]{}
    \renewcommand{\headrulewidth}{0pt}
    \fancyfoot{}
}

\pagestyle{fancy}
\fancyhead[R]{\thepage}
\fancyhead[L]{}
\renewcommand{\headrulewidth}{0pt}
\fancyfoot{}

%------------------------------------------------------------------

\newcommand{\lA}{\mathcal{A}}
\newcommand{\bA}{\mathbb{A}}
\newcommand{\lB}{\mathcal{B}}
\newcommand{\bB}{\mathbb{B}}
\newcommand{\lC}{\mathcal{C}}
\newcommand{\rC}{\mathscr{C}}
\newcommand{\kC}{\mathfrak{C}}
\newcommand{\bC}{\mathbb{C}}
\newcommand{\lD}{\mathcal{D}}
\newcommand{\lE}{\mathcal{E}}
\newcommand{\bE}{\mathbb{E}}
\newcommand{\lF}{\mathcal{F}}
\newcommand{\lG}{\mathcal{G}}
\newcommand{\lH}{\mathcal{H}}
\newcommand{\mH}{\mathrm{H}}
\newcommand{\lI}{\mathcal{I}}
\newcommand{\lJ}{\mathcal{J}}
\newcommand{\lK}{\mathcal{K}}
\newcommand{\lL}{\mathcal{L}}
\newcommand{\mL}{\mathrm{L}}
\newcommand{\lM}{\mathcal{M}}
\newcommand{\bN}{\mathbb{N}}
\newcommand{\mN}{\mathrm{N}}
\newcommand{\lN}{\mathcal{N}}
\newcommand{\fO}{\mathbf{O}}
\newcommand{\rO}{\mathscr{O}}
\newcommand{\lO}{\mathcal{O}}
\newcommand{\lP}{\mathcal{P}}
\newcommand{\bP}{\mathbb{P}}
\newcommand{\kP}{\mathfrak{P}}
\newcommand{\bQ}{\mathbb{Q}}
\newcommand{\bR}{\mathbb{R}}
\newcommand{\mR}{\mathrm{R}}
\newcommand{\lR}{\mathcal{R}}
\newcommand{\lS}{\mathcal{S}}
\newcommand{\bS}{\mathbb{S}}
\newcommand{\lT}{\mathcal{T}}
\newcommand{\lU}{\mathcal{U}}
\newcommand{\lV}{\mathcal{V}}
\newcommand{\lX}{\mathcal{X}}
\newcommand{\lY}{\mathcal{Y}}
\newcommand{\lW}{\mathcal{W}}
\newcommand{\bZ}{\mathbb{Z}}

\newcommand{\op}{\mathrm{op}} 
\newcommand{\im}{\mathrm{im}}
\newcommand{\id}{\mathrm{id}}
\newcommand{\set}{\lS\mathrm{et}} % category of sets
\newcommand{\poset}{\lP\lO\set} % category of partially ordered sets
\newcommand{\Grpd}{\lG\mathrm{rpd}} % category of groupoids
\newcommand{\cat}{\lC\mathrm{at}} % category of categories
\newcommand{\Fun}{\lF\mathrm{un}} % category of functors
\newcommand{\Psh}{\lP\Shv} % category of pre-sheaves
\newcommand{\Shv}{\lS\mathrm{h}} % category of sheaves
\newcommand{\hi}{\mathrm{hi}}
\newcommand{\Ho}{\mathrm{h}} % homotopy category
\newcommand{\sset}{\mathrm{s}\set} % category of simplicial sets
\newcommand{\Top}{\lT\op} % category of topological spaces 
\newcommand{\Kan}{\lK\mathrm{an}} % category of Kan complexes
\newcommand{\WKan}{\lW\Kan} % category of weak Kan complexes
\newcommand{\Open}{\lO\mathrm{pen}} % category of open subspaces
\newcommand{\Sing}{\lS\mathrm{ing}} % singular simplicial set
\newcommand{\scat}{\mathrm{s}\cat} % category of simplicial categories
\newcommand{\Path}{\mathrm{Path}}
\newcommand{\PSp}{\lP\Sp} % pre-spectra
\newcommand{\Sp}{\lS\mathrm{p}} % category of spectra
\newcommand{\st}{\mathrm{st}}
\newcommand{\Hom}{\lH\mathrm{om}} % enriched hom-spaces, hom-spectrum
\renewcommand{\hom}{\mathrm{Hom}} % hom-set, hom-space
\newcommand{\Mon}{\lM\mathrm{on}} % category of monoid objects
\newcommand{\CMon}{\lC\Mon} % category of commutative monoid objects
\newcommand{\Grp}{\lG\mathrm{rp}} % category of groups
\newcommand{\CGrp}{\lC\lG\mathrm{on}} % category of commutative group objects
\newcommand{\Ab}{\lA\mathrm{b}} % category of abelian groups
\newcommand{\Ch}{\lC\mathrm{h}} % category of chain complexes
\newcommand{\ch}{\mathrm{ch}}
\newcommand{\Bdl}{\lB\mathrm{dl}} % bundles
\renewcommand{\Line}{\lL\mathrm{ine}} % category of lines
\newcommand{\LineBdl}{\Line\Bdl} % category of line bundles
\newcommand{\GrbBdl}{\lG\mathrm{rb}\Bdl} % category of bundle gerbes
\newcommand{\spin}{\lS\mathrm{pin}} % spin (group)
\newcommand{\Euc}{\lE\mathrm{uc}} % site of Euclidean manifolds
\newcommand{\Mfd}{\lM\mathrm{fd}} % site of smooth manifolds
\newcommand{\sm}{\mathrm{sm}}
\newcommand{\Vect}{\lV\mathrm{ect}} % category of vector spaces
\newcommand{\Mod}{\lM\mathrm{od}} % category of modules
\newcommand{\Pic}{\lP\mathrm{ic}} % Picard groupoid
\newcommand{\cl}{\mathrm{cl}}
\newcommand{\pos}{\lP\mathrm{os}}
\newcommand{\Fred}{\lF\mathrm{red}}
\newcommand{\Ind}{\mathrm{Ind}} % Ind 
\newcommand{\PT}{\lP\lT} % Pontryagin-Thom collapse map

\DeclareMathOperator*\colim{colim} % colimit

% left adjoint on top of right adjoint: 
\newcommand{\adjun}[4]{
\begin{tikzcd}[column sep=2cm,ampersand replacement=\&]
  #1 \arrow[r,shift left,"#3"] \&
  #2 \arrow[l,shift left,"#4"] 
  \end{tikzcd}
}

\newcommand{\simpset}[7]{
 \begin{tikzcd}[row sep=0.5in, column sep=0.5in]
   #1 \arrow[r, shorten >=1ex,shorten <=1ex]
   \pgfmatrixnextcell #2 
   \arrow[l, shift left=1.2, "#5"] \arrow[l, shift right=1.2, "#4"'] 
   \arrow[r, shift right, shorten >=1ex,shorten <=1ex ] \arrow[r, shift left, shorten >=1ex,shorten <=1ex] 
   \pgfmatrixnextcell #3 
   \arrow[l] \arrow[l, shift left=2, "#7"] \arrow[l, shift right=2, "#6 "'] 
   \arrow[r, shorten >=1ex,shorten <=1ex] \arrow[r, shift left=2, shorten >=1ex,shorten <=1ex] \arrow[r, shift right=2, 
   shorten >=1ex,shorten <=1ex]
   \pgfmatrixnextcell \cdots 
   \arrow[l, shift right=1] \arrow[l, shift left=1] \arrow[l, shift right=3] \arrow[l, shift left=3] 
 \end{tikzcd}
}


% N.R. notes
\newcommand{\nrnote}[1]{\todo[color=green!40,linecolor=green!40!black,size=\tiny]{#1}}
\newcommand{\nrmpar}[1]{\todo[noline,color=green!40,linecolor=green!40!black,size=\tiny]{#1}}
\newcommand{\nrnoteil}[1]{\ \todo[inline,color=green!40,linecolor=green!40!black,size=\normalsize]{#1}}

\newtheorem{theorem}[equation]{Theorem}
\newtheorem{lemma}[equation]{Lemma}
\newtheorem{proposition}[equation]{Proposition}
\newtheorem{corollary}[equation]{Corollary}
% \newtheorem{statement}[section]{Statement}

\theoremstyle{definition}
\newtheorem{definition}[equation]{Definition}
\newtheorem{example}[equation]{Example}
% \newtheorem{attone}[equation]{Attention}

\theoremstyle{remark}
\newtheorem{remark}[equation]{Remark}
% \newtheorem{intone}[equation]{Intuition}
\newtheorem{notation}[equation]{Notation}
\newtheorem{question}[equation]{Question}
% \newtheorem{conjone}[equation]{Conjecture}
% \newtheorem{warning}[equation]{Warning}

\numberwithin{equation}{section}

\begin{document}

\tableofcontents

\chapter{Why Differential Cohomology? \texorpdfstring{\\}{ } \normalsize \emph{Talk by Konrad Waldorf}}


This talk is primarily an overview. We are trying to understand historically why we need differential cohomology as a generalization of the classical cohomology theory. Roughly speaking, the story goes through work of Deligne \cite{deligne1971mixedhodge}, Beilinson \cite{beilinson1984regulators}, and Brylinski \cite{brylinski1993geomquant}. It then proceeds with work of Cheeger--Simons \cite{cheegersimons1985diffchar}, Freed \cite{freed2000diffcoh} and Hopkins--Singer \cite{hopkinssinger2005diffcoh}. It makes significant advances with Simons--Sullivan \cite{simonssullivan2008diffcoh} and Bunke--Schick \cite{bunkeschick2009smoothk} and finally Bunke--Nikolaus-V{\"o}lkl \cite{bunkenikolausvoelkl2016diffcoh}, getting us to our modern definition.

\section{Why are Cohomology Theories not Enough?}
Consider the following fact about integral singular cohomology: Let $M$ be a smooth manifold, $\bC^*\Bdl(M)$ the category of $\bC^*$-principal bundles over $M$, and $\LineBdl$ the category of complex line bundles over $M$, then there is a sequence of isomorphisms:
\begin{equation}\label{eq:cech}
	\pi_0(\LineBdl(M))\simeq\pi_0(\bC^*\Bdl(M))\xrightarrow{\delta}\check{H}^2(M,\bZ)\simeq H^2(M,\bZ)
\end{equation}
\begin{remark}
	A short exact sequence of abelian sheaves $0\to\lF\to\lG\to\lH\to 0$ induces an exact sequence of \v Cech cohomology groups, see \cite{project2026cechsing}, including a connecting homomorphism $\delta:\check{H}^1(M,\lH)\to\check{H}^2(M,\lF)$. The morphism $\delta$ in \cref{eq:cech} is the connecting homomorphism associated to the sequence $0\to2\pi i\bZ\to\bC\to\bC^*\to0$. 
\end{remark}
In particular, isomorphism classes of line bundles only depend on the underlying topological space of $M$. To recover more information about the smooth geometry of $M$, let us consider the categories $\bC^*\Bdl_\nabla(M)$ and $\LineBdl_\nabla(M)$ of $\bC^*$-principal bundles, resp. line bundles, with connection and connection-preserving isomorphisms, then 
\[\pi_0(\LineBdl_\nabla(M))\simeq\pi_0(\bC^*\Bdl_\nabla(M))\]
\begin{question}
Is there a generalized version of cohomology theory $E^*$ such that $E^2(M)\simeq\pi_0(\bC^*\Bdl_\nabla(M))$? 
\end{question}
Imitating \v Cech and de Rham cohomology, consider the following modified de Rham chain complex of sheaves: Let $\Theta\in\Omega^1(\bC^*)$ be the Maurer-Cartan form, 
\begin{equation}\label{eq:delignecomp}
	\begin{tikzcd}
		\lD(n):C^\infty_{\bC^*}\arrow[r," d\log"] & \Omega^1 \arrow[r," d"] & \cdots \arrow[r," d"] & \Omega^{n-1}
	\end{tikzcd}
\end{equation}where $C^\infty_{\bC^*}$ is the sheaf of $\bC^*$-valued functions on $M$ and $ d\log$ maps $f\in C^{\infty}(U,\bC^*)$ to $f^*\Theta\in\Omega^1(U)$. Given a cover $U\to M$, construct $\lD(2)(U)$ as the total complex of the following bi-complex 
\begin{equation}\label{eq:bicomp}
	\begin{tikzcd}
		\vdots & \vdots & \vdots   \\
		C^\infty(U^{[3]},\bC^*) \arrow[r]\arrow[u] & \Omega^1(U^{[3]}) \arrow[r]\arrow[u] &  \Omega^2(U^{[3]}) \arrow[u] \arrow[r] & \cdots \\ 
		C^\infty(U^{[2]},\bC^*) \arrow[r," d^{1,1}"]\arrow[u,"\delta^{0,2}"] & \Omega^1(U^{[2]}) \arrow[u]\arrow[r] & \Omega^2(U^{[2]}) \arrow[u]\arrow[r] & \cdots\\
		C^\infty(U,\bC^*) \arrow[r]\arrow[u] & \Omega^1(U) \arrow[r," d^{2,0}"]\arrow[u,"\delta^{1,1}"] & \Omega^2(U) \arrow[u] \arrow[r] & \cdots
	\end{tikzcd}
\end{equation}where $U^{[n]}$ denotes the $n$-fold fiber product over $M$. Given $(L,\nabla)\in\bC^*\Bdl_\nabla(M)$, we construct a class in $C^\infty(U^{[2]},\bC^*)\oplus\Omega^1(U)$ as follows: Take $U=L\to M$ as cover, $\mu:L^{[2]}=L\times_ML\to\bC^*$ the difference map, characterized by $\mu(p_1,p_0)\triangleright p_0=p_1$, and $\nabla\in\Omega^1(L)$ the connection 1-form. Unfolding the definitions, notice the following:
\begin{enumerate}
	\item $\delta^{0,2}\mu=1$ is equivalent to the cocycle condition for $\mu$.
	\item $ d^{1,1}\mu=\delta^{1,1}\nabla$ is equivalent to $\nabla$ being $\bC^*$-equivariant.
	\item $ d^{2,0}\nabla=0$ appears only for $n>2$, and is equivalent to flatness of $\nabla$.  
\end{enumerate}
Notice that (1)-(2) are satisfied by definition, while (3) is an additional condition on $\nabla$.
\begin{definition}
	Let $M$ be a smooth manifold, denote by $\lH^n(M)$ the colimit over all covers $U\to M$ of the $n$-th cohomology group of $\lD(n)(U)$, called the \emph{$n$-th Deligne cohomology group}. 
\end{definition}
\begin{theorem}
	$\lH^2(M)\simeq\pi_0(\bC^*\Bdl_\nabla(M))$.
\end{theorem} 

\section{From Exact Sequences to the Hexagon}
A differential cohomology theory is meant as a coupling of ordinary cohomology classes with differential data, such as connections. The coherence conditions of this coupling are captured by axioms, in the style of Steenrod-Eilenberg axioms for ordinary cohomology theories, which include the existence of a certain diagram of natural transformations, called the \emph{hexagon}.

Taking a closer look at $\lH^n(M)$ we see the following: A $n$-cocycle representing a class in $\lH^n(M)$ consists of a cover $U\to M$, together with \[\mu:U^{[n]}\to\bC^*,\lA_1\in\Omega^1(U^{[n-1]}),\lA_2\in\Omega^2(U^{[n-2]}),\cdots,\lA_{n-1}\in\Omega^{n-1}(U)\]
The condition $\delta^{0,n}\mu=1$ is equivalent to $\mu$ being a $(n-1)$-cocycle for $\bC^*$-valued \v Cech cohomology of $M$, therefore there is a map $\lH^n(M)\to\check{H}^{n-1}(M,\bC^*)\simeq H^n(M,\bZ)$. Using \[\delta^{n-1,1}\lA_{n-1}= d^{n-1,1}\lA_{n-2}\]we can show that $\delta^{n,1} d^{n,0}\lA_{n-1}=0\in\Omega^n(U^{[2]})$, i.e. $ d^{n,0}\lA_{n-1}$ is constant on fibers of $\pi:U\to M$, so there is a unique $F\in\Omega^{n}(M)$ such that $\pi^*F= d^{n,0}\lA_{n-1}$ and is closed, therefore there is a map $\lH^n(M)\to\Omega^{n}_{\cl}(M)$.\nrnote{How to prove $F$ has integral periods?} The 
\begin{remark}\label{rmk:close}
	$F$ being closed can be proved as follows: $\pi^* d F= d^{n+1,0}\pi^*F= d^{n+1,0} d^{n,0}\lA_{n-1}=0$, by uniqueness we conclude $dF=0$. 
\end{remark}
\begin{remark}[\textcolor{red}{Achtung}] The differential $d^{n+1,0}$ does not appear in the complex \cref{eq:bicomp}, which stops at $d^{n,0}$ by definition, see \cref{eq:delignecomp}. The equation appearing in \cref{rmk:close} is interpreted in the bicomplex for $\lD(n+1)$. 
\end{remark}
\begin{remark}
	For $n=2$, we recover the usual notion of curvature 2-form. 
\end{remark}
For semplicity, consider $n=2$ and a $2$-cocycle of $\lH^2(M)$ such that $\mu$ is a \v Cech $2$-coboundary. This means that the underlying $\bC^*$-principal bundle is trivial, the choice of a section corresponds to the choice of $\beta\in C^\infty(U,\bC^*)$ such that $\delta^{0,1}\beta=\mu$. In that case, $(\mu,\lA)$ is cohomologous to $(0,\lA-d^{1,0}\beta)$ and identified by a unique $\rho\in\Omega^1(M)$ such that $\pi^*\rho=\lA-d^{1,0}\beta$. With a bit of work, one can prove that if $\lA-d^{1,0}\beta$ is exact, then $\rho$ is $d\log$-exact. In conclusion, there is a short exact sequence \begin{equation}
	\begin{tikzcd}
		0 \arrow[r] & \frac{\Omega^1(M)}{\mathrm{Im}d\log} \arrow[r] & \lH^2(M) \arrow[r] & H^2(M,\bZ) \arrow[r] & 0
	\end{tikzcd}
\end{equation}
\begin{remark}
	$\rho$ exists by the same argument as $F$ above, since $\delta^{1,1}(\lA-d^{1,0}\beta)=\delta^{1,1}\lA-d^{1,1}\delta^{0,1}\beta=\delta^{1,1}\lA-d^{1,1}\mu=0$ in $\Omega^1(U^{[2]})$. Moreover, $\pi^*d\rho=d^{2,0}\lA$, so $\rho$ is closed if and only if $\lA$ is flat, i.e. $d^{2,0}\lA=0$. 
\end{remark}
\begin{remark}
	Surjectivity is a corollary of the fact that every line bundle admits a connection, proven using partitions of unity. 
\end{remark}
Similarly, consider the curvature map $\lH^2(M)\to\Omega^{2}_\cl(M)$ and $\lA\in\Omega^1(U)$ a horizontal boundary, i.e. $\beta\in C^\infty(U,\bC^*)$ exists such that $d^{1,0}\beta=\lA$, then $\mu-\delta^{0,1}\beta$ is a 1-cocycle for \v Cech cohomology of $M$. In conclusion there is a short exact sequence \begin{equation}
	\begin{tikzcd}
		0 \arrow[r] & \check{H}^1(M,\bC^*) \arrow[r] & \lH^2(M) \arrow[r] & \Omega^2_{\cl,\bZ}(M) \arrow[r] & 0 
	\end{tikzcd}
\end{equation}where $\Omega^n_{\cl,\bZ}(M)$ denotes the sub-space of closed, integral $n$-forms.
\begin{remark}
	A closed $n$-form is integral if the corresponding cohomology class in real cohomology, via the de Rham isomorphism, belong to the image of $H^n(M,\bZ)\to H^n(M,\bR)$. 
\end{remark}
%----------TO BE CONTINUED------------%
These exact sequences are not independent from one another, but are linked via the de Rham differential $d:\Omega^1(M)\to\Omega^2(M)$ and the connecting homomorphism $H^n(M,\bC^*)\to H^{n+1}(M,\bZ)$. The diagram incorporating all these maps is:
\begin{equation}
	\begin{tikzcd}[column sep={2cm,between origins},row sep={1.5cm,between origins}]
            & H^{n-1}(M,\bC^*) \arrow[dr]\arrow[rr] & & H^n(M,\bZ) \arrow[dr] & & \\
            H^{n-1}(M,\bC)\arrow[ru]\arrow[rd] & & \lH^n(M) \arrow[dr]\arrow[ru] & & H^n(M,\bC) & \\
            & \frac{\Omega^{n-1}(M)}{\mathrm{Im}d_{\lD(n)}} \arrow[ru]\arrow[rr,"d_{\lD(n)}"] & & \Omega^n_{\cl,\bZ}(M) \arrow[ru] & & 
        \end{tikzcd}
\end{equation}
where $d_{\lD(n)}$ is the differential for $\lD(n)$, to distinguish from the de Rham differential. This diagram was first discovered by Simons and Sullivan in the 2000s \cite{simonssullivan2008diffcoh}.

\section{Holonomy and Characters}
We have seen that differential cohomology is a generalization of ordinary cohomology. We can hence wonder how differential cohomology interacts with other aspects of ordinary cohomology theory, such as holonomy and characters. Holonomy of a line bundle with connection is obtained by exponential of the integral of a connection 1-form over loops, we wish to generalize this from loops to maps with a generic orientable $(n-1)$-dimensional as domain. Let $f: \Sigma \to M$ be a smooth map, where $\Sigma$ is an closed $(n-1)$-manifold (i.e. $\Sigma$ is connected, compact, and without boundary), and $\mu\in\lH^n(M)$ be a class in Deligne cohomology. Since $H^n(\Sigma,\bZ)\simeq0$, for dimensional reasons, $f^*(\mu)\in\lH^n(\Sigma)$ must be in the image of \begin{equation}\label{eq:inclus}
	\begin{tikzcd}
		\frac{\Omega^1(\Sigma)}{\mathrm{Im}d_{\lD(n)}} \arrow[r] & \lH^n(\Sigma)
	\end{tikzcd}
\end{equation}Pick $\omega\in\Omega^1(\Sigma)$ in the equivalence class that is mapped to $f^*(\mu)$ via \cref{eq:inclus}. 
%---------------TO BE CONTINUED--------------------
\[Hol_\xi(\Phi) = e^{2\pi i\int_{\Sigma}\varphi} \in \bC^\times\] 
Here the exponential is crucial as the integral is only well-defined up to integer multiples of $2\pi i$. Indeed, if $\varphi, \varphi'$ are in the fiber, then $\int_\Sigma \varphi - \int_\Sigma \varphi' \in \bZ$\nrnote{Have to check with Konrad, but if $\Sigma$ is without boundary, then there should be no issue.}. We can now apply Stokes' theorem to this situation. Concretely, if $\Sigma = \partial B$, where $B$ is an $n$-dimensional manifold and $\phi = \Phi|_\Sigma$, then we have

\[ Hol(\phi) = e^{2\pi i \int_B \Phi^*\mathrm{curv}(\xi)}\]

Let us look at a manifestation thereof. In the WZW (Wess–Zumino–Witten) Model one has a surface $\Sigma$, $\phi\colon \Sigma \to M$, and a $B$-field $B \in \Omega^2(M)$ . Then we have $S(\phi) = \int_\Sigma \phi^*B$. \nrnote{How does that relate to the previous explanation?}  

If $M = G$ is a Lie group, then WZW model is not globally defined i.e. $dB = H = <\theta \wedge [\theta \wedge \theta]>$. The solution is to take a differential cohomology class $\xi \in \hat{H}^3(G)$, such that  $dB = \mathrm{curv}(\xi)$. Then $S(\phi) = Hol_\xi(\phi)$. This indeed matches our intuition, as $Hol_{i(\alpha)}(\phi) = e^{2\pi i \int_\Sigma \phi^*\alpha}$. Witten showed that if $G$ is a simply connected Lie group, then the setup admits an extension to $3$ dimensions (we can skip dimension $2$), which implies that every map $\Sigma \to G$ must be nul-homotopic. Finally, let $[\Sigma]$ be the fundamental class of $\Sigma$. Then $\phi_*[\Sigma] \in H_{n-1}(M)$ and $\xi\colon H_{n-1} \to \bC^\times$, reconstructs differential cohomology classes.

\chapter{The Underlying Machine of Higher Category Theory \texorpdfstring{\\}{ } \normalsize \emph{Talk by Matthias Frerichs}}

In this lecture we introduce the machinery of higher category theory needed for those talks that come after. For the $\infty$-categorical background, we mainly follow \cite{lurie2009htt} and \cite{land2021introduction}. 

\section{Background}

\subsection{Model Categories}

We recall some definitions regarding model categories. For more details, see \cite{hirschhorn2003modelcategories} (we will also refer to chapter 11 of \cite{riehl2014categoricalhomotopytheory}). A \emph{model category} consists of a complete and cocomplete category $\lC$ together with three classes of maps, called \emph{weak equivalences, fibrations}, and \emph{cofibrations}. The standard axioms (\cite[Definition 7.1.3]{hirschhorn2003modelcategories}) on these data are the following:  
\begin{enumerate}
	\item Weak equivalences contain all isomorphisms and satifying $2003$ (2-out-of-3), i.e., for every pair $f,g$ of composable morphisms, if any two of $f,g$ and $gf$ are weak equivalences, so is the third.
	\item All classes are closed under retracts, i.e., $f$ is a retract of $g$ (\cite[Definition 7.1.1]{hirschhorn2003modelcategories}) and $g$ is a weak equivalence, fibration, or cofibration, the so is $f$.
	\item For every solid diagram \begin{equation}\label{diag:lifting}
    \begin{tikzcd}
      X \arrow[d,"i"']\arrow[r] & Y\arrow[d,"f"]  \\
      Z\arrow[ru,dashed]\arrow[r] & W
    \end{tikzcd}
  \end{equation}there is a dashed arrow making the diagram commute if $i$, resp. $f$, is a cofibration, resp. fibration, and either one is a weak equivalence. 
  \item Every morphism admits a factorization consisting of a cofibrantion followed by an acyclic fibration, and a factorization consisting of an acyclic cofibration followed by a fibration. 
\end{enumerate}
\begin{definition}
	Given a pair of morphisms $i,f$, we say $f$ is \emph{right orthogonal to $i$} and $i$ is \emph{left orthogonal to $f$}, if for every solid diagram \begin{equation}
    \begin{tikzcd}
      X \arrow[d,"i"']\arrow[r] & Y\arrow[d,"f"]  \\
      Z\arrow[ru,dashed]\arrow[r] & W
    \end{tikzcd}
  \end{equation}there is a dashed arrow making the diagram commute
\end{definition}
The last three axioms are equivalent to saying that acyclic cofibration with fibrations, and cofibrations with acyclic fibrations, form a \emph{weak factorization system} (\cite[Definition 11.2.1]{riehl2014categoricalhomotopytheory}), see \cite[\S 11.3]{riehl2014categoricalhomotopytheory}. 

Denote by $\emptyset$, resp. $1$, the initial, resp. terminal, object. An object $X$ is \emph{cofibrant}, resp. \emph{fibrant}, if the canonical morphism $\emptyset\to X$, resp. $X\to1$, is a cofibration, resp. fibration. An object both cofibrant and fibrant is called \emph{bifibrant}. The full sub-category of bifibrant objects is denoted as $\lC^\circ$. 

\subsection{Simplicial Sets}

We introduce some notations, definitions, and constructions regarding simplicial sets. For more details, see \cite{hovey1999modelcategories,goerssjardine1999simplicialhomotopytheory}, or \cite[Tag 0004]{kerodon2026}. For each natural number $n$, denote by $[n]$ the linearly ordered set $\{0,\cdots,n\}$. Let $\Delta$ denote the category of linearly ordered sets $[n]$ and (non-strict) monotone functions between them. Functors $\Delta^\op\to\set$ are called \emph{simplicial sets}. The category of simplicial sets is denoted as $\sset$. 

Simplicial sets are connected to topological spaces via the singular simplicial set and geometric realization functors. Given $n$, denote by $|\Delta^n|$ the sub-space $\subseteq[0,1]^{n+1}$ of tuples $(t_0,\cdots,t_n)$ such that $t_0+\cdots+t_n=1$. Given a topological space $X$, denote by $\Sing X$ its \emph{singular simplicial set}. The nerve functor admits a left adjoint, sending a simplicial set $S$ to its \emph{geometric realization} $|S|$.

On the other hand, simplicial sets are connected to categories via the nerve functor. Given $n$, denote by $[n]$ the category associated to the linearly ordered set $[n]$. Given a category $\lC$, denote by $\mN\lC$ its \emph{nerve}. The nerve functor admits a left adjoint, sending a simplicial set $S$ to its \emph{homotopy category} $\pi(S)$, see \cite[Definition 1.2.5]{land2021introduction} or \cite[Tag 004M]{kerodon2026}. As we'll see, the homotopy category admits a simpler description for a special class of simplicial sets. 

Among the morphisms of $\Delta$, there are two special classes: 
\begin{enumerate}
	\item Given $0\leq i\leq n$, we denote by $d^i_n$ the unique injective, monotone function $[n-1]\to[n]$ such that $i$ is not contained in the image. Given a simplicial object $S$, we'll write $d_i$ for the map $S_n\to S_{n-1}$, called \emph{face map}, induced by $d^i_n$ (leaving implicit the dependency on $S$ and $n$).
	\item Given $0\leq i\leq n$, we denote by $s^i_n$ the unique surjective, monotone function $[n+1]\to[n]$ such that $i,i+1$ are mapped to $i$. Given a simplicial object $S$, we'll write $s_i$ for the map $S_{n+1}\to S_n$, called \emph{degeneracy map}, induced by $s^i_n$ (leaving implicit the dependency on $S$ and $n$). 
\end{enumerate}The sets $S_n$, for every $n\geq0$, together with face and degeneracy maps, completely determine a simplicial set, see \cite[Tag 04FW]{kerodon2026} (see also \cite[\S VII.5, Proposition 2]{maclane1998categories}, for a vintage reference). 
\begin{remark}
	The category $\sset$ is closed cartesian, hence self-enriched. We use the exponential notation for the internal homs of $\sset$, i.e., given simplicial sets $X,Y$, we write $Y^X$ for the simplicial set of morphisms from $X$ to $Y$. 
\end{remark}
\begin{remark}
	Given a partially ordered set $J$, denote by $\Delta^J$ the simplicial set represented by $J$, i.e., a $n$-simplex consists of a monotone function $[n]\to J$.  
\end{remark}
\begin{remark}
Given $n$, a simplicial set $S$, and a $0$-simplex $X$, we'll abuse notation and denote by $X$ the image of $X$ under the map $S_0\to S_n$ induced by the terminal function $[n]\to[0]$. Given a 1-simplex $f$, we'll write $f:X\to Y$ to mean $X$, resp. $Y$, is the 1-face, resp. 0-face, of $f$.
\end{remark}
\begin{definition}
	Given $0\leq i\leq n$, denote by $\Lambda^n_i$ the \emph{$i$-th horn}, characterized as follows: A $k$-simplex of $\Lambda^n_i$ is a monotone function $f:[k]\to[n]$ such that $\{i\}\cup f([k])\neq[n]$. Horn corresponding to $0\leq i<n$, resp. $0<i\leq n$, are called \emph{left}, resp. \emph{right}. Horns that are both left and right are called \emph{inner}. 
\end{definition}
\begin{definition}
	A simplicial set $S$ is called a \emph{Kan complex} if the following condition holds: For every $0\leq i\leq n$ and solid diagram \begin{equation}
    \begin{tikzcd}
      \Lambda^n_i \arrow[d,hook]\arrow[r] & S  \\
      \Delta^n \arrow[ru,dashed]
    \end{tikzcd}
  \end{equation}there is a dashed arrow making the diagram commute. Denote by $\Kan$ the full sub-category of $\sset$ consisting of Kan complexes. 
\end{definition}

% \begin{definition}
%	Given $n$, denote by $\partial\Delta^n$ the \emph{$n$-boundary}, characterized as follows: A $k$-simplex of $\partial\Delta^n$ is a monotone function $f:[k]\to[n]$ such that $f([k])\neq[n]$. 
% \end{definition}
% \begin{definition}
%	A morphism of simplicial sets is called a \emph{weak homotopy equivalence} if its geometric realization is a weak homotopy equivalence of topological spaces. 
% \end{definition}
% The geometric realization of $\Lambda^n_i\subseteq\Delta^n$ is (up to homeomorphism, see \cite[Tag 050E]{kerodon2026}) a deformation retract (see \cite[Tag 002K]{kerodon2026}), therefore horn inclusions are weak homotopy equivalences.  

\section{Models of Higher Category Theory}

\subsection{\texorpdfstring{$\infty$}{oo}-Categories} 

\begin{definition}
	A simplicial set $S$ is called a \emph{$\infty$-category} if the following condition holds: For every $0< i< n$ and solid diagram \begin{equation}
    \begin{tikzcd}\label{eq:diag}
      \Lambda^n_i \arrow[d,hook]\arrow[r] & S  \\
      \Delta^n \arrow[ru,dashed]
    \end{tikzcd}
  \end{equation}there is a dashed arrow making the diagram commute. Denote by $\WKan$ the full sub-category of $\sset$ consisting of $\infty$-categories. 
\end{definition}
\begin{example}
  The nerve of a category is an $\infty$-category. Conversely, an $\infty$-category is isomorphic to the nerve of a category if and only if, for every solid diagram like (\ref{eq:diag}), there is a \emph{unique} dashed arrow making the diagram commute (see \cite[Proposition 1.1.2.2]{lurie2009htt} or \cite[Theorem 1.1.52]{land2021introduction}). 
\end{example}
\begin{definition}
	Given a simplicial set $S$ and 1-simplices $f:X\to Y,g:Y\to Z,h:X\to Z$, we say \emph{$h$ is a composition of $f$ and $g$} if there is a 2-simplex with boundary 
\begin{equation}
    \begin{tikzcd}
      & Y \arrow[dr,"g"] & \\
      X \arrow[ru,"f"] \arrow[rr,"h"'] & {} & Z 
    \end{tikzcd}
\end{equation}Given 1-simplices $f,h:X\to Y$, we say \emph{$h$ is homotopic to $f$} if $h$ is a composition of $f$ and $Y$ (viewed as a 1-simplex).  
\end{definition}
We recall briefly the construction of the category $\Ho S$ of a $\infty$-category $S$ (details can be found in \cite{lurie2009htt}, \cite[Tag 0049]{kerodon2026}, or \cite[Definition 1.2.5]{land2021introduction}):
\begin{enumerate}
	\item Objects of $\Ho S$ are the 0-simplices of $S$.
	\item Given 0-simplices $X,Y$, being homotopic is an equivalence relation on 1-simplices from $X$ to $Y$ (\cite[Tag 003Z]{kerodon2026}). Denote by $\hom_{\Ho S}(X,Y)$ the homotopy classes of 1-simplices from $X$ to $Y$.
	\item Given 1-simplices $f:X\to Y,g:Y\to Z$, there is a 1-simplex $h:X\to Z$, unique up to homotopy, such that $h$ is a composition of $f$ and $g$ (\cite[Tag 0043]{kerodon2026}). Moreover, the homotopy class of $h$ only depends on the homotopy classes of $f$ and $g$ (\cite[Tag 0048]{kerodon2026}). Therefore, there is a unique function \begin{equation}
		\begin{tikzcd}[column sep=1.5cm]
			\hom_{\Ho S}(X,Y)\times\hom_{\Ho S}(Y,Z) \arrow[r,"\circ_{X,Y,Z}"] & \hom_{\Ho S}(X,Z)
		\end{tikzcd}
	\end{equation}  such that $[h]=[g]\circ_{X,Y,Z}[f]$ whenever $h$ is a composition of $f$ and $g$. 
	\item The composition laws $\circ_{X,Y,Z}$ are associative and, given a 0-simplex $X$, the homotopy class of $X$ (viewed as a 1-simplex) acts as the identity morphism of $X$ (\cite[Tag 004B]{kerodon2026}).
\end{enumerate}The following is a special case of \cite[Corollary 1.2.13]{land2021introduction}. 
\begin{theorem}Given an $\infty$-category $S$, the category $\Ho S$ is isomorphic to $\pi(S)$.
\end{theorem}We will also refer to $\Ho S$ as the \emph{homotopy category} of $S$. 
\begin{remark}
	We switch to using calligraphic upper-case letters to denote $\infty$-categories. Given a $\infty$-category $\lC$, we refer to 0-simplices, resp. 1-simplices, as \emph{objects}, resp. \emph{morphisms}. 
\end{remark}
\begin{definition}
	Given a $\infty$-category $\lC$, a morphism is called an \emph{equivalence} if its image in $\Ho\lC$ is an isomorphism. 
\end{definition}
\begin{definition}
	A $\infty$-category $\lC$ is called an \emph{$\infty$-groupoid} if $\Ho\lC$ is a groupoid, i.e., every morphism is an equivalence. 
\end{definition}
The following is a result of A. Joyal.
\begin{theorem} A $\infty$-category $\lC$ is a Kan complex if and only if it is a $\infty$-groupoid. 
\end{theorem}
For the proof, see \cite[Corollary 1.4]{joyal2002quasicatkancomplexes}, \cite[Tag 019D]{kerodon2026}, \cite[Proposition 1.2.5.1]{lurie2009htt}, or \cite[Corollary 2.1.12]{land2021introduction}.
\begin{theorem}
	Given a simplicial set $S$ and $\infty$-category $\lC$, the hom simplicial set $\lC^S$ is an $\infty$-category. 
\end{theorem}
For the proof, see \cite[Corollary 1.3.38]{land2021introduction}, \cite[Tag 0066]{kerodon2026}, or \cite[Corollary 2.3.2.5]{lurie2009htt}.
\begin{definition}
	Given $\infty$-categories $\lC,\lD$, denote by $\Fun(\lC,\lD)$ the simplicial set $\lD^\lC$, called the \emph{$\infty$-category of functors from $\lC$ to $\lD$}.  
\end{definition}
\begin{definition}
	Given a simplicial set $S$ and simplices $X,Y$, denote by $\hom_S(X,Y)$ the pullback of the following diagram: \begin{equation}
		\begin{tikzcd}
			& S^{\Delta^1} \arrow[d,"{(d_1,d_0)}"]\\
			\Delta^0 \arrow[r,"{(X,Y)}"'] & S\times S
		\end{tikzcd}
	\end{equation}
\end{definition}
\begin{theorem}
	Given an $\infty$-category $\lC$ and objects $X,Y$, the simplicial set $\hom_\lC(X,Y)$ is an $\infty$-groupoid.
\end{theorem}
For the proof, see \cite[Corollary 2.2.4]{land2021introduction}, or \cite[Tag 01P1]{kerodon2026}. 
\begin{definition}
	Given an $\infty$-category $\lC$ and objects $X,Y$, we call $\hom_\lC(X,Y)$ the \emph{mapping space from $X$ to $Y$}.
\end{definition}
\begin{definition}
	Given $\infty$-categories $\lC,\lD$, a functor $F:\lC\to\lD$ is called an \emph{Joyal equivalence} if there is $G:\lD\to\lC$ such that $F\circ G$, resp. $G\circ F$, is equivalent to the identity of $\lD$, resp. $\lC$. 
\end{definition}
Given a simplicial set $S$, denote by $S^\op$ the simplicial set having, for every $n$, $S^\op_n=S_n$, and for every $0\leq i\leq n$, \begin{equation}
		\begin{tikzcd}[row sep=0.25cm]
		(S^\op_n \arrow[r,"d_i^\op"] &  S^\op_{n-1}):=(S_n\arrow[r,"d_{n-i}"] &  S_{n-1})\\
		(S^\op_n \arrow[r,"s_i^\op"] &  S^\op_{n+1}):=(S_n\arrow[r,"s_{n-i}"] &  S_{n+1})
		\end{tikzcd}
	\end{equation}
\begin{theorem}[{\cite[Tag 003S]{kerodon2026}}]
Given an $\infty$-category $\lC$, the simplicial set $\lC^\op$ is an $\infty$-category. 
\end{theorem} 
\begin{definition}
	Given an $\infty$-category $\lC$, we call $\lC^\op$ the \emph{opposite $\infty$-category of $\lC$}. 
\end{definition}
Given a linearly ordered set $I$, a \emph{cut} consists of a partition $J,J'$ of $I$ such that $j<j'$, for every pair $j\in J$ and $j'\in J'$. The cuts $(\emptyset,I)$ and $(I,\emptyset)$ are allowed. Notice that a cut of $I$ is equivalent to a monotone function $I\to[1]$. 
\begin{definition}
	Given a pair of simplicial sets $X,Y$, denote by $X\star Y$ the \emph{join of $X$ and $Y$}, characterized as follows: A $n$-simplex of $X\star Y$ is a triple consisting of a cut $(J,J')$ of $[n]$ and morphisms of simplicial sets $\Delta^{J}\to X$ and $\Delta^{J'}\to Y$. Given a monotone function $p:[k]\to[n]$, the corresponding function on simplices sends a triple \begin{equation}
		\begin{tikzcd}
			(J,J') & \Delta^{J} \arrow[r,"\sigma"] & X & \Delta^{J'} \arrow[r,"\sigma'"] & Y
		\end{tikzcd}
	\end{equation}to
	\begin{equation}
		\begin{tikzcd}
			(p^{-1}(J),p^{-1}(J')) & \Delta^{p^{-1}(J)} \arrow[d]\arrow[r] & X & \Delta^{p^{-1}(J')} \arrow[d]\arrow[r] & Y\\
			& \Delta^{J} \arrow[ru,"\sigma"'] & & \Delta^{J'} \arrow[ru,"\sigma'"']
		\end{tikzcd}
	\end{equation}where the vertical arrows are induced by the restriction of $p$ to $p^{-1}(J)$ and $p^{-1}(J')$. 
\end{definition}
It can be checked directly that $\star$ underlies a monoidal structure on $\sset$, with monoidal unit $\emptyset$ (see \cite[Tag 016K]{kerodon2026}). Given a simplicial set $X$, denote by $\emptyset_X$ the initial functor $\emptyset\to X$
\begin{definition}
	Given a morphism of simplicial sets $\alpha:S\to X$, denote by $X_{\alpha/}$, resp. $X_{/\alpha}$, the simplicial set characterized as follows: A $n$-simplex of $X_{\alpha/}$, resp. $X_{/\alpha}$, is a commutative square \begin{equation}
		\begin{tikzcd}[row sep=1cm]
			S\star\emptyset\arrow[r,"\cong"]\arrow[d,"S\star\emptyset_{\Delta^n}"'] & S \arrow[d,"\alpha"]  \\
			S\star\Delta^n \arrow[r] & X
		\end{tikzcd}\text{, resp.}\qquad
		\begin{tikzcd}[row sep=1cm]
			\emptyset\star S \arrow[r,"\cong"]\arrow[d,"\emptyset_{\Delta^n}\star S"'] & S \arrow[d,"\alpha"] \\
			\Delta^n\star S \arrow[r] & X
		\end{tikzcd}
	\end{equation}
\end{definition}
\begin{remark}
	Notice that $X_{\alpha/}$ has a canonical projection to $X$, sending a $n$-simplex to the following composition \begin{equation}
		\begin{tikzcd}[column sep=1cm]
			\Delta^n \arrow[r,"\cong"] & \emptyset\star\Delta^n \arrow[r,"\emptyset_S\star\Delta^n"] & S\star\Delta^n \arrow[r] & X
		\end{tikzcd}
	\end{equation}A similar projection exists for $X_{/\alpha}$. 
\end{remark}
\begin{theorem}
	Given a $\infty$-category $\lC$ and morphism of simplicial sets $\alpha:S\to\lC$, the simplicial sets $\lC_{\alpha/}$ and $\lC_{/\alpha}$ are $\infty$-categories. 
\end{theorem}For the proof, see \cite[Corollary 2.1.2.2]{lurie2009htt}, or \cite[Corollary 1.4.24]{land2021introduction}.
\begin{definition}
	Given a $\infty$-category $\lC$ and morphism of simplicial sets $\alpha:S\to\lC$, we call $\lC_{\alpha/}$, resp. $\lC_{/\alpha}$, the \emph{$\infty$-category of objects under}, resp. \emph{over}, $\alpha$. 
\end{definition}

\subsection{Simplicial Categories}

\begin{definition}
	Categories enriched over $\sset$ (with cartesian monoidal structure) are called \emph{simplicial categories}. Denote by $\cat_{\sset}$ the category of simplicial categories. 
\end{definition}
\begin{remark}
	The functor $\delta:\set\to\sset$ sending a set $S$ to $\coprod_{S}\Delta^0$ is fully faithful onto discrete simplicial sets. Since $\delta$ is strong monoidal, it allows us to identify locally small categories with simplicial categories having discrete hom simplicial sets. 
\end{remark}
\begin{remark}
	The functor $\pi_0:\sset\to\set$ is strong symmetric monoidal (see \cite[Lemma 1.2.47]{land2021introduction} or \cite[Tag 00GX]{kerodon2026}), hence it induces a change of enrichment $\pi_*:\cat_{\sset}\to\cat$. Given a simplicial category $\lC$, there is a natural functor $\lC\to\pi_*\lC$.
\end{remark}
\begin{definition}
	A functor of $\infty$-categories is called a \emph{categorical equivalence} if it is an isomorphism in $\pi_*\WKan$. A morphism of simplicial sets $f:X\to S$ is called a \emph{categorical equivalence} if, for every $\infty$-category $\lC$, the induced functor \begin{equation}
		\begin{tikzcd}
			\Fun(S,\lC) \arrow[r,"f^*"] & \Fun(X,\lC)
		\end{tikzcd}
	\end{equation}is a categorical equivalence. 
\end{definition}
We recall briefly the construction of the \emph{simplicial nerve} of a simplicial category $\lC$ (details can be found in \cite[\S 1.1.5]{lurie2009htt} and \cite[Tag 00KM]{kerodon2026}). Given a partial ordered set $J$, define the simplicial category $\kC^J$ as follows: 
\begin{enumerate}
	\item Objects of $\kC^J$ are the elements of $J$. 
	\item Given a pair of objects $i,j$, let $\hom_{\kC^J}(i,j)$ be the nerve of the (category associated to the) partially ordered set $\kP_{i,j}$ of $I\subseteq\{k\in J|i\leq k\leq j\}$ containing $\{i,j\}$, ordered by reversed inclusion $\supseteq$.
	\item Given a triple of objects $i,j,k$, composition is defined as the nerve applied to the monotone map \begin{equation}
		\begin{tikzcd}[column sep=1.5cm]
			\kP_{i,j}\times\kP_{j,k} \arrow[r,"\cup"] & \kP_{i,k}
		\end{tikzcd}
	\end{equation}
	\item Given an object $i$, the singleton $\{i\}$ acts as the identity. 
\end{enumerate}
Given a simplicial category $\lC$, denote by $\mN\lC$ its \emph{simplicial nerve}, the $n$-simplices of which are enriched functors $\kC[\Delta^n]\to\lC$. 
\begin{remark}
	The simplicial nerve of a category is isomorphic to the ordinary nerve, see \cite[Tag 00KZ]{kerodon2026} or \cite[Lemma 1.2.64]{land2021introduction}.
\end{remark}
\begin{theorem} Given a simplicial set $S$ and a Kan complex $\lC$, the hom simplicial set $\lC^S$ is a Kan complex. 
\end{theorem}For the proof, see {\cite[Corollary 1.3.38]{land2021introduction}}. In particular, $\Kan$ is self-enriched. The following is a result of Cordier and Porter.
\begin{theorem}
	Given a category $\lC$ enriched over Kan complex, the simplicial nerve $\mN\lC$ is an $\infty$-category. 
\end{theorem}
For the proof, see \cite[Lemma 1.2.70]{land2021introduction}, \cite[Proposition 1.1.5.10]{lurie2009htt}, or \cite[Tag 00LJ]{kerodon2026}.
\begin{theorem}
	Given a simplicial category $\lC$, the nerve of $\lC\to\pi_*\lC$ factors through an isomorphism $\pi(\mN\lC)\to\pi_*\lC$.
\end{theorem}For the proof (if $\lC$ is enriched over Kan complexes), see {\cite[Tag 00M4]{kerodon2026}}. The general case is \cite[Exercise 50]{land2021introduction}, which can be solved using the same argument as \cite[Tag 00M4]{kerodon2026}, as a consequence of \cite[Observation 1.2.8 and Lemma 1.2.70]{land2021introduction}. 
\begin{definition}
	Denote by $\lS$ the $\infty$-category $\mN\Kan$, called the \emph{$\infty$-category of $\infty$-groupoids}.
\end{definition}
\begin{definition}[{\cite[Definition 5.1.0.1]{lurie2009htt}}]
	Given a $\infty$-category $\lC$, denote by $\Psh(\lC)$ the $\infty$-category of functors $\lC^\op\to\lS$, called the \emph{$\infty$-category of pre-sheaves on $\lC$}. 
\end{definition}

% To Do: Yoneda embedding. Straightening and Unstraightening. Twisted arrow category. 

\subsection{Group Actions on Spaces}

The following is a result of Moore. 
\begin{theorem}
	The underlying simplicial set of a simplicial group is a Kan complex. 
\end{theorem}For a proof, see \cite[Ch. 1, Lemma 3.4]{goerssjardine1999simplicialhomotopytheory}. In particular, given a simplicial groupoid $\lC$ with a single object, $\mN\lC$ is a $\infty$-groupoid.


% Definition of G-torsor

% Simplicial groups are Kan complexes, the corresponding oo-category with one object is a oo-groupoid and is the delooping of the group, i.e., equivalent to the W-bar construction of G. Given a simplicial group G, there is a canonical equivalence between cofibrant-fibrant G-simplicial sets (projective model structure) and functors BG->S. The quotient functor -/G is then equivalent to the colimit functor Fun(BG,S)->S. 

\section{Accessible and Presentable \texorpdfstring{$\infty$}{oo}-Categories}

% Bla Bla, adjoint functor theorem, presentable categories, bla bla...

\begin{definition}
	Given a cardinal $\kappa$, a simplicial set is called \emph{$\kappa$-small} if the set of all its non-degenerate simplices has cardinality $<\kappa$. A $\infty$-category is called $\kappa$-small if its underlying simplicial set is. 
\end{definition}
Let $\kappa$ be a regular cardinal. 
\begin{definition}
	A $\infty$-category $\lC$ is \emph{$\kappa$-filtered} if every diagram indexed by a $\kappa$-small simplicial set has a cocone. The colimit of a diagram indexed by a small $\kappa$-filtered $\infty$-category is called a \emph{$\kappa$-filtered colimit}. 
\end{definition}
\begin{definition}
	A functor of $\infty$-categories $F:\lC\to\lD$ is called \emph{$\kappa$-accessible} if it preserves all $\kappa$-filtered colimits. Denote by $\Fun^\kappa(\lC,\lD)$ the full sub-category of $\Fun(\lC,\lD)$ consisting $\kappa$-accessible functors.
\end{definition}
\begin{remark}
	$\kappa$-accessible functors are called \emph{$\kappa$-finitary functors} (\cite[Tag 05Y6]{kerodon2026}). 
\end{remark}
\begin{definition}
	Given an $\infty$-category $\lC$, an object $X$ is called \emph{$\kappa$-compact} if $\hom_\lC(X,-)$ is $\kappa$-accessible. Denote by $\lC^\kappa$ the full sub-category of $\lC$ consisting of $\kappa$-compact objects. 
\end{definition}
\begin{definition}
	A $\infty$-category $\lC$ is called \emph{$\kappa$-compactly generated} if the following hold: 
	\begin{enumerate}
		\item $\lC$ admits all small $\kappa$-filtered colimits. 
		\item Every objects can be written as a small $\kappa$-filtered colimit of $\kappa$-compact objects. 
	\end{enumerate}
\end{definition}
\begin{definition}
	Given  a functor of $\infty$-categories $F:\lC\to\lD$, we say \emph{$F$ exhibits $\lD$ as the $\Ind_\kappa$-completion of $\lC$} if the following hold:
	\begin{enumerate}
		\item $\lD$ admits all small $\kappa$-filtered colimits.
		\item For every $\infty$-category $\lX$ that admits all small $\kappa$-filtered colimit, pre-composition with $F$ induces an equivalence \begin{equation}
			\begin{tikzcd}
				\Fun^\kappa(\lD,\lX) \arrow[r,"F^*"] & \Fun(\lC,\lX)
			\end{tikzcd}
		\end{equation}
	\end{enumerate}
	Given a $\kappa$-accessible functor $G:\lD\to\lX$ and functor $H:\lC\to\lX$, an equivalence $\eta:H\to G\circ F$ is said to \emph{exhibit $G$ as the $\Ind_\kappa$-extension of $H$}.  
\end{definition} 
\begin{theorem}\label{thm:condind}
	A functor of $\infty$-categories $F:\lC\to\lD$ exhibits $\lD$ as the $\Ind_\kappa$-completion of $\lC$ if and only if the following hold: 
	\begin{enumerate}
		\item The $\infty$-category $\lD$ admits all small $\kappa$-filtered colimits. 
		\item The functor $F$ is fully faithful. 
		\item For every object $X\in\lC$, its image is $\kappa$-compact. 
		\item The $\infty$-category $\lD$ is generated under small $\kappa$-filtered colimits by the image of $F$. 
	\end{enumerate}
\end{theorem}For the proof, see \cite[Tag 0653]{kerodon2026}. 
\begin{proposition}
	Given a $\infty$-category $\lC$, the following are equivalent:
	\begin{enumerate}
		\item $\lC$ is $\kappa$-compactly generated. 
		\item The inclusion $\lC^\kappa\hookrightarrow\lC$ exhibits $\lC$ as the $\Ind_\kappa$-completion of $\lC^\kappa$. 
	\end{enumerate}
\end{proposition}
For the proof, see \cite[Tag 067B]{kerodon2026}. 
\begin{proposition}
	Given a $\kappa$-compactly generated $\infty$-category $\lC$, the restricted Yoneda embedding $h:\lC\to\Psh(\lC^\kappa)$ is $\kappa$-accessible and fully faithful. 
\end{proposition}
\begin{proof}
	$\kappa$-accessibility of the restricted Yoneda embedding comes from $\lC^\kappa$ being a full sub-category of $\kappa$-compact objects. By definition, there is an equivalence $\eta:h^\kappa\to h|_{\lC^\kappa}$, where $h^\kappa$ denotes the Yoneda embedding of $\lC^\kappa$, then $\eta$ exhibits $h$ as the $\Ind_\kappa$-extension of $h^\kappa$. By \cite[Tag 0651]{kerodon2026}, we conclude that $h$ is fully faithful, since $h^\kappa$ is fully faithful. 
\end{proof}
\begin{definition}
	A $\infty$-category is called \emph{$\kappa$-accessible} if there is a small category $\lD$ together with a functor $F:\lD\to\lC$ exhibiting $\lC$ as the $\Ind_\kappa$-completion of $\lD$.
\end{definition}
\begin{proposition}
	Given a $\kappa$-accessible $\infty$-category $\lC$, the full sub-category $\lC^\kappa$ is essentially small. 
\end{proposition}
For the proof, see \cite[Tag 06K7]{kerodon2026}.
\begin{definition}
	An $\infty$-category $\lC$ is ($\kappa$-)\emph{presentable} if ($\kappa$-)accessible and cocomplete. 
\end{definition}
\begin{definition}
	A functor of $\infty$-categories $F:\lC\to\lD$ is called a \emph{localization} if it admits a fully faithful right adjoint. If the right adjoint to $F$ is ($\kappa$-)accessible, we call $F$ a ($\kappa$-)\emph{accessible localization}.
\end{definition}
\begin{theorem}\label{thm:loccomplete}
	The localization of a cocomplete category is cocomplete. 
\end{theorem}
For the proof, see \cite[Tag 04JW]{kerodon2026}.
\begin{theorem}\label{lem:imofcompact}
	Given a $\kappa$-accessible localization $F:\lC\to\lD$, the image of a $\kappa$-compact object is $\kappa$-compact. 
\end{theorem}
\begin{proof}
	Let $G$ be the $\kappa$-accessible right adjoint to $F$. For every object $X\in\lC$, the functor $\hom_\lD(F(X),-)$ is equivalent to $\hom_\lC(X,G(-))$, which is $\kappa$-accessible if $X$ is $\kappa$-compact. 
\end{proof}
\begin{theorem}
	Given a $\kappa$-presentable $\infty$-category, there is a $\kappa$-accessible localization $F:\Psh(\lC^\kappa)\to\lC$. 
\end{theorem}
\begin{proof}
	We already know that the restricted Yoneda emebdding $h:\lC\to\Psh(\lC^\kappa)$ is $\kappa$-accessible and fully faithful. By completeness and $\lC^\kappa$ being essentially small, $h$ admits a left adjoint $F$, given by left Kan extension of the identity of $\lC$ along $h$. We conclude that $F$ is the desired $\kappa$-accessible localization. 
\end{proof}
\begin{theorem}\label{lem:indcompletion}
	The $\kappa$-accessible localization of a $\kappa$-accessible $\infty$-category is $\kappa$-accessible.  
\end{theorem}
\begin{proof}
	Let $F:\lC\to\lD$ be a $\kappa$-accessible localization, with fully faithful left adjoint $G$. Let $\lC$ be the $\Ind_\kappa$-completion of a small sub-category $\lX$. Denote by $\lY$ the full sub-category of $\lD$ consisting of the images under $F$ of objects in $\lX$. By Theorem \ref{lem:imofcompact}, $\lY$ is a full sub-category of $\kappa$-compact objects. Since $G$ is fully faithful, $F$ is essentially surjective (see \cite[Tag 02FF]{kerodon2026}). Since $F$ preserves colimits, $\lD$ is generated under small $\kappa$-filtered colimits by $\lY$. The conclusion follows from Theorem \ref{thm:condind}.
\end{proof}
\begin{theorem}\label{thm:locpres}
	The $\kappa$-accessible localization of a $\kappa$-presentable $\infty$-category is $\kappa$-presentable. 
\end{theorem}
\begin{definition}
	Given an $\infty$-category $\lC$, a class of morphisms $S$, and a functor $F:\lC\to\lD$, we say \emph{$F$ exhibits $\lD$ as the localization of $\lC$ at $S$} if the following hold: 
	\begin{enumerate}
		\item $F$ sends every morphism in $S$ to an equivalence in $\lD$.
		\item For every $\infty$-category $\lX$, pre-composition with $F$ induces a fully faithful functor \begin{equation}
			\begin{tikzcd}
				\Fun(\lD,\lX) \arrow[r,"F^*"] & \Fun(\lC,\lX)
			\end{tikzcd}
		\end{equation}with essential image consisting of functors sending every morphism in $S$ to an equivalence in $\lX$. 
	\end{enumerate}
\end{definition}
\begin{theorem}
	Given a functor of presentable $\infty$-categories $F:\lC\to\lD$, the following are equivalent: 
	\begin{enumerate}
		\item $F$ is a localization. 
		\item $F$ is cocontinuous and there is a class of morphisms $S$ of $\lC$ such that $F$ exhibits $\lD$ as the localization of $\lC$ at $S$. 
	\end{enumerate}
\end{theorem}
For the proof, see \cite[Tag 06UW]{kerodon2026}. Finally, we state the central property of presentable $\infty$-categories. 
\begin{theorem}[Adjoint Functor Theorem]
	Given a functor of presentable $\infty$-categories $F:\lC\to\lD$, the following hold: 
	\begin{enumerate}
		\item $F$ admits a right adjoint if and only if it preserves small colimits. 
		\item $F$ admits a left adjoint if and only if it is accessible and preserves small limits. 
	\end{enumerate}
\end{theorem}
For the proof, see \cite[Corollary 5.5.2.9]{lurie2009htt}, \cite[Tag 06Q4]{kerodon2026joyal}, or \cite[Theorem 5.2.14]{land2021introduction} (see \cite[Remark 5.2.16]{land2021introduction}).


\section{Stable \texorpdfstring{$\infty$}{oo}-Categories and Spectra}

% Stable homotopy theory is the study of homotopy types once the (suspension,loop) adjunction is made into an equivalence.

\begin{definition}
  Let $\lC$ be an $\infty$-category, an object is called \emph{zero} if both initial and terminal. An $\infty$-category $\lC$ with a zero object is called \emph{pointed}. 
\end{definition}

\newpage

\section{Uhmmmm}
We now use the $\infty$-categorical framework to study spectra. Here are some result that justify the study of spectra. 
\nrnote{Examples of stable phenomena? Freudenthal's suspension theorem, Alexander duality of manifolds, etc.}

\begin{proposition}
  Let $\lD$ be a pointed $\infty$-category. Evaluation at the 0-sphere $S^0$ induces an equivalence \begin{equation}
    \begin{tikzcd}
      \Fun^\mL(\lS_*,\lC) \arrow[r] & \lC
    \end{tikzcd}
  \end{equation}
\end{proposition}
We now introduce the notion of a {triangle}.
\begin{definition}
  Let $\lC$ be a pointed $\infty$-category. A \emph{triangle} in $\lC$ is a commutative diagram of the form \begin{equation}\label{eq:triangle}
    \begin{tikzcd}
      X \arrow[r] \arrow[d] & Y \arrow[d] \\
      0 \arrow[r] & Z
    \end{tikzcd}
  \end{equation}A triangle is \emph{exact}, resp. \emph{coexact}, if it is a pullback, resp. pushout, square. 
\end{definition}
\begin{definition}
  Let $\lC$ be a pointed $\infty$-category. Denote by $\lC^{\Sigma}$, resp. $\lC^\Omega$, the full sub-category of $\Fun(\Delta^1\times\Delta^1,\lC)$ of coexact, resp. exact, triangles of the form 
  \begin{equation}\label{eq:exact}
  \begin{tikzcd}
   X \arrow[r]  \arrow[d] & 0 \arrow[d] \\ 
   0 \arrow[r] & Y
  \end{tikzcd},
\end{equation}
\end{definition}
If $\lC$ is finitely cocomplete, resp. complete, for every object $X$, resp. $Y$, there is a contractible space of coexact, resp. exact, triangles as \cref{eq:exact}. In particular, $\lC\simeq\lC^\Sigma$ and $\lC\simeq\lC^\Omega$.  
\begin{proposition}
  If $\lC$ is a finitely complete and cocomplete pointed $\infty$-category, define the following functors Then the functors
  \begin{equation}
    \begin{tikzcd}
      \Sigma:\lC \arrow[r] & \lC^\Sigma \arrow[r,"\mathrm{ev}_{(1,1)}"] & \lC  & & \Omega:\lC \arrow[r] & \lC^\Omega \arrow[r,"\mathrm{ev}_{(0,0)}"] & \lC
    \end{tikzcd}
  \end{equation}are adjoint ($\Sigma$ is left adjoint to $\Omega$). 
\end{proposition}

\begin{theorem}\label{thm:stable}
  Let $\lC$ be a finitely bicomplete pointed $\infty$-category. The following are equivalent:
  \begin{enumerate}
    \item A triangle is exact if and only if it is coexact. 
    \item $(\Sigma,\Omega)$ is an adjoint equivalence.
    \item A commutative square is a pullback if and only if it is a pushout. 
  \end{enumerate}
\end{theorem}
\begin{definition}
  A finite bicomplete pointed $\infty$-category $\lC$ satisfying any of the equivalent conditions in \cref{thm:stable} is called \emph{stable}.
\end{definition}
If $\bA$ denotes a nice abelian category, there is a \emph{derived $\infty$-category} of $\bA$, denoted $\lD(\bA)$ such that $\Ho\lD(\bA)=D(\bA)$ is the ordinary derived category of $\bA$, i.e. the localization of chain complex at quasi-isomorphisms. From homological algebra, it is known that $D(\bA)$ underlies the structure of a triangulated category, which turns out to be the 1-categorical reflection of $\lD(\bA)$ being stable. 
\begin{proposition}[{\cite[3.11]{lurie2017ha}}]
  If $\lC$ is a stable $\infty$-category, then $\Ho\lC$ has a canonical structure of a triangulated category. 
\end{proposition}To construct the stabilization of a pointed $\infty$-category, there are several approaches, such as reduced excisive functors on $\lS^\mathrm{fin}_*$, the category of pointed, finite spaces, see \cite[1.4.2.8]{lurie2017ha}. Here we consider the more explicit approach using spectrum objects.
\begin{definition}
  Let $\lC$ be a pointed $\infty$-category. A \emph{pre-spectrum object in $C$} consists of a functor $E:\bZ\times\bZ\to\lC$ such that $E(n,m)\simeq0$, for all $n\neq m$. Denote by $\PSp(\lC)$ the category of pre-spectrum objects. The functor $\Omega^{\infty-n}:\PSp(\lC)\to\lC$ is defined as evaluation at $(n,n)$. 
\end{definition}
For every $n$, the diagram \begin{equation}
  \begin{tikzcd}
    E(n,n) \arrow[r]\arrow[d] & 0 \arrow[d] \\
    0 \arrow[r] & E(n+1,n+1) 
  \end{tikzcd}
\end{equation}determines a pair of adjoint morphisms \[\alpha_n:\Sigma E(n,n)\to E(n+1,n+1),\qquad\beta_n:E(n,n)\to\Omega E(n+1,n+1)\]
\begin{definition}
  Let $\lC$ be a pointed $\infty$-category. A \emph{spectrum object in $\lC$} consists of a pre-spectrum object $E$ such that $\beta_n$ is an equivalence, for all $n$. Denote by $\Sp(\lC)\subseteq\PSp(\lC)$ the full sub-category of spectrum objects. 
\end{definition}
\begin{theorem}
  Let $\lC$ be a presentable pointed $\infty$-category, then $\Omega^{\infty-n}:\Sp(\lC)\to\lC$ admits a left adjoint $\Sigma^{\infty-n}:\lC\to\Sp(\lC)$, for every $n$. 
\end{theorem}
In particular, $\Sigma^\infty:\lC\to\Sp(\lC)$ has the following universal property. 
\begin{theorem}
 Let $\lC$ be a presentable pointed $\infty$-category. Given a stable $\infty$-category $\lD$, pre-composition by $\Sigma^\infty$ induces an equivalence 
 \begin{equation}
  \begin{tikzcd}
    \Fun^\mL(\Sp(\lC),\lD) \arrow[r] & \Fun^\mL(\lC,\lD)
  \end{tikzcd}
 \end{equation}In particular, for $\lC=\lS_*$, evaluation at the \emph{sphere spectrum} $\bS=\Sigma^\infty S^0$ induces an equivalence
 \begin{equation}
  \begin{tikzcd}
    \Fun^\mL(\Sp(\lS_*),\lD) \arrow[r] & \lD
  \end{tikzcd}
 \end{equation}
\end{theorem}
\begin{definition}
  The \emph{$\infty$-category of spectra} is the category of spectrum objects in pointed spaces. 
\end{definition}

\section{Generalized Cohomology Theories}

We shall now use the language of $\infty$-categories to reformulate the concept of generalized cohomology theory \` a-l\`a Eilenberg-Steenrod. In this new context, we recall a representability theorem for cohomology theories by spectrum object. 
\begin{remark}
  Denote by $\set^\bZ$ the category of $\bZ$-indexes families of sets. Given an object $S$ and $n\in\bZ$, denote by $\Sigma^n S$ the shifted family $(\Sigma^nS)_i=S_{i-n}$.  
\end{remark}
\begin{definition}[{\cite[1.4.1.6]{lurie2017ha}}]\label{def:cohthe}
  Let $\lC$ be a finitely cocomplete pointed $\infty$-category, $\Sigma_\lC:\lC \to \lC$ the induced suspension functor. A \emph{generalized cohomology theory} is a functor $H:\Ho\lC^{\op}\to\set^\bZ$ together with a natural isomorphism $\partial:\Sigma H\to H\Sigma_\lC$ such that:
  \begin{enumerate}
    \item $H$ preserves arbitrary products. In particular, $H^n(0)$ is the one-point set. Given an object $X$, the unique morphism $X\to0$ induces an element $*\simeq H^n(0)\to H^n(X)$, which we denote by $0$. 
    \item Given a coexact triangle $X'\to X\to X''$, if $\eta\in H^n(X)$ has image $0\in H^n(X'')$, then it lies in the image of $H^n(X')\to H^n(X)$. 
  \end{enumerate}
\end{definition}
\begin{theorem}[{\cite[1.4.1.10]{lurie2017ha}}]
 Let $\lC$ be a nice $\infty$-category and $(H,\partial)$ a generalized cohomology theory, then, for every $n$, the functor $H^n$ is representable by an object $E(n)$. 
\end{theorem}
The natural isomorphism $\partial$ translates into an equivalence $E(n)\simeq\Omega E(n+1)$, which is then used to construct a spectrum object representing the cohomology theory $H^n$, see \cite[1.4.1.11]{lurie2017ha}.
\begin{remark}
  For $\lC=\lS_*$, the above definition of cohomology theory reduces to the classical Eilenberg-Steenrod definition. Since $\lS_*$ is nice, we thus recover the classical \emph{Brown representability theorem}. 
\end{remark}

\chapter{Descent in(to) Higher Category Theory \texorpdfstring{\\}{ } \normalsize \emph{Talk by Hannes Berkenhagen}}

In this lecture we cover the basics of sheaf theory and learn about differential cohomology theories as sheaves of spectra. We begin by recalling the basics of $\set$-valued sheaves over an ordinary category or topological space, progressing then to sheaves of spaces over an $\infty$-category, and finally to sheaves of spectra. At the end, we'll take a closer look to sheaves over the site of smooth manifolds.  
	
\section{Sheaves in Category Theory}

	Let us review sheaves in classical category theory. A good reference remains \cite{maclanemoerdijk1994topos}.
	\begin{definition}
		Let $\lC$ be a category. A \emph{pre-sheaf} on $\lC$ is a functor $F:\lC^{\op}\to\set$. 
	\end{definition}
	\begin{definition}
		Let $X$ be a topological space, we denote by $\Open_X$ the partial order of open subsets of $X$, viewed as a category. We call a pre-sheaf on $\Open_X$ a \emph{pre-sheaf on $X$}. Given $F:\Open_X^{\op}\to\set$ and $u\in F(U)$, for every $V\subseteq U$, we write $u|_V$ for the image of $u$ under the map $F(U)\to F(V)$. 
	\end{definition}	
	\begin{definition}
		Let $X$ be a topological space. A pre-sheaf on $X$ is a \emph{sheaf} if, for every open $U\subseteq X$ and open cover $\{U_i\subseteq U\}_i$, the following is an equalizer diagram:
		\[
		\begin{tikzcd}
			F(U) \arrow[r] & \prod_iF(U_i) \arrow[r, shift left=1] \arrow[r, shift right = 1] &   \prod_{i,j} F(U_i \cap U_j) 
		\end{tikzcd}
		\]
	\end{definition}
	Now, we wish to extend the notion of sheaf to pre-sheaves on a generic category $\lC$.
	To do so, we first generalize the concept of cover, for examples by specifying for each object $C$ a class of families of maps $\{f_i:C_i\to C\}$ that {cover} $C$. Finally, in the sheaf condition, we replace intersections $U_i\cap U_j$ with fiber products (assuming $\lC$ has the necessary limits). Then a pre-sheaf $F:\lC^{\op}\to\set$ is a sheaf if, for every object $C$ and any covering family $\{f_i:C_i\to C\}_i$, the following is an equalizer diagram:\[
		\begin{tikzcd}
			F(C) \arrow[r] & \prod_i F(C_i) \arrow[r, shift left=1] \arrow[r, shift right = 1] &   \prod_{i,j} F(C_i \times_C C_j) 
		\end{tikzcd}
		\]
	The above idea leads to the notion of \emph{coverage}, here we will use another language, namely that of \emph{sieves}. 
	
	\section{Sheaves via Grothendieck Topologies}
	
	We now generalize from sheaves on a topological space to sheaves on a category via the additional data of a Grothendieck topology. We begin by recalling a few definitions. Denote by $y$ the \textit{Yoneda embedding} $\lC\to\lP(\lC)$, mapping $C$ to the pre-sheaf $y(C):=\hom_\lC(-,C)\colon\lC^{\op}\to\set$. 
	\begin{definition}
		A \emph{sub-functor} to a pre-sheaf $F\colon\lC^{\op}\to\set$ consists of a pre-sheaf $S:\lC^{\op}\to\set$ such that $S(C)\subseteq F(C)$, for all objects $C$, and the sub-set inclusions define a natural transformation $S\to F$. The last condition is equivalent to the following: Given $f\colon D\to C$ and $u\in S(C)$, then $F(f)(u)\in S(D)$. 
	\end{definition}
	\begin{definition}
		Given a pre-sheaf $F\colon\lC^{\op}\to\set$, denote by $\lC_{/F}$ the category of pairs $(C,c)$, consisting of an object $C$ and a morphism of pre-sheaves $y(C)\xrightarrow{c} F$.
	\end{definition}
	
	\begin{theorem}[Density Theorem]\label{thm:density}
		$F$ is canonically isomorphic to the colimit of the diagram \begin{equation}
			\begin{tikzcd}
				\lC_{/F} \arrow[r] & \lC \arrow[r,"y"] & \lP(\lC)
			\end{tikzcd}
		\end{equation}where the first functor is the canonical projection.
	\end{theorem}
	\begin{definition}
		Let $\lC$ be a category. A \emph{sieve} on an object $C$ is a sub-functor of $y(C)$, with $y(C)$ being the \textit{maximal sieve}. Given a morphism $f\colon D\to C$ and a sieve $S$ on $C$, denote by $f^*(S)$ the \textit{pullback sieve}, mapping $X$ to the set of arrows $g\colon X\to D$ such that $fg\in S(X)$. 
	\end{definition}
	\begin{definition}\label{def:generatedsieve}
		Let ${\lF}$ be a set of morphisms with codomain $C\in\lC$. Define the \emph{sieve generated by $\lF$} as the sieve $S_\lF\subseteq\hom_\lC(-,C)$ mapping $D$ to the set of morphisms $D\to C$ that factor through some element of $\lF$. 
	\end{definition}
	\begin{definition}
		Let $\lC$ be a category. A \emph{Grothendieck topology} on $\lC$ is a collection $J(C)$ of sieves, called \textit{covering sieves}, on every object $C$ in $\lC$, satisfying the following axioms:
		\begin{enumerate}
			\item The maximal sieve is a covering sieve.
			\item Given a covering sieve $S$ and a morphism $f\colon D \to C$, the pullback sieve $f^*(S)$ is a covering sieve on $D$.
			\item A sieve $S$ on $C$ is a covering sieve if, for some covering sieve $S'$ on $C$ and every map $f\in S'$, the pullback sieve $f^*(S)$ is also a covering sieve. 
		\end{enumerate}
		A \emph{Grothendieck site} is a pair $(\lC, J)$ where $\lC$ is a category and $J$ is a Grothendieck topology on $\lC$.
	\end{definition}
	\begin{example}If every sieve is a covering sieve, we obtain the \emph{discrete Grothendieck topology}. If the only covering sieves are the maximal ones, we obtain the \textit{indiscrete Grothendieck topology}.
	\end{example}
	\begin{example}
		Let $X$ be a topological space. A sieve on an object $U$ of $\Open_X$ corresponds to a collection of open subsets of $U$. Define a sieve to be a covering sieve on $U$ if the corresponding family of subsets covers $U$ in the classical sense.
	\end{example}
	\begin{example}\label{ex:top}
		Let $\Top$ be the category of topological spaces and continuous maps. Define a sieve $S$ to be a covering sieve if generated, in the sense of \cref{def:generatedsieve}, by a family $\lO$ of jointly surjective, open embeddings. 
	\end{example}
	\begin{example}Let $\Mfd$ be the category of smooth manifolds and smooth maps. We can define a topology in the same way as \cref{ex:top}, using families of jointly surjective, smooth open embeddings. 
	\end{example}
	\begin{definition}\label{def:indtop}
		Let $(\lC,J)$ be a site and $\lD\subseteq\lC$ a full sub-category. The \emph{sub-category Grothendieck topology} on $\lD$ assigns to each object $C\in\lD$ the collection of sieves $S\subseteq\hom_{\lD}(-,C)$ that are the restriction to $\lD$ of covering sieves of $C$ in $\lC$.
	\end{definition}
	\begin{definition}\label{def:slicetop}
		Let $(\lC,J)$ be a site and $C\in\lC$. Given $f:D\to C$, there is a canonical identification $\alpha_f:(\lC_{/C})_{/f}\simeq\lC_{/D}$. The \emph{slice Grothendieck topology} on $\lC_{/C}$ assigns to each object $f:D\to C$ the collection of sieves $\lD\subseteq(\lC_{/C})_{/f}$ that correspond to covering sieves of $D$ under $\alpha_f$. 
	\end{definition}
	We are now ready to define sheaves on a Grothendieck site. Given a sieve $S$ on $C$, notice that the category $\lC_{/S}$ can be identified with the full sub-category of $\lC_{/C}$ spanned by morphisms $f\in S$.
	\begin{remark}
		A sieve on $C$ is equivalent to the datum of a sub-category $\lD\subseteq\lC_{/C}$ closed un pre-composition, i.e. if $f\in\lD$, then $fg\in\lD$, for every $g$ composable with $f$. 
	\end{remark} 
	\begin{definition}\label{map}
		Given a pre-sheaf $F\colon \lC^{\op}\to\set$, denote by $\hom(\lC_{/S},F)$ the limit of the following diagram:
	\begin{equation}\label{eq:comparison}
		\begin{tikzcd}
			(\lC_{/S})^{\op}\arrow[r] & \lC^{\op}\arrow[r,"F"] & \set
		\end{tikzcd}
	\end{equation}The inclusion $\lC_{/S}\subseteq\lC_{/C}$ induces a natural \textnormal{comparison map} $\hom(\lC_{/C},F)\to\hom(\lC_{/S},F)$.
	\end{definition}
	 Notice that $\hom(\lC_{/C},F)$ is the limit of a diagram indexed by the category $(\lC_{/C})^{\op}$, which has an initial object, the identity of $C$, therefore the limit is canonically isomorphic to $F(C)$, the value of the diagram at the initial object.
	\begin{definition}
		Let $(\lC, J)$ be a Grothendieck site. A {pre-sheaf} $F : \lC^{\op} \to \set$ is a \emph{sheaf} if, for every object $C$ and covering sieve $S$ of $C$, the canonical morphism
		\[ F(C)\cong\hom(\lC_{/C},F) \to \hom(\lC_{/S},F)\]
		is an isomorphism. Explicitly, an object in $\hom(\lC_{/S},F)$ consists of an element $x_f\in F(D)$, for every map $f\in S$ (where $D$ is the domain of $f$), such that $x_{fg}=F(g)(x_f)$, for every $f$ and $g$ composable. The canonical morphism takes $x\in F(C)$ to the tuple with $x_f:=F(f)(x)$. 
	\end{definition}Using \cref{thm:density} and Yoneda's lemma, one can construct a canonical isomorphism $\hom(\lC_{/S},F)\simeq\hom_{\lP(\lC)}(S,F)$, such that the following diagram commutes \begin{equation}
		\begin{tikzcd}
			F(C) \arrow[d,"\cong"] \arrow[r] & \Hom(\lC_{/S},F) \arrow[d,"\cong"] \\
			\Hom_{\lP(\lC)}(y(C),F) \arrow[r] & \Hom_{\lP(\lC)}(S,F)
		\end{tikzcd}
	\end{equation}where the bottom map is the restriction map induced by the inclusion $S\subseteq y(C)$. In particular, we can rephrase the sheaf condition as follows:
	\begin{corollary}
		A pre-sheaf $F$ on a Grothendieck site $(\lC, J)$ is a sheaf if and only if for every object $C$ and covering sieve $S$ of $C$, the canonical morphism
		\[ F(C)  \cong \hom_{\lP(\lC)}(y(C),F)\to \hom_{\lP(\lC)}(S,F) \]
		is an isomorphism. 
	\end{corollary}
	
	\section{Sheaves on \texorpdfstring{$\infty$}{oo}-Categories}
	Now that we know how  to define sheaves on 1-categories, we generalize the definition of sheaves to the case of $\infty$-categories. Here we largely follow \cite{lurie2009htt}. Recall that a sieve on $C$ determines a full sub-category of $\lC_{/C}$. Conversely, a full sub-category $\lD$ of $\lC_{/C}$ determines a sieve if and only if $f\in\lD$ implies $fg\in\lD$, for every $g$ composable with $f$. 
	\begin{definition}
		Let $\lC$ be an $\infty$-category and $C$ an object of $\lC$. A \textit{sieve} on $C$ is a full sub-$\infty$-category $\lD$ of $\lC_{/C}$ such that $f\in\lD$ implies $fg\in\lD$, for every $g$ composable with $f$.  
	\end{definition}
	A sieve $\lD\subseteq\lC_{/C}$ induces a sieve $\lD'\subseteq\Ho(\lC)_{/C}$ by taking homotopy classes of 1-morphisms in $\lD$. 
	\begin{definition}
		Let $\lC$ be an $\infty$-category. A \emph{Grothendieck topology} on $\lC$ is a collection $J(C)$ of sieves for every object $C$, such that it induces a Grothendieck topology on $\Ho(\lC)$. 
	\end{definition}
	
	Notice, we have the following compatibility observation.
	\begin{lemma}
		Let $\lC$ be a $1$-category. A Grothendieck topology on $\lC$ is precisely a Grothendieck topology on $\lC$ seen as an $\infty$-category.
	\end{lemma}
	\begin{definition}
		Let $\lC$ be a $\infty$-category and $C$ an object of $\lC$. Given a sieve $\lD\subseteq\lC_{/C}$ and a pre-sheaf $F:\lC^{\op}\to\lS$, denote by $\hom(\lD,F)$ the limit of the diagram \begin{equation}
		\begin{tikzcd}
			\lD^{\op}\arrow[r] & \lC^{\op}\arrow[r,"F"] & \lS
		\end{tikzcd}
	\end{equation}where the first functor is the opposite of the canonical projection $\lC_{/C}\to\lC$ restricted to $\lD$. Using the same argument as in the 1-categorical case, there is a comparison map $\hom(\lC_{/C},F)\to\hom(\lD,F)$ and a canonical equivalence $\hom(\lC_{/C},F)\simeq F(C)$. 
	\end{definition}
	\begin{definition}
		Let $\lC$ be an $\infty$-category. A pre-sheaf $F\colon\lC^{\op}\to\lS$ is a \emph{sheaf} if for every object $C$ and covering sieve $\lD$ of $C$, the canonical morphism
		\[ F(C)\simeq\hom(\lC_{/C},F) \to \hom(\lD,F) \]
		is an equivalence. Denote by $\Shv(\lC,J)$ the full sub-$\infty$-category of $\Fun(\lC^{\op},\lS)$ spanned by the sheaves on $(\lC,J)$.
	\end{definition}
	
	Note we defined sheaves as local objects in a locally presentable $\infty$-category. Hence, using the formalism of locally presentable $\infty$-categories, we immediately have the following result. 
	
	\begin{proposition}
		The $\infty$-category of sheaves on $\Shv(\lC,J)$ is locally locally presentable.
	\end{proposition}
	
\section{Sheaves with arbitrary Values}
We started with sheaves valued in sets. Then generalized to $\infty$-categorical sheaves valued in spaces. However, we want $\infty$-categorical sheaves valued in spectra. Hence the next step is to generalize the values of our $\infty$-categorical sheaves. Abstractly we obtain such sheaves via the \emph{tensor product} of locally presentable $\infty$-categories, which we review now.

\begin{theorem}[{\cite[Proposition 4.8.1.15]{lurie2017ha}}]
  Let $\lC, \lD$ be locally presentable $\infty$-categories. Then there exists a locally presentable $\infty$-category $\lC \otimes \lD$ together with a functor $F:\lC\times\lD\rightarrow\lC\otimes\lD$ such that 
  \begin{enumerate}
	\item The functor $F$ preserves colimits component-wise.
	\item For any locally presentable $\infty$-category $\lE$, pre-composition with $F$ induces an equivalence between the sub-$\infty$-category $\Fun^\mL(\lC \otimes \lD , \lE)\subseteq\Fun(\lC \otimes \lD , \lE)$ of colimit preserving functors, and the sub-$\infty$-category $\Fun^{\mL,\mL}(\lC \times \lD, \lE)\subseteq\Fun(\lC\times\lD,\lE)$ of functors preserving colimits in each component. 
  \end{enumerate}
\end{theorem}
\begin{theorem}[{\cite[Proposition 4.8.1.17]{lurie2017ha}}]For $\lC$ and $\lD$ locally presentable $\infty$-categories, there is a canonical equivalence between $\lC\otimes\lD$ and $\Fun^R(\lC^{\op},\lD)$, of limit preserving functors $\lC^{\op}\to\lD$. 
\end{theorem}
\begin{definition}
  Let $(C,J)$ be $\lC$ be a Grothendieck site, and $\lD$ be a locally presentable $\infty$-category. A \emph{sheaf on $(\lC,J)$ with values in $\lD$} is a limit preserving functor $F\colon \Shv(\lC,J)^{\op} \to \lD$, which corresponds to an object in $\Shv(\lC,J) \otimes \lD$. 
\end{definition}

Of course, this description is very abstract and ideally we want a more explicit description that we can use in the future. That is the aim of the next sections.

\section{An Explicit Description of Sheaves valued in Spectra}

We have given a formal definition of sheaves valued in spectra, via the tensor product of locally presentable $\infty$-categories. We now want a more explicit description thereof. For this we make use of the following result.

\begin{proposition}[{\cite{lurie2017ha}, see also \cite{gepnergrothnikolaus2015infiniteloopspacemachine}}]
 The inclusion $\Pr^\mL_{\st} \to \Pr^\mL$ of stable locally presentable $\infty$-categories into the category of locally presentable $\infty$-categories admits a left adjoint, which is explicitly given by tensoring with $\Sp$.
\end{proposition}

We can recall from our previous talk that the stabilization of a locally presentable $\infty$-category $\lC$ is given by the $\infty$-category $\Sp(\lC)$ of spectrum objects in $\lC$, therefore $\lC\otimes\Sp  \simeq \Sp(\lC)$, for every locally presentable $\infty$-category $\lC$. In particular, the $\infty$-category $ \Shv(\lC,J)\otimes \Sp$ is equivalent to the category of spectrum objects in $\Shv(\lC,J)$. Finally, since the inclusion $\Shv(\lC,J) \to \Fun(\lC^{\op},\lS)$ preserves and reflects limits, a commutative square in $\Shv(\lC,J)$ is a pullback square if and only if it is a pullback square in $\Fun(\lC^{\op},\lS)$.

\begin{proposition}
  A spectrum object in sheaves on the Grothendieck site $(\lC,J)$ is given by a spectrum object in pre-sheaves on $(\lC,J)$ such that each component is a sheaf.
\end{proposition}

This gives us the following explicit description of sheaves valued in spectra. Recall that $\Omega^{\infty-n}\colon\Sp\to\lS$ is the functor mapping a spectrum object into its $n$-th component.
\begin{theorem}
  There is a canonical equivalence between $\Shv(\lC,J)\otimes\Sp$ and the category of functors $F\colon \lC^{\op} \to \Sp$ such that $\Omega^{\infty - n} F\colon \lC^{\op} \to \lS$ is a sheaf of spaces, for every $n$.
\end{theorem}

\section{An Explicit Description of Sheaves on the Site of Manifolds}
We continue our analysis of sheaves with the aim of providing explicit descriptions. Now we focus on the case of sheaves on the site of manifolds. Recall that a sieve $S$ on a manifold $M$ is a covering sieve if an open cover $\lO$ of $M$ exists, such that every morphism in $S$ factors through the inclusion $U\subseteq M$ of some element $U\in\lO$. From the definition of the topology on $\Mfd$ we have the following: Given a covering sieve $S$ generated by an open cover $\lO$. Denote by $\lI(\lO)$ the poset (viewed as a category) consisting of finite intersections of elements of $\lO$.
\begin{lemma}
	The functor $i:\lI(\lO)\rightarrow\Mfd_{/S}$ is cofinal.
\end{lemma}
In particular, since $\hom(\Mfd_{/S},F)$ is defined as the limit $(\Mfd_{/S})^{\op}\to\Mfd^{\op}\xrightarrow F\lS$ and $i^{\op}$ is \emph{final}, there is a canonical equivalence \[\hom(\Mfd_{/S},F)\simeq\lim(\lI(\lO)^{\op}\to\Mfd^{\op}\xrightarrow F\lS)=\textrm{lim}_{U\in\lI(\lO)}F(U)\]
\begin{proposition}\label{pro:sheaf}
  A pre-sheaf $F\colon \Mfd^{\op} \to \Sp$ is a sheaf if and only if for every manifold $M$ and every open cover $\lO$ of $M$, the canonical morphism $F(M) \to \lim_{U \in \lI(\lO)}F(U)$ is an equivalence.
\end{proposition}

In fact we can check the sheaf condition for sheaves on the site of manifolds with a very specific type of covering families. For a proof of the following result, we refer to \cite[Proposition 3.6.6]{amabeldebrayhaine2021diffcoh}.

\begin{proposition}
  A pre-sheaf $F\colon \Mfd^{\op} \to \Sp$ is a sheaf if and only if
  \begin{enumerate}
    \item $F(\emptyset)$ is terminal.
    \item For every $M = U \cup V$, the canonical square
    \[
      \begin{tikzcd}
        F(M) \ar[r] \ar[d] & F(U) \ar[d] \\
        F(V) \ar[r] & F(U \cap V)
      \end{tikzcd}
    \]is a pullback square. 
    \item For a sequential family of opens $U_1 \subseteq U_2 \subseteq \cdots$ covering $M$, the canonical map $F(M) \to \lim_{n} F(U_n)$ is an equivalence.
  \end{enumerate}
\end{proposition}

\section{Enough Points and Equivalences of Sheaves}
We now have established a very solid understanding of sheaves on the site of manifolds with arbitrary values as pre-sheaves with suitable limit conditions. We now want to use these explicit descriptions to better understand equivalences of sheaves on the site of manifolds.
\begin{definition}[{\cite[Remark 6.5.4.8]{lurie2009htt}}]\label{def:point}
	Given a $\infty$-topos $\lX$, a \emph{point} of $\lX$ is a functor $p:\lS\to\lX$ admitting a left exact left adjoint $p^*$. We say $\lX$ has \emph{enough points} if the following holds: Given a morphism $f$, if $p^*(f)$ is an equivalence for every point $p$, then $f$ is an equivalence.
\end{definition}
Here we are mainly interested in $\infty$-topoi of sheaves, i.e. $\lX=\Shv(\lC)$, for some site $\lC$.
\begin{remark}
	A functor $p:\lS\to\lX$ admitting a left exact left adjoint is called a \emph{geometric morphism}, see \cite[Definition 6.3.1.1]{lurie2009htt}. 
\end{remark}
\begin{definition}\label{def:stalk}
	Let $X$ be a topological space and $\lC$ a locally presentable category. Given a point $x\in X$, denote by $x^*$ the functor $\Shv(X,\lC)\to\lC$ given by $x^*(F)=\colim_{x\in U}F(U)$. The functor $x^*$ is called 
	\emph{stalk at $x$} and is left adjoint to the \emph{skyscraper at $x$} functor, denoted as $x_*$ and given by $x_*(S)(U)=S$, if $x\in U$, and $=1$, otherwise (where $1\in\lC$ is the terminal object).
\end{definition}
\begin{remark}\label{rmk:pointsheaf}
	Using that filtered colimits commute with finite limits (see \cite[Example 7.3.4.7]{lurie2009htt}), one can prove that $x^*$ is left exact and $x_*$ is a point of $\Shv(X)$ according to \cref{def:point}. In particular, every point $x\in X$ induces a point of $\Shv(X)$. 
\end{remark}
\begin{theorem}\label{thm:enougpoints}
	Let $X$ be a smooth manifold, then $\Shv(X)$ has enough points. 
\end{theorem}
\begin{proof}
	\cite[Corollary 7.2.1.17]{lurie2009htt} and \cite[Theorem 7.2.3.6]{lurie2009htt} reduce the claim to $X$ having finite Lebesgue's covering dimension (see \cite[Definition 7.2.3.1]{lurie2009htt}). The Fundamental Urysohn's Identity states that, for separable metric spaces (such as smooth paracompact manifolds), Lebesgue's covering dimension coincides with the \emph{small inductive dimension}, or SID, for short. The SID of a space $X$, denoted by $\mathrm{ind}X$, is defined inductively as follows: 
	\begin{enumerate}
		\item $\mathrm{ind}\emptyset=-1$.
		\item $\mathrm{Ind}X=n$ is the smallest natural number such that every point has a neighborhood basis of open sets $U$ such that $\mathrm{ind}\partial U\leq n-1$. If no such natural number $n$ exists, we set $\mathrm{ind}X=\infty$.
	\end{enumerate}One can check by induction that $\mathrm{ind} X=n$, for every $n$-dimensional topological manifold $X$.
\end{proof}
\begin{remark}\label{rmk:geopoints}
	The proof of \cite[Corollary 7.2.1.17]{lurie2009htt} actually shows that the collection of points $x_*:\lS\to\Shv(X)$, for all $x\in X$, is enough to detect equivalences, namely $f$ is an equivalence if $x^*(f)$ is an equivalence, for every $x\in X$. 
\end{remark}
% Do smooth manifolds have finite Lebesgue dimension or am I tripping?
\begin{theorem}\label{thm:mfdenoughpoints}
	The $\infty$-category $\Shv(\Mfd)$ has enough points. 
\end{theorem}
\begin{proof}
	Let $M$ be a smooth manifold, there is a faithful functor $\Open_M\to\Mfd$ that is a morphism of sites (see \cite[Definition A.10]{pstragowski2022syntheticspectracellularmotivic}) and satisfies the \emph{covering lifting property} (see \cite[Definition A.12]{pstragowski2022syntheticspectracellularmotivic}), therefore the restriction functor $-|_{M}:\Shv(\Mfd)\to\Shv(M)$ is a left adjoint (see \cite[Proposition A.13]{pstragowski2022syntheticspectracellularmotivic}) and a right adjoint (hence left exact). In particular, a point $p:\lS\to\Shv(M)$, for any manifold $M$, induces a point of $\Shv(\Mfd)$ as the composition of $p:\lS\to\Shv(M)$ with the right adjoint $\Shv(M)\to\Shv(\Mfd)$. The left adjoint is $\Shv(\Mfd)\xrightarrow{-|_M}\Shv(M)\xrightarrow{p^*}\lS$. 
	
	Let $f$ be a morphism such that $p^*(f)$ is an equivalence for every point $p$ of $\Shv(\Mfd)$. In particular, $p^*(f|_M)$ is an equivalence for every point of $\Shv(M)$, hence $f|_M$ is an equivalence in $\Shv(M)$, by \cref{thm:enougpoints}. Finally, if $f|_M$ is an equivalence, for every $M$, then $f$ is an equivalence. 
\end{proof}
The proof of \cref{thm:mfdenoughpoints} together with \cref{rmk:geopoints} implies the following: 
\begin{corollary}\label{cor:points}
  A morphism $f$ in $\Shv(\Mfd)$ is an equivalence if and only if for every manifold $M$ and every point $x \in M$, the map on stalks $x^*(f)$ is an equivalence.
\end{corollary}
Note, for every $n$-manifold $M$ and $x\in M$, there is an open embedding $x':\bR^n\to M$ such that $x'(0)=x$. The induced map $\Open_{\bR^n,0}^{\op}\to\Open_{M,x}^{\op}$ is final (since $\{x'(U)|0\in U\subseteq\bR^n\}$ forms a neighborhood basis for $x$), therefore it induces a natural equivalence $0^*\simeq x^*$. Denote by $0_n$ the origin of $\bR^n$, the previous argument together with \cref{cor:points} implies the following:
\begin{corollary}
  A morphism $f$ of sheaves on $\Mfd$ is an equivalence if and only if for all $n \geq 0$, the induced map on stalks $(0_n)^*(f) $ is an equivalence.
\end{corollary}

\section{Sheaves of Manifolds via the Euclidean site}

According to \cref{def:indtop}, every full sub-category $\lB\subseteq\lC$ of a site $\lC$ inherits a Grothendieck topology. If $\lB$ is \emph{dense} (see \cite[Definition 2.2.1]{johnstone2002elephanti}) then the restriction functor $\Shv(\lC,\set)\to\Shv(\lB,\set)$ is an equivalence, this is the \emph{comparison lemma} (\cite[Theorem 2.2.3]{johnstone2002elephanti}). However, this is no longer true with sheaves valued in $\infty$-categories, i.e. a dense sub-site $\lB\subseteq\lC$ does not induce an equivalence between $\infty$-categories of sheaves (see \cite[Counterexample 20.4.0.1]{lurie2018sag}). But, the comparison lemma from $\infty$-sheaves \emph{does} hold for $\lC=\Mfd$ because of \emph{hypercompleteness}, as we'll show. This comparison lemma allows us to relate statements about sheaves on dense sub-categories of $\Mfd$ to statements about sheaves on all manifolds.

Let $\lC$ be a 1-site and $\lP_\Delta(\lC)=\Fun(\lC^{\op},\sset)$ the ordinary category of simplicial pre-sheaves on $\lC$. Given a pre-sheaf $F:\lC^{\op}\to\set$, denote by $F_\delta$ the simplicial pre-sheaf mapping $C$ to the constant simplicial set with value $F(C)$. 
\begin{definition}\label{def:locepi}
	Given a morphism of pre-sheaves of sets $f:F\to\hom_\lC(-,C)$, the \emph{image of $f$} is the sieve $S_f$  mapping $D$ to the image of the map of sets $F(D)\to\hom_C(D,C)$. If $C\in\lC$ and $\lC$ is a site, we call $f$ a \emph{local epimorphism} if $S_f$ is a covering sieve of $C$. A morphism of pre-sheaves $F\to G$ is a local epimorphism if, for every $y(C)\to G$, the projection $y(C)\times_GF\to y(C)$ is a local epimorphism. 
\end{definition}
\begin{definition}[{\cite[Lemma 4.9]{dugger2003hypercoverssimplicialpresheaves}}]
	Given $C\in\lC$, an \emph{hypercover} of $C$ consists of a simplicial pre-sheaf $U:\lC^{\op}\to\sset$ together with a morphism $U\to y(C)_\delta$, called \emph{augmentation}, such that the maps $U_0\to y(C)$ and $U_n\simeq \hom_{\sset}({\Delta^n},U)\to\hom_{\sset}({\partial\Delta^n},U)$, for every $n$, are local epimorphisms, see \cref{def:locepi}. An hypercover has \emph{height $k$} if $U_n\to\hom_{\sset}(\partial\Delta^n,U)$ is an isomorphism, for all $n\geq k$. 
\end{definition}
\begin{remark}\label{rmk:cech}
	An hypercover has height 0 if and only if it is the \v Cech nerve of the morphism $U_0\to y(C)$. 
\end{remark}
\begin{definition}
	Denote by $\widehat{\Shv}(\lC)\subseteq\Shv(\lC)$ the full sub-$\infty$-category of sheaves satisfying \emph{hyperdescent}, meaning that, for every object $C$ and hypercover $U$ of $C$, the augmentation induced an equivalence 
	$$F(C)\simeq\hom_{\lP(\lC)}(y(C),F)\to\lim_n\hom_{\lP(\lC)}(U_n,F)$$ $\Shv(\lC)$ is called \emph{hypercomplete} if every sheaf satisfies hyperdescent. 
\end{definition}
\begin{definition}[{\cite[Definition 3.12.2]{barwick2020exodromy}}]\label{def:dense}
	Let $\lC$ be a site. A full sub-category $\lB\subseteq\lC$ is \emph{dense} if for every object $C\in\lC$ there is a family $\lF$ of morphisms with codomain $C$ and domain in $\lB$, such that $S_\lF$ (see \cref{def:generatedsieve}) is a covering sieve. 
\end{definition}
\begin{theorem}\label{thm:mfdhypercomplete}
	$\Shv(\Mfd)$ is hypercomplete. If $i:\lD\subseteq\Mfd$ is a full, dense sub-site, then $\Shv(\lD)$ is also hypercomplete. Moreover, the adjunction $(i_*,i^*):\Shv(\lD)\to\Shv(\Mfd)$ is an equivalence.
\end{theorem}
\begin{proof}
	Hypercompleteness of $\Shv(\Mfd)$ follows from \cref{thm:mfdenoughpoints}, since having enough points implies hypercompleteness, see \cite[Example 3.11.10]{barwick2020exodromy}. The second and third claims are a direct consequence of \cite[Corollary 3.12.13]{barwick2020exodromy} and the first claim. By \cite[Proposition 3.12.11]{barwick2020exodromy}, pre-sheves restriction $i^*$ preserves sheaves satisfying hyperdescent. 
\end{proof}Recall that, given a locally presentable category $\lE$, there is a canonical equivalence $\Shv(\lC,\lE)\simeq\Shv(\lC)\otimes\lE$. In particular, the following is a simple corollary of \cref{thm:mfdhypercomplete}:
\begin{corollary}\label{cor:equivalence}
For every locally presentable category $\lE$ and full, dense sub-site $\lD\subseteq\Mfd$, the restriction functor $\Shv(\Mfd,\lE)\to\Shv(\lD,\lE)$ is an equivalence. 	
\end{corollary}
We will mainly focus on the following sub-site of $\Mfd$.
\begin{definition}
	Denote by $\Euc$ the full sub-category of $\Mfd$ generated by Euclidean manifolds, i.e. $\bR^n$, for every $n$. $\Euc\subseteq\Mfd$ is dense in the sense of \cref{def:dense}. 
\end{definition}
Related to the notion of hypercompleteness is the notion of \emph{Postnikov completeness}. A space $S$ is \emph{$n$-truncated}, $n\geq0$, if $\pi_i(S)\simeq0$, for all $i>n$. We say $S$ is \emph{$(-1)$-truncated} if $S$ if it's either empty or contractible. Denote by $\tau_{\leq n}\lS\subseteq\lS$ the full sub-category of $n$-truncated spaces. The inclusion functor admits a left adjoint $\tau_{\leq n}$, and so every space has a natural map $S\to\tau_{\leq n}S$. A classical result of algebraic geometry is that $S\simeq\lim_n\tau_{\leq n}S$. The notion of truncated object makes sense in a general $\infty$-category. 
\begin{definition}[{\cite[Definition 5.5.6.1]{lurie2009htt}}]
	Let $n\geq-1$, an object $C\in\lC$ is $n$\emph{-truncated} if $\hom_\lC(D,C)$ is a $n$-truncated space, for every $D\in\lC$. Denote by $\tau_{\leq n}\lC$ the full sub-category of $n$-truncated objects. A morphism $f:C\to D$ \emph{exhibit $D$ as the $n$-truncation of $C$} if $D$ is $n$-truncated and the induced morphism $\tau_{\leq n}C\to D$ is an equivalence. 
\end{definition}
If $\lC$ is a locally presentable $\infty$-category, $\tau_{\leq n}\lC$ can be characterized as a class of local objects (see \cite[Proposition 5.5.6.18]{lurie2009htt}). In particular, the inclusion functor $\tau_{\leq n}\lC\subseteq\lC$ admits a left adjoint $\tau_{\leq n}:\lC\to\tau_{\leq n}\lC$. Given an object $C\in\lC$, its truncations form a sequence 
\begin{equation}
	\cdots\to\tau_{\leq2}C\to\tau_{\leq1}C\to\tau_{\leq0}C
\end{equation}
where $\tau_{\leq n+1}C\to\tau_{\leq n}C$ is the $(n+1)$-truncation of the map $C\to\tau_{\leq n}C$ together with the fact that $\tau_{\leq n+1}E\simeq E$ naturally, for every $n$-truncated $E$. Let $\bZ_{\geq0}$ be the partial order of natural numbers. 
\begin{definition}[{\cite[Definition A.7.2.1]{lurie2018sag} and \cite[Definition 5.5.6.23]{lurie2009htt}}]
	A \emph{Postnikov pre-tower} is a functor $X:\bZ_{\geq0}^{\op}\to\lC$ such that $X_{n+1}\to X_n$ exhibits $X_n$ as the $n$-truncation of $X_{n+1}$, for every $n\geq0$. Denote by $\pos(\lC)$ the category of Postnikov pre-towers, then there is a comparison functor $\tau:\lC\to\pos(\lC)$. We say $\lC$ is \emph{Postnikov complete} if $\tau$ is an equivalence. 
\end{definition}One can show that $\tau$ is left left adjoint to the limit functor $\lim:\Fun(\bZ_{\geq0}^\mathrm{op},\lC)\to\lC$ restricted to Postnikov pre-towers. In particular, if $\lC$ is Postnikov complete, then $C\simeq\lim_n\tau_{\leq n}C$. Let $\lX=\Shv(\lC)$, for some site $\lC$, and $\widehat{\lX}=\widehat{\Shv(\lC)}$, then $\tau_{\leq n}\lX\subseteq\widehat{\lX}$, for every $n$, see \cite[Lemma 6.5.2.9]{lurie2009htt}. In particular, if $\lX$ is Postnikov complete, then every object $C$ is equivalent to the limit of its truncations, which implies $C$ is hypercomplete (Since $\widehat{\lX}$ is a category of local objects in $\lX$, it is closed under limits in $\lX$).
\begin{theorem}
	If $\Shv(\lC)$ is Postnikov complete, then it is hypercomplete. 
\end{theorem} 

\chapter{Sheaves of Spectra and Deligne Cohomology \texorpdfstring{\\}{ } \normalsize \emph{Talk by Alessandro Nanto}}

In the last talk we learned the definition of a differential cohomology theory, as a sheaf valued in spectra on the site of manifolds. This talk continues our journey through differential cohomology theories, and focuses on the following three topics:
\begin{enumerate}
  \item We want to learn how to construct non-trivial examples out of sheaves valued in chain complexes.
  \item We want to understand how we can extend classical cohomology operations to the setting of differential cohomology theories.
  \item We want to introduce suitable analogues of fiber-wise integration.  
\end{enumerate}
%-----TO BE CONTINUED (This is one can of worms I can't deal with right now -_-)----------------

\section{Abelian Groups, Spectra and the Heart}
Let us start by reviewing the relation between abelian groups, rings and spectra. 
\begin{definition}
    Let $n\in\bZ$ and $X$ be a spectrum, define $\pi_n(X):=\pi_0(\Omega^{\infty+n}X)=\pi_0(X_{-n})$. We call $\pi_n$ the \textit{$n$-th homotopy group} of $X$. 
\end{definition}
\begin{remark}
 Note that since $X_n\simeq\Omega^2X_{n+2}$, for any $n$, the set $\pi_0(X_n)$ underlies the structure of an abelian group.
\end{remark}

  The category $\Sp$ underlies the structure of a symmetric monoidal $\infty$-category (\cite[Corollary 4.8.2.19]{lurie2017ha}). Following \cite{lurie2017ha}, we denote by $\otimes$ the tensor product on $\Sp$.
  \begin{definition}
    A (commutative) algebra object in $\Sp$ is called an ($\bE_\infty$-)\emph{ring spectrum}, see \cite[Definition 7.1.0.1]{lurie2017ha}. Given a ring spectrum $R$, denote by $\Mod_R$ the corresponding category of \emph{left $R$-module spectra}, see \cite[Definition 7.1.1.2]{lurie2017ha}. 
  \end{definition}
  \begin{remark}
    The sphere spectrum $\bS$ acts as the monoidal unit of $\Sp$, therefore it is a $\bE_\infty$-ring spectrum. The category $\Mod_{\bS}$ is canonically equivalent to $\Sp$. 
  \end{remark}
    \begin{definition}
    Denote by $\Sp_{\geq0}\subseteq\Sp$ the full sub-category generated by \emph{connective spectra}, i.e. spectra $X$ such that $\pi_n(X)\simeq0$, for all $n<0$. Denote by $\Sp^\heartsuit\subseteq\Sp_{\geq0}$ the \emph{heart of spectra}, i.e. the full sub-category generated by spectra $X$ such that $\pi_n(X)\simeq0$, for all $n>0$.
  \end{definition}
  We have the following result relating connective spectra and the heart, which follow immediately.
  \begin{lemma}
  Let $X$ be a connective spectrum. The following are equivalent:
  \begin{enumerate}
    \item $X$ is in the heart. 
    \item $\pi_n(\Omega^\infty X)=0$, for all $n>0$.
    \item $\Hom_{\lS_*}(S,\Omega^\infty X)\simeq0$, for all connected, pointed spaces $S$.
    \item $X$ is local with respect to the class of maps $\Sigma^\infty S\to0$, for every connected pointed space $S$.
  \end{enumerate}
  \end{lemma}

  The category $\Sp_{\geq0}$ is presentable and $\pi_0$ induces an equivalence between the heart and $\Ab$ (\cite[Proposition 1.4.3.6]{lurie2017ha}). The heart is a sub-category of local objects of connective spectra, therefore the inclusion $\Ab\simeq\Sp^\heartsuit\subseteq\Sp_{\geq0}$ is a right adjoint. The category $\Sp_{\geq0}$ is closed under $\otimes$ and, given $X,Y$ connective spectra, 

\begin{align}\label{eq:stablehurewicz}
  \pi_0(X\otimes Y)\simeq\pi_0(X)\otimes\pi_0(Y)
\end{align}
see \cite[Theorem 2.3.28]{davies2024atii}

\begin{definition}Given an abelian group $A$, denote by $HA$ the (unique up to equivalence) spectrum of the heart such that $\pi_0(HA)\simeq A$. We call $HA$ the \emph{Eilenberg-Mac Lane spectrum} of $A$.
\end{definition}
Using \cref{eq:stablehurewicz} and the adjunction between $H$ and $\pi_0$, one can prove $H$, viewed as a functor $\Ab\to\Sp$, is lax monoidal. In particular, if $R$ is a commutative ring, then $HR$ is a connective $\bE_\infty$-ring spectrum. On the other hand, if $R$ is a connective $\bE_\infty$-ring spectrum and $M$ a connective module, then $\pi_0(M)$ is a $\pi_0(R)$-module. 
\begin{definition}
  Given a commutative ring $R$, denote by $\Ch(R)=\Ch(\Mod_{R})$ the ordinary category of unbounded chain complexes. Let $\lD(R)$ be the $\infty$-localization of $\Ch(R)$ at the class of quasi-isomorphisms, called the \emph{derived category}.  
\end{definition} 
Similar to the heart of spectra, given an $\bE_\infty$-ring spectrum $R$, denote by $\Mod_{R}^\heartsuit\subseteq \Mod_{R}$ the full sub-category generated by $R$-modules such that the underlying spectrum belongs to the heart of spectra. 
\begin{theorem}[Stable Dold-Kan Correspondence]\label{thm:stabledk}
  Let $R$ be a commutative ring. 
  \begin{enumerate}
    \item $\Mod_{R}\simeq \Mod_{HR}^\heartsuit$ via taking Eilenberg-Mac Lane spectra.
    \item The above equivalence extends to an equivalence $H:\lD(R)\simeq \Mod_{HR}$ of symmetric monoidal $\infty$-categories.
  \end{enumerate} 
\end{theorem}
\begin{proof}
  (1) is \cite[Proposition 7.1.1.13]{lurie2017ha}, while (2) is \cite[Theorem 7.1.2.13]{lurie2017ha}.
\end{proof}
An interesting consequence of \cref{thm:stabledk} is the following:
\begin{corollary}\label{cor:homotopyvshomology}
  Given $F\in\lD(R)$, then $\pi_n(HF)\simeq H_n(F)$, for all $n\in\bZ$.
\end{corollary}
\begin{proof}
  \begin{align*}
    \pi_n(HF) & =\pi_0(\Omega^{\infty+n}HF)\\
    & \overset{\Circled{1}}{\simeq} \pi_0(\Hom_{\Sp}(\Sigma^n\bS,HF)) \\
    & \overset{\Circled{2}}{\simeq} \pi_0(\Hom_{\Mod_{HR}}(\Sigma^n HR,HF)) & \\
    & \overset{\Circled{3}}{\simeq} \pi_0(\Hom_{\lD(R)}(R[n],F))\\
    & \overset{\Circled{4}}{\simeq} H_n(F)
  \end{align*} \Circled{1} The functor $\Omega^{\infty+n}$ is corepresented by the shifted sphere spectrum $\Sigma^n\bS$. \Circled{2} The forgetful functor $\Mod_{HR}\to\Mod_\bS\simeq\Sp$ is right adjoint to tensoring with $HR$ and $HR\otimes(\Sigma^n\bS)\simeq \Sigma^nHR$. \Circled{3} \cref{thm:stabledk} \Circled{4} $\pi_0$ of the mapping space $\Hom_{\lD(R)}(R[n],F)$ is equivalent to the mapping space $R[n]\to F$ in the \emph{ordinary} derived category of $R$, i.e. homotopy classes of maps $R[n]\to F$, which correspond exactly to classes in $H_n(F)$.
\end{proof}

\section{Locally constant sheaves on manifolds}\label{sec:constant sheaves}

Let $\lC$ be a presentable $\infty$-category. The $\infty$-categorical background given in previous talks allows to conclude the existence of a number of functors. Here we give a (somewhat) explicit formula for one. 
\begin{remark}
  Let $\lB$ be a full, dense sub-site of $\Mfd$, recall then that, for every presentable category $\lC$, the restriction functor induces an equivalence $\Shv(\Mfd,\lC)\simeq\Shv(\lB,\lC)$. 
\end{remark}
Evaluation at $\{0\}$ induces an adjunction $(\Gamma^*,\Gamma):\lC\to\Shv(\Mfd,\lC)$, where the functor $\Gamma$ is evaluation at $\{0\}$, while the left adjoint $\Gamma^*$ maps $C\in\lC$ to the sheafification of the constant pre-sheaf with value $C$. 
\begin{remark}\label{rmk:cotensor}
Every presentable $\infty$-category $\lC$ is uniquely \emph{cotensored over} $\lS$, see \cite[Remark 5.5.2.6]{lurie2009htt}. More explicitly, for every space $S$ and object $C$, there is an object $C^S$ together with a natural equivalence \[\hom_\lS(S,\hom_\lC(-,C))\simeq\hom_\lC(-,C^S)\]
\end{remark}
\begin{definition}\label{def:flat}
  Denote by $\Sing$ the functor $\Mfd\to\lS$ mapping a manifold to its underlying space of \emph{smooth} simplexes. Given a presentable $\infty$-category $\lC$, denote by $\flat$ the composition $\lC\to\Fun(\lS^{\op},\lC)\to\Fun(\Mfd^{\op},\lC)$, the first functor coming from \cref{rmk:cotensor}, the second being pre-composition with $\Sing^{\op}$.
\end{definition}
Explicitly, given an object $C\in\lC$, the associated pre-sheaf $\flat C$ maps a manifold $M$ to $C^{\Sing(M)}$.
\begin{remark}
  Given a topological space $X$, denote by $\widehat{\Sing}(X)$ the corresponding space, i.e. its singular simplicial set. If $X=M$ is a smooth manifold, then $\Sing(M)\subseteq\widehat{\Sing}(M)$ and the inclusion is a homotopy equivalence by Whitney's Approximation Theorem, see \cite[Theorem 1.6]{zuoqin2021whitney}. 
\end{remark}
% Smooth singular simplicial set is homotopy equivalent to continuous singular simplicial set. Lurie proves continuous sss is a cosheaf, but this is naturally equivalent to $\Sing$. 

\begin{lemma}[{\cite[Corollary 6.46]{bunkegepner2021differential}}]\label{lem:flatsheaf}
  $\Sing$ is a cosheaf. In particular,
  $\flat$ factors through $\Shv(\Mfd,\lC)\subseteq\Fun(\Mfd^{\op},\lC)$.
\end{lemma}
\cref{lem:flatsheaf} is essentially the consequence of a generalized version of Seifert-van Kampen theorem, namely \cite[Proposition A.3.2]{lurie2017ha}, stating that, given a topological space $X$ and a covering sieve $\rO$, the space $\widehat{\Sing}(X)$ is the colimit of $\widehat{\Sing}(U)$, over $U\in\lI(\rO)$.
\begin{theorem}\label{thm:leftadjoint}
  $\flat:\lC\to\Shv(\Mfd,\lC)$ is left adjoint to $\Gamma$.
\end{theorem}
\begin{proof}
  The composition $\lC\xrightarrow{\flat}\Shv(\Mfd,\lC)\xrightarrow{|_{\Euc}}\Shv(\Euc,\lC)$ maps an object $C$ to the sheaf $\flat C$ restricted to Euclidean spaces. Since $\bR^n$ is contractible, $(\flat C)(\bR^n)=C^{\Sing(\bR^n)}\simeq C$ and so $\flat$ restricted to $\Euc$ is equivalent to $L$, the functor taking $C$ to the the pre-sheaf with constant value $C$, which is left adjoint to $\Gamma$ restricted to $\Euc$.
\end{proof}
\begin{remark}\label{rmk:constsheaf}
  The proof of \cref{thm:leftadjoint} shows that, given an object $C\in\lC$ and a dense sub-site $\lB\subseteq\Mfd$ of \emph{contractible} manifolds, the constant pre-sheaf with value $C$ is equivalent to the sheaf $\flat C$ restricted to $\lB$, hence it is already a sheaf. 
\end{remark}

\section{Chain complexes and sheaves}

Denote by $\Ch$ the category of chain complexes of abelian groups and $\lD(\bZ)$ the associated derived category. Moreover, denote by $\Mod$ the category of $H\bZ$-module spectra. The equivalence of \cref{thm:stabledk} applied point-wise induces an equivalence of sheaf categories
\[\Shv(\Mfd,\lD(\bZ))\simeq\Shv(\Mfd,\Mod)\] 
\begin{definition}\label{def:EMsheaf}
  Given a sheaf $V:\Mfd^{\op}\to\lD(\bZ)$, denote by $HV$ the corresponding sheaf of $H\bZ$-module spectra, called \emph{Eilenberg-Mac Lane sheaf}. 
\end{definition}
In this section, we introduce some technical lemmas regarding sheaves on manifolds valued in $\lD(\bZ)$. We implicitly identify (sheaves of) cochain complexes $V^*$ with (sheaves of) chain complexes with reversed indexing, i.e. $V_n:=V^{-n}$.
\begin{definition}[{\cite[Definition 7.14]{bunkenikolausvoelkl2016diffcoh}}]Given $n\in\bZ$ and a (sheaf of) cochain complex $V^*$, denote by $V^{\geq n}$, resp. $V^{\leq n}$, the \emph{naive truncations} of $V^*$ \[\cdots\to0\to V^n\to V^{n+1}\to\cdots, \qquad \mbox{resp. }\cdots\to V^{n-1}\to V^n\to0\to\cdots\]
\end{definition}
Denote by $C^0$ the sheaf of {sets} associated to $\bR$ by the Yoneda embedding, mapping a manifold $M$ to the set of smooth maps $M\to\bR$. The sheaf $C^\infty$ is a ring object in $\Shv(\Mfd,\set)$, a $C^\infty$-module consisting of a sheaf $V$ together with a natural $C^\infty(M)$-module structure on $V(M)$, for every $M$.
\begin{lemma}[{\cite[Lemma 7.12]{bunkenikolausvoelkl2016diffcoh}}]\label{lem:sheafloc}
  Let $V:\Mfd^{\op}\to\Ch$ a sheaf of chain complexes of $C^\infty$-modules, then the composition $\Mfd^{\op}\xrightarrow{V}\Ch\to\lD(\bZ)$ is a sheaf. 
\end{lemma}
\begin{definition}\label{def:forms}
  Denote by $\Omega^*$ the sheaf $\Mfd^{\op}\to\Ch$ mapping $M$ to its de Rham complex $\Omega^*(M)$. 
\end{definition}
\cref{lem:sheafloc} ensures that the sheaves in \cref{def:forms} is a sheaf after post-composition with the localization functor $\Ch\to\lD(\bZ)$. 

\section{Higher de Rham theorem}

\begin{remark}
  Given a smooth manifold $M$, denote by $C_*^\sm(M,\bZ)$ the singular complex generated by $\Sing(M)$. Given an abelian group $A$, denote by $C^*_\sm(M,A)$ the singular cochain complex valued in $A$ associated to $C_*^\sm(M,\bZ)$. 
\end{remark}
The de Rham homomorphism is a natural chain homomorphism $\Omega^*\to C_\sm^*(-,\bR)$. By a theorem of de Rham, the induced morphism on cohomology is an isomorphism of graded commutative algebras. A stronger result, which we mention in this section, proves that the morphism of chain complexes is also compactible with the DG algebra structures on both complexes, but {up to coherent homotopies}. 
\begin{remark}\label{rmk:Dcotensor}
  Since $\lD(\bZ)$ is presentable, we know that it is cotensored over $\lS$. Given a space $S$ and a chain complex $V_*$, the cotensor $V_*{}^S$ is the cochain complex (viewed as a chain complex) of graded linear maps $C_*(S,\bZ)\to V_*$, from the (normalized) singular chain complex of $S$ to $V_*$, see \cite[Definition 1.3.2.1]{lurie2017ha}. In particular, let $V_*=V$ be concentrated in degree 0, then $V_*{}^S$ is the singular cochain complex of $S$ with values in $V$. 
\end{remark}
Recall the functor $\flat$ defined in \cref{def:flat}. By definition of $\Sing:\Mfd\to\lS$ and \cref{rmk:Dcotensor}, given $V_*$ chain complex, $\flat V$ maps a manifold $M$ to $V_*{}^{\Sing(M)}$. 
\begin{definition}\label{def:derhammorf}
Consider the natural morphism $\Omega^*(M)\to (\flat\bR)(M)=C_\sm^*(M,\bR)$ taking a form $\omega\in\Omega^n(M)$ to the linear map $\int\omega:C^\sm_{n}(M,\bZ)\to\bR$. We call $\mathrm{dr}:\Omega^*\to\flat\bR$ the \emph{de Rham transformation}.  
\end{definition}
\begin{remark}\label{rmk:ainfmor}
  For the explicit data and conditions that determine an $A_\infty$-morphism between $A_\infty$-algebras, see \cite[Proposition 10.2.12]{loday2012algebraic}. Every $A_\infty$-homomorphism $X\to Y$, where $X$ and $Y$ are $A_\infty$-algebras in cochain complexes, has an underlying cochain momorphism. If $\phi:X\to Y$ is the underlying cochain morphism of an $A_\infty$-homomorphism, there is a specified morphism $\mu_2:X\otimes X\to Y$ such that $$\phi(vw)-\phi(v)\phi(w)=\mu_2(dv,w)+(-1)^{|v|}\mu_2(v,dw)-d\mu_2(v,w)$$for all $v,w\in X$. 
\end{remark}
\begin{theorem}[{\cite[Theorem 3.25]{abad2010ainftyrhamtheoremintegration}}]\label{thm:derham} 
  $\mathrm{dr}_M:\Omega^*(M)\to C^*_\sm(M,\bR)$ underlies an $A_\infty$-homomorphism of DG algebras. 
\end{theorem}

\section{Deligne Cohomology}

In this section, we give the definition of the $k$-th Deligne sheaf as the Eilenberg-Mac Lane sheaf (see \cref{def:EMsheaf}) associated to a sheaf $V(k):\Mfd^{\op}\to\lD(\bZ)$. Following that, we give explicit cochain complexes that are quasi-isomorphic to the value of $V(k)$ at a manifold $M$. 
\begin{definition}\label{def:deligne}
  Given $k\in\bN$, define ${\bZ}(k):\Mfd^{\op}\to\lD(\bZ)$ as the pullback 
  \[\begin{tikzcd}
    {\bZ}(k)\arrow[r]\arrow[d] & \Omega^{\geq k}\arrow[d,"\mathrm{dr}"] \\
    \flat\bZ \arrow[r] & \flat\bR
  \end{tikzcd}\]
  We call the corresponding Eilenberg-Mac Lane spectrum $H{\bZ}(k)$ the $k$\emph{-th Deligne sheaf}.
\end{definition}
\subsection{Model A}\label{mod:A} See {\cite[\S 3.2]{hopkinssinger2005diffcoh}}. Let $\bA(k)^*(M)$ to the cochain complex where degree $n$ elements are triples $(c,h,\omega)\in C_\sm^n(M,\bZ)\oplus C_\sm^{n-1}(M,\bR)\oplus\Omega^n(M)$ for which $\omega=0$ if $n<k$, with differential $\delta(c,h,\omega)=(\delta c,\mathrm{dr}_M(\omega)-c-\delta h,d\omega)$. The complex $\bA^*(k)(M)$ fits into a diagram
\[\begin{tikzcd}
    \bA^*(k)(M)\arrow[d]\arrow[r] & \Omega^{\geq k}(M)\arrow[d] \\
    C_\sm^*(M,\bZ) \arrow[r] & C_\sm^*(M,\bR)
  \end{tikzcd}\]which commutes up to homotopy given by the projections $\bA^n(k)(M)\to C_\sm^{n-1}(M,\bR)$. Applying the localization functor $\Ch\to\lD(\bZ)$ to the above diagram we obtain a pullback square, hence $\bA^*(k)(M)\simeq\bZ(k)(M)$. 

\subsection{Model B}\label{mod:B} Recall that every cofiber sequence is a fiber sequence in a stable $\infty$-category, and that $\flat$ preserves the zero object and pushout squares. Consider the diagram
\[
  \begin{tikzcd}
    {\bZ}(k)\arrow[r]\arrow[d] & \Omega^{\geq k}\arrow[d]\arrow[dd,bend left=5em,color=red] \\
    \flat\bZ \arrow[r]\arrow[d] & \flat\bR\arrow[d] \\
    0\arrow[r] & \flat(\bR/\bZ)
  \end{tikzcd}
\]
Since both inner squares are pullbacks, ${\bZ}(k)$ is equivalent to the fiber of the \textcolor{red}{red arrow}. Similar to \cref{mod:A}, consider a smooth manifold $M$, let $\bB(k)^*(M)$ be the cochain complex where degree $n$ elements are pairs $(\chi,\omega)\in C_\sm^{n-1}(M,\bR/\bZ)\oplus\Omega^n(M)$ for which $\omega=0$ if $n<k$, with differential $\delta(\chi,\omega)=(e^{2\pi i\mathrm{dr}(\omega)}-\delta\chi,d\omega)$. Similar to \cref{mod:A}, the complex $\bB(k)^*(M)$ fits into the above diagram so that the outer square is an pullback square in $\lD(\bZ)$, hence it is equivalent to $\hat{\bZ}(k)(M)$.
\begin{definition}[{\cite[Definition 3.4]{hopkinssinger2005diffcoh}}, see also {\cite[Chapter 5]{barbecker2014diffchar}}] Consider a manifold $M$, a \emph{differential character} of degree $n$ consists of a character $\chi:Z^\sm_{n}(M,\bZ)\to\bR/\bZ$ on the group of smooth $n$-cycles of $M$ together with a closed $n$-form $\omega\in\Omega^{n+1}(M)$, such that, for every smooth $(n+1)$-chain $c$, \[\chi(\partial c)=e^{2\pi i\int_c\omega}\]
\end{definition}
It is immediate to see that a cocycle in $\bB(k)^n(M)$ (for $n\geq k$) determines a differential character of degree $n$. 
\subsection{Interlude}\label{int:goodcover}
\begin{definition}
  Given a full sub-site $\lB\subseteq\Mfd$ and a manifold $M$, a $\lB$\emph{-good open cover} of $M$ consists of an open cover $\rO$ such that every finite intersection of elements in $\rO$ is either empty or diffeomorphic to an object of $\lB$. If $\lB=\Euc$ the sub-site of Euclidean spaces, we call a $\lB$-good open cover a \emph{differentiably good open cover}, see \cite[Definition 6.3.9]{fiorenza2011cechcocyclesdifferentialcharacteristic}.
\end{definition}If $\lB$ is the sub-site of contractible manifolds, we recover the notion of a \emph{good open cover}. Let $V$ be a sheaf on manifolds and $W$ a sheaf on $\lB$ together with an equivalence $V|_\lB\simeq W$, then $V$ can be evaluated on a manifold $M$ using $W$, provided that $M$ admits a $\lB$-good open cover, since \[V(M)\simeq\lim_{U\in\lI(\rO)}V(U)\simeq\lim_{U\in\lI(\rO)}W(U)\]
If every manifolds admits a $\lB$-good open cover, then $V$ is completely determined by $W$. Since we mainly care for $\lB=\Euc$, we need the following result. 
\begin{theorem}[{\cite[Theorem A.1]{fiorenza2011cechcocyclesdifferentialcharacteristic}}]
  Every paracompact smooth manifold admits a differentiably good open cover. 
\end{theorem}
\subsection{Model C}\label{mod:C} See {\cite[Lemma 7.3.4]{amabeldebrayhaine2021diffcoh}}. Let $\lB\subseteq\Mfd$ be a dense sub-site of contractible manifolds, consider the following diagram in the category of {sheaves} on $\lB$
\[\begin{tikzcd}
     {\bZ}(k)|_\lB\arrow[d]\arrow[r] & \Omega^{\geq k}|_\lB\arrow[d] \arrow[r] & 0\arrow[d] \\
    \bZ \arrow[r]\arrow[rr,bend right=2.5em,color=red] & \Omega^*|_\lB\arrow[r] & \Omega^{\leq k-1}|_\lB
\end{tikzcd}\] The left square is obtained by: $\Circled{1}$ Apply the restriction to $\lB$ functor to the pullback diagram of \cref{def:deligne} $\Circled{2}$ Substitute $\flat\bR|_\lB$ with $\Omega^*|_\lB$ (see \cref{thm:derham}) and $\flat\bZ|_\lB$ with $\bZ$ (see \cref{rmk:constsheaf}). Since the right square is a pullback, ${\bZ}(k)|_\lB$ is equivalent to the fiber of the \textcolor{red}{red arrow}. Let $\bC(k)^*$ be the sheaf of chain complexes $\bZ\to\Omega^0\to\cdots\to\Omega^{k-1}\to0\to\cdots$, where $\bZ$ is in degree 0 and $\bZ\hookrightarrow\Omega^0$ is the inclusion as constant functions. The complex $\bC(k)^*$ fits into the above diagram, so that the outer square is an pullback, therefore ${\bZ}(k)|_\lB\simeq\bC(k)^*$. 

Since $\bZ(k)|_\lB\simeq\bC(k)^*$, we can use the argument from \cref{int:goodcover} to conclude that, assuming $M$ has a $\lB$-good open cover, 
\[
    {\bZ}(k)(M)\simeq\lim_{U\in\lI(\mathscr O)} \bC(k)^*(U)
  \]
  The limit is then equivalent to the limit of the cosimplicial object $\bC(k)^{*,*}:\Delta\to\lD(\bZ)$ defined as follows:
  \begin{enumerate}
    \item $\bC(k)^{*,n}$ is the product of all $\bC(k)^*(U_0\cap\cdots\cap U_n)$, ranging over the $(n+1)$-tuples of $\rO^{n+1}$. 
    \item Given a monotone map $\alpha:\langle n\rangle\to\langle m\rangle$, define $\alpha^*:\bC(k)^{*,n}\to\bC(k)^{*,m}$ as follows: Given $\theta\in\bC(k)^{*,n}$, the $(U_0,\cdots,U_m)$-component of $\alpha^*(\theta)$ is the $(U_{\alpha(0)},\cdots,U_{\alpha(n)})$-component of $\theta$ restricted to $U_0\cap\cdots\cap U_m$.
  \end{enumerate}
  We can evaluate the limit by lifting $\bC(k)^{*,*}$ to a cosimplicial object in $\Ch$ and apply \cite[Lemma 7.10]{bunkenikolausvoelkl2016diffcoh}, the limit is then equivalent to (the image under localization $\Ch\to\lD(\bZ)$ of) the total complex of the cochain bicomplex $\cdots\to 0\to \bC(k)^{*,0}\to\bC(k)^{*,1}\to\cdots$, with differential \[\delta^{*,n}=\sum_{\ell=0}^{n+1}(-1)^\ell(\ell\mbox{-th face map})^*:\bC(k)^{*,n}\to\bC(k)^{*,n+1}\] 

\section{Multiplicative structure on \texorpdfstring{$\bZ(k)$}{Z(k)}}

Consider the following diagram:
\begin{center}
  \begin{tikzcd}[column sep={2.5cm,between origins},row sep={1.5cm,between origins}]
    & \bZ(k)\otimes\bZ(\ell)\arrow[dr]\arrow[dl]\arrow[d,color=red] & \\
    \flat\bZ\otimes\flat\bZ\arrow[d,"\cup"]\arrow[dr] & \bZ(k+\ell)\arrow[dr]\arrow[dl] & \Omega^{\geq k}\otimes\Omega^{\geq\ell}\arrow[d,"\wedge"]\arrow[dl]\\
    \flat\bZ\arrow[dr] & \flat\bR\otimes\flat\bR\arrow[d,"\cup"] & \Omega^{\geq k+\ell}\arrow[dl]\\
    & \flat\bR & 
  \end{tikzcd}
\end{center}The vertical-outward facing-left square is commutes on the nose, while the vertical-outward facing-right square commutes up to homotopy (see \cref{thm:derham}). Since the horizontal-bottom square is a pullback, there is a contractible space of \textcolor{red}{red arrows} making the above cube commute. Using model A (see \cref{mod:A}) for $\bZ(k)$ and $\bZ(\ell)$, consider the following morphism $\mu:\bZ(k)\otimes\bZ(\ell)\to\bZ(k+\ell)$ \[(c,h,\omega)\otimes(b,u,\eta)\mapsto(c\cup b,h\cup \mathrm{dr}(\eta)+(-1)^{|c|}c\cup u+\mu_2(\omega,\eta),\omega\wedge\eta)\]where $\mu_2$ is the homotopy between $\mathrm{dr}(-)\cup\mathrm{dr}(-)$ and $\mathrm{dr}(-\wedge-)$, see \cref{thm:derham} and \cref{rmk:ainfmor}. One can check that $\mu$ fits into the above diagram, making it commute up to homotopy. 

\chapter{The Road So Far \texorpdfstring{\\}{ } \normalsize \emph{Talk by Nima Rasekh}}

In this talk we summarize what we covered until now and discuss possible future directions.

\section {Summary}

We have established that the classical differential cohomology theory called \emph{Deligne cohomology} can be captured in the language of sheaves on the site $\Mfd$ and valued in the stable $\infty$-category of spectra. This naturally leads to several relevant questions, that were the focus of our talks:
\begin{enumerate}
  \item How do sheaves of spectra related to other frameworks for differential cohomology, namely Hopkins and Singer's hexagon diagram\nrnote{citation}?
  \item How, given a ordinary cohomology theory $E$, can we construct an interesting differential cohomology refining $E$ in our framework?
  \item What are concrete benefits of this approach in contrast to others? 
\end{enumerate} We shall postpone discussing the last point to a later talk. 

\section{Theoretical framework}

We saw that the theoretical framework is fundamentally that of presentable $\infty$-categories, stable $\infty$-categories, and $\infty$-categorical sheaves, as developed by Lurie \cite{lurie2009htt}, among others. 

As it primarily a background, for what follows we will take it for granted, and refer to the relevant sources.

\section{Relation to other approaches}
A lot of historical development of differential cohomology theories has focused on constructing them via specific data, and concretely the input often consists of an ordinary cohomology theory and some geometric data. This perspective is completely absent in the definition we just gave, so how can we reconcile them? This is the central theme of the \emph{fracture square}. This already came up so let us quickly summarize. Let $\lC$ be a presentable $\infty$-category. 

\begin{definition}
  A pre-sheaf $F$ is called \emph{$\mathbb{R}$-invariant} if, for any manifold $M$, the projection $M\times\bR\to M$ induces an equivalence $F(M)\to F(M\times\bR)$. Denote by $\lP_{\bR}(\Mfd,\lC)$ the full sub-category of $\bR$-invariant pre-sheaves, and $\Shv_\bR(\Mfd,\lC)$ the full sub-category of $\bR$-invariant sheaves. 
\end{definition}
Recall from \cite{berkenhagen2025diffcoh3} that the restriction $\Shv(\Mfd,\lC)\to\Shv(\Euc,\lC)$ is an equivalence. 
\begin{lemma}\label{lem:inv}
  A sheaf $F$ is $\bR$-invariant if and only if every smooth map $\bR^n\to\bR^m$ induces an equivalence $F(\bR^m)\to F(\bR^n)$. 
\end{lemma}
\begin{proof}
  The terminal map $\bR^n\to\bR^0$ is the composition of projections $M\times\bR\to M$, therefore if $F$ is a $\bR$-invariant sheaf, $F(\bR^0)\to F(\bR^n)$ is an equivalence. For any map $f:\bR^n\to\bR^m$, there is a commutative diagram \begin{equation}
    \begin{tikzcd}
      F(\bR^0) \arrow[rd,"\simeq"'] \arrow[r,"\simeq"] & F(\bR^m) \arrow[d,"f^*"] \\
      & F(\bR^n)
    \end{tikzcd}
  \end{equation}implying $f^*$ is an equivalence. 

  On the other hand, let $F$ be a sheaf such that, for every smooth map $f$ of Euclidean manifolds, $f^*$ is an equivalence. Let $M$ be a smooth manifold and $\lO$ a differentiably good open cover (see \cite[Theorem 5.4]{nanto2025diffcoh4}), then $\{U\times\bR|U\in\lO\}$ is a differentiably good open cover of $M\times\bR$ and there is a commutative square \begin{equation}
    \begin{tikzcd}
      F(M) \arrow[r,"\simeq"] \arrow[d] & \lim_{U\in\lI(\lO)}F(U) \arrow[d,"\simeq"] \\
      F(M\times\bR) \arrow[r,"\simeq"] & \lim_{U\in\lI(\lO)} F(U\times\bR) 
    \end{tikzcd}
  \end{equation}The horizontal maps are equivalences by the sheaf condition, while the right-vertical map is an equivalence because $U\times\bR\to U$ is a smooth map of Euclidean manifolds, we then conclude that $F(M)\to F(M\times\bR)$ is an equivalence. 
\end{proof}
Recall from \cite{nanto2025diffcoh4} that the left adjoint $\flat:\lC\to\Shv(\Mfd,\lC)$ to the global section functor $\Gamma$ has the following description: Given an object $C\in\lC$, $\flat C$ maps $M$ to $C^{\Sing M}$, where $\Sing M$ denotes the (smooth) singular simplicial set of $M$.
\begin{corollary}\label{cor:equiv}
  The functor $\flat:\lC\to\Shv(\Mfd,\lC)$ factors through $\Shv_\bR(\Mfd,\lC)$. The adjunction \begin{equation}
    \begin{tikzcd}
      \lC \arrow[rr,shift right,"\flat"'] & & \Shv_\bR(\Mfd,\lC) \arrow[ll,shift right,"\Gamma"']
    \end{tikzcd}
  \end{equation}is an equivalence. 
\end{corollary}
\begin{proof}
  Since every Euclidean manifold has contractible underlying homotopy type, $\flat C$ maps every map $f$ of Euclidean spaces to an equivalence. By \cref{lem:inv}, we conclude that $\flat C$ is $\bR$-invariant. 

  Let $F$ be a $\bR$-invariant sheaf, set $C:=F(\bR^0)$, then $F|_{\Euc}$ is equivalent to the constant pre-sheaf with value $C$ (see \cite[Remark 2.7]{nanto2025diffcoh4}), which is $\flat C|_{\Euc}$, therefore $F\simeq \flat C$.
\end{proof}
\begin{theorem}\label{thm:adj}
  There is a 3-terms adjunction 
  \begin{equation}\label{eq:adjunction}
    \begin{tikzcd}
      \Shv(\Mfd,\lC) \arrow[rr,shift left=2.5,"\Pi_\infty"] \arrow[rr,shift right=2.5,"\Gamma"'] & & \lC \arrow[ll,hook,"\flat" description]
    \end{tikzcd}
  \end{equation}
  where $\Gamma$ is the global sections function, i.e. evaluation at $\bR^0$, $\flat$ is the locally constant sheaf functor, $\Pi_\infty$ is in the proof below. There is an induced 3-terms adjunction \begin{equation}
    \begin{tikzcd}
      \Shv(\Mfd,\lC) \arrow[rr,shift left=2.5,"L_\hi"] \arrow[rr,shift right=2.5,"R_\hi"'] & & \Shv_\bR(\Mfd,\lC) \arrow[ll,hook]
    \end{tikzcd}
  \end{equation}where $L_\hi(F)$, resp. $R_\hi(F)$, is the locally constant sheaf associated to $\Pi_\infty(F)$, resp. $\Gamma(F)=F(\bR^0)$. 
\end{theorem}
\begin{proof}
  Since the composition $\lC\xrightarrow{\flat}\Shv(\Mfd,\lC)\to\Shv(\Euc,\lC)$ is the constant functor, it has as left adjoint the colimit functor $\Shv(\Euc,\lC)\to\lC$. Define $\Gamma_\sharp$ as the pre-composition of the colimit functor with $-|_{\Euc}$. The induced adjunction comes from composing the adjunction in \cref{eq:adjunction} with the adjoint equivalence \begin{equation}
    \begin{tikzcd}
      \lC \arrow[rr,shift left=2.5,"\flat"] \arrow[rr,shift right=2.5,"\flat"'] & & \Shv_\bR(\Mfd,\lC) \arrow[ll,"\Gamma" description]
    \end{tikzcd}
  \end{equation}
\end{proof}
\begin{remark}
  In particular, $\Shv_\bR(\Mfd,\lC)$ is both a reflective and coreflective sub-category of $\Shv(\Mfd,\lC)$. Reflectivity could also be deduced from the definition of $\bR$-invariant sheaves as local objects with respect to the class of projections $M\times\bR\to M$. 
\end{remark}
\begin{definition}
  Let $\lC$ be a category, $\lD\subseteq\lC$ a full sub-category. An object $C$ is \emph{right orthogonal to $\lD$} if $\hom_\lC(D,C)$ is contractible, for every $D\in\lD$. Denote by $\lD^\perp$ the full sub-category of objects right orthogonal to $\lD$. 
\end{definition}
\begin{definition}
  A right orthogonal sheaf to $\Shv_\bR(\Mfd,\lC)$ is called \emph{pure}.
\end{definition}
\begin{lemma}
  $F$ is pure if and only if $\Gamma(F)$ is terminal. 
\end{lemma}
\begin{proof}
  By \cref{cor:equiv}, every $\bR$-invariant sheaf is of the form $\flat C$, therefore $F$ is pure if and only if the space $\hom_{\Shv(\Mfd,\lC)}(\flat C,F)\simeq\hom_{\lC}(C,\Gamma(F))$ is contractible, for all $C$, which is equivalent to $\Gamma(F)$ being terminal. 
\end{proof}
%---------------TO BE CONTINUED------------------
\begin{definition}
  Given a sheaf $F$, denote by $\mathrm{Def}(F)$ the fiber of the unit $F\to L_\hi(F)$. 
\end{definition}
  
\begin{definition}
  A sheaf is \emph{pure} if the value at the point is the terminal spectrum. We denote the full sub-category of pure sheaves by $\Shv_{\mathrm{pure}}(\Mfd)$.
\end{definition}

Similarly, the category of pure sheaves is closed under limits and colimits, so we have the following result:

\begin{theorem}
 There is a diagram of adjunctions
 \[
 \begin{tikzcd}[column sep=2cm]
 \Shv_{\mathrm{pure}}(\Mfd) \arrow[r, shift left = 2, "\mathrm{Def}", hookrightarrow] \arrow[r, leftarrow, "\mathrm{Cyc}" description] \arrow[r, shift right = 2, hookrightarrow] & \Shv(\Mfd)
 \end{tikzcd}
 \]
\end{theorem}

\begin{remark}
  The notational choice is not coincidental. The left adjoint to $\mathrm{Cyc}$ is precisely the functor $\mathrm{Def}$, where the domain is restricted to pure sheaves.
\end{remark}

Notice we obviously have the following result.

\begin{lemma}
  A sheaf is \emph{trivial} if it is $\mathbb{R}$-invariant and pure.
\end{lemma}

So, pure and $\bR$-invariant sheaves are ``disjoint''. Even better, they cover everything, which is the gist of the fracture square.

\begin{remark}
Notice in both adjunction diagrams there is one fully faithful functor that is not named and we will abuse notation and directly consider objects in the full sub-category as objects in the larger category.
\end{remark}
\begin{theorem}
 Let $E$ be a differential cohomology theory. Then the following is a pullback square
 \[
  \begin{tikzcd}
    E \arrow[r] \arrow[d] & \mathrm{Cyc} E \arrow[d] \\ 
    L_{hi}E \arrow[r] & L_{hi}\mathrm{Cyc} E
  \end{tikzcd}
 \]
\end{theorem}

In fact we can further expand this pullback square to several other pullback squares:
 \[
  \begin{tikzcd}
    \Sigma^{-1}L_{hi}\mathrm{Cyc} E \arrow[d] \arrow[r] & R_{hi}E \arrow[r] \arrow[d] & 0 \arrow[d] \\ 
    \mathrm{Def} E \arrow[r] \arrow[d] & E \arrow[r] \arrow[d] & \mathrm{Cyc} E \arrow[d] \\ 
    0 \arrow[r] & L_{hi}E \arrow[r] & L_{hi}\mathrm{Cyc} E
  \end{tikzcd}
 \]

We can restructure this diagram to more resemble the fracture square, which is the following:
\[
\begin{tikzcd}
 & R_{hi}(E) \arrow[dr] \arrow[rr] & & L_{hi}(E) \arrow[dr] & \\
 \Sigma^{-1}L_{hi}\mathrm{Cyc} E \arrow[ur] \arrow[dr ]& & E \arrow[ur] \arrow[dr] & & L_{hi}\mathrm{Cyc} E \\
  & \mathrm{Def}(E) \arrow[ur] \arrow[rr] & & \mathrm{Cyc} E \arrow[ur] & 
\end{tikzcd}
\]

 The proof of all these results has as of yet been postponed.


\section{Construction of differential cohomology theories}
Having advanced theory is of course interesting, but not enough, we also want explicit examples. Indeed we saw in past weeks several ways to explicitly construct differential cohomology theories.

\begin{example}
  Given a spectrum (cohomology theory) $E$, we can define the constant sheaf, which is by definition a differential cohomology theory. Even better, it is an $\bR$-invariant sheaf, and every $\bR$-invariant sheaf is obtained this way.

  Notice in this case the pure part is trivial, meaning we get the following, somewhat trivial, diagram of pullback squares:
  \[
  \begin{tikzcd}
    0 \arrow[d] \arrow[r] & E \arrow[r] \arrow[d] & 0 \arrow[d] \\ 
     0 \arrow[r] \arrow[d] & E \arrow[r] \arrow[d] & 0 \arrow[d] \\ 
    0 \arrow[r] & E \arrow[r] & 0
  \end{tikzcd}
 \]
\end{example}

This means we have a very good understanding of $\bR$-invariant sheaves, using classical algebraic topology. However, of course these examples have no geometric content, so we want to go beyond them.

\begin{example}
  Given a sheaf of abelian groups, we can construct the differential cohomology theory given by post-composing with the functor that takes an abelian group $A$ to the associated Eilenberg-MacLane spectrum $H(A)$. 
\end{example}

\begin{example}
  There is a non-trivial way to extend a functor from abelian groups to spectra to one from chain complexes to spectra, which is called the \emph{stable Dold-Kan embedding}. Hence, again via post-composition, every sheaf of chain complexes gives rise to a differential cohomology theory.
\end{example}

\begin{example}
 An important example of the previous example are truncated forms. Given $k \geq 0$, $\Omega^{\geq k}$ is a sheaf that associates to a manifold $M$ the chain complex of $k$-truncated forms on $M$. Using the previous example we hence get a differential cohomology theory which we also denote $\Omega^{\geq k }$.

 Notice, if $k > 0$ then $\Omega^{\geq k}$ is in fact pure, and also not trivial, which means it cannot be $\bR$-invariant.
\end{example}

The approach we used until now helps us generate examples that either pure or $\bR$-invariant, but we would also like to mix them. Here we use the fracture square. The fracture square is not only an abstract theoretical tool, it also gives us very explicit ways to construct differential cohomology theories, using methods closer to the historical approaches. Concretely, we now have the following results:

\begin{corollary}
  A differential cohomology theory is uniquely determined by the following three pieces of data:
  \begin{enumerate}
    \item An $\bR$-invariant sheaf $E_{\bR}$, which is equivalently a spectrum, or equivalently a cohomology theory.
    \item A pure sheaf $E_{\mathrm{pure}}$. 
    \item A morphism of spectra $E_{\bR} \to L_{hi}E_{\mathrm{pure}}$.
  \end{enumerate}
\end{corollary}

More specifically, we often face the situation that we have a fixed cohomology theory $E$, and want to find a lift of sorts.

\begin{definition}
  Let $E$ be a cohomology theory. A \emph{differential refinement of $E$} is a differential cohomology $\hat{E}$, such that $L_{hi}\hat{E} \simeq E$.
\end{definition}

Putting the definition and corollary together, we get the final result to generate examples of interest.

\begin{corollary}
 Let $E$ be a cohomology theory (spectrum). The differential refinement $\hat{E}$ is uniquely determined by the following pieces of data:
\begin{enumerate}
  \item A pure sheaf $E_{\mathrm{pure}}$.
  \item A morphism of spectra $E \to L_{hi}E_{\mathrm{pure}}$.
\end{enumerate}
\end{corollary}

These results suggests the following algorithm for constructing differential cohomology theories of interest:
\begin{enumerate}
  \item Pick a pure sheaf $E_{\mathrm{pure}}$, that encodes geometric data of interest.
  \item Compute the spectrum $L_{hi}E_{\mathrm{pure}}$.
  \item Then for an arbitrary spectrum $E$, every map of spectra $E \to L_{hi}E_{\mathrm{pure}}$ gives rise to a differential refinement of $E$.
\end{enumerate}

Let us implement this algorithm in practice. As we saw, the key example a pure sheaf is $\Omega^{k \geq}$, the sheaf of $k$-truncated forms. We now have the following result due to \cite{bunkenikolausvoelkl2016diffcoh}.
% With these results at hand, we can look at non-trivial examples of interest. The key step is hence to pick a pure sheaf and compute its associated $\bR$-invariant sheaf.

\begin{theorem}
  Let $\Omega^{k \geq}$ be the pure sheaf of $k$-truncated forms. Then, $L_{hi}\Omega^{k \geq} \simeq H\bR$. 
\end{theorem}

\begin{example}
 Let $E$ be a spectrum with a map of spectra $E \to H\bR$. Then there is a differential refinement, $\hat{E}$, given via the pullback square 
 \[ 
 \begin{tikzcd}
 \hat{E} \arrow{r} \arrow{d} & \Omega^{k \geq}  \arrow{d} \\
 E \arrow{r} & H\bR
 \end{tikzcd}
 \]
\end{example}

\begin{example}
 Let us see an example of the example. Let $E = H\bZ$, i.e. singular cohomology, and $H\bZ \to H\bR$ the evident inclusion. Then the resulting differential refinement of singular cohomology recovers \emph{Deligne cohomology} $\hat{\bZ}[l]$. It has the property that the $k$-th sheaf cohomology group recovers the classical Deligne cohomology group $\hat{H}^k(M)$, however at other degrees we get trivial groups. This means the classical Deligne cohomology groups sit inside short exact sequences.
\end{example}

So, in hindsight the fact that historically people have been looking at maps into $H\bR$ is not a coincidence, every differential cohomology theory whose pure part is $\Omega^{k \geq}$ is obtained this way.

\begin{question} \label{que:deligne}
 Given the example, here is a natural question one might wonder. Is there a single differential cohomology theory, such that its $k$-th sheaf cohomology group recovers the classical Deligne cohomology group $\hat{H}^k(M)$ for all $k \geq 0$? It seems the answer would have to be no, because one single cohomology theory would not decompose into a collection of short exact sequences.
\end{question}


\section{Future Directions}
For future talks we can consider the following topics.

\subsection{Fancy Definition of Cohomology Groups}

As we saw in \cref{que:deligne}, we might wonder whether it is possible to recover all Deligne cohomology groups out of one differential cohomology. 

While this might not be possible with the naive regular definition of cohomology groups, there appears to be an alternative definition due to Bunke--Gepner \cite{bunkegepner2021diffktheory}, which is able to recover all Deligne cohomology groups, via a refined method.

Understanding this approach can help us understand ways to extract cohomological data in a non-formal way that is more helpful in geometrically motivated applications.

\subsection{Further Examples of Differential Cohomology Theories}
We already saw Deligne cohomology as a differential refinement of singular cohomology. However, many other cohomologies also admit differential refinements, that merit further study.
\begin{enumerate}
  \item We could look at differential $K$-theory. It is understood as the differential refinement of $ku$, the $K$-theory spectrum. Following the algorithm we want:
  \begin{enumerate}
   \item A pure sheaf $\Omega^{\geq k}(- ;\mathbb{C}[u^{\pm 1}])$.
   \item The computation $L_{hi}\Omega^{\geq k}(- ;\mathbb{C}[u^{\pm 1}]) \simeq H\bC[u^{\pm 1}]$.
   \item A map of spectra $ku \to H\bC[u^{\pm 1}]$, which is the \emph{Chern character}.
  \end{enumerate}
  This was studied with classical means by Hopkins--Singer \cite{hopkinssinger2005diffcoh}.
  \item There are further differential refinements, such as \emph{differential algebraic $K$-theory} \cite{bunkegepner2021diffktheory} or \emph{differential complex cobordism} \cite{bunkeschickschroederwiethaup2009landweber}, however, they seem to have been studied before the sheaf-theoretic framework was developed.
\end{enumerate}

\subsection{Abstract Proof and further Applications of Fracture Square}
One question is whether we want to explicitly go through the proof of the fracture square, and the fact that $\bR$-invariant sheaves are precisely spectra. This involves understanding advanced aspects of sheaf theory, such as \emph{recollements}.

The essential step in the proof is the following very technical result.

\begin{proposition}
  Assume we have the following data and assumptions:
  \begin{enumerate}
    \item A Grothendieck site $(\lC,J)$, such that $\lC$ has a terminal object.
    \item A stable $\infty$-category $\lT$.
    \item The inclusion functor $\Delta\colon \lT \to \Shv_\lT(\lC,J)$ is fully faithful and admits a left adjoint $L_{const}$. 
  \end{enumerate}
  Then the following holds:
  \begin{enumerate}
    % \item $\Delta$ admits a right adjoint $R_{const}$, given by evaluation at the terminal object.
    \item The full sub-category of $\Shv_\lT(\lC,J)$ consisting of sheaves $P$, such that $P(*)$ is the point admits a left adjoint $L^{\bot}\colon \Shv_\lT(\lC,J) \to \Shv_\lT(\lC,J)^\bot$.
    \item For every sheaf $P$ in $\Shv_\lT(\lC,J)$, there is a pullback square
    \[
    \begin{tikzcd}
      P \arrow[r] \arrow[d] & L^{\bot}P \arrow[d] \\
      L_{const}P \arrow[r] & L_{const}L^{\bot}P
    \end{tikzcd},
    \]
    inside $\Shv_\lT(\lC,J)$, where the inclusions are left implicit.
  \end{enumerate} 
\end{proposition}

This is a very general result. We use it for the following specific case:
\begin{lemma}
 Let $(\Euc,J)$ be the Euclidean site and $\Sp$ the stable $\infty$-category of spectra. Then the inclusion functor $\Delta\colon \Sp \to \Shv_{\Sp}(\Euc,J)$ is given by the constant presheaf functor.
\end{lemma}

In other words, the constant presheaf is already a sheaf. This directly implies that $\Delta$ is fully faithful and that it admits a left adjoint via colimit. So we can directly apply the result above to get the fracture square.

Beyond looking at the proof, we can also look at manifestations of this general result in the context of other sites, such as:
\begin{enumerate}
  \item The site of topological spaces.
  \item The site of coarse spaces.
  \item The site of diffeological spaces.
  \item The site of Lie groupoids.
\end{enumerate}

Concretely, we can pursue the following questions:
\begin{enumerate}
  \item Do these Grothendieck sites satify the assumptions of the proposition?
  \item If yes, what does the existence of a fracture square imply in these cases?
  \item If not, what are are the obstructions inherent to the theory?
\end{enumerate}

\subsection{Twisted (differential) cohomology}
Twisted cohomology theories further generalize cohomology theories. They have been refined to a differential version by Bunke--Nikolaus \cite{bunkenikolaus2019twisted}. Given their myriad applications, both twisted cohomology and its differential version merit a careful analysis.

\subsection{Applications to Physics}
Twisted differential cohomology theories have found concrete applications in physics, which we should explore.

\subsection{Further Differential Refinements of Cohomological Notions}
Twisted K-theory has known connections to the concept of \emph{loop group representations} \cite{freedhopkinsteleman2011loopgroupsi}. While twisted K-theory has successfully been lifted to a differential cohomology theory, the loop group representations have not. This is an interesting open problem that should be further explored.


\subsection{Differential Cohomology in Cohesive \texorpdfstring{$\infty$}{oo}-Topoi}
The collective work of Grady, Sati, Schreiber, et al. has focused on developing differential cohomology theories in the context of cohesive $\infty$-topoi and study their applications to physics \cite{schreiber2013cohesive,fiorenzasatischreiber2024charactermap}, which merits further study.

\chapter{Seminar Plan \texorpdfstring{\\}{ } \normalsize \emph{Talk by Nima Rasekh}}

This talk is a quick review and plans for this semester.

% In this talk we summarize what we covered until now and discuss possible future directions.

\section{Summary}
In our last talk we already did a rather detailed summary of what we covered until now \cite{rasekh2025diffcoh5}, so here do a quick bullet point summary.
\begin{enumerate}
  \item Differential cohomology theories should be defined as sheaves valued in the $\infty$-category of spectra valued on the site of manifolds. 
  \item We can use abstract $\infty$-topos theoretic properties to prove that every differential cohomology theory decomposes into an $\bR$-invariant part (a purely homotopical differential cohomology) and a pure part (a purely geometric differential cohomology), via a pullback square called the \emph{fracture square}. 
  \item This fracture square relates the original definition to the ``hexagon approach'' going back to Simons--Sullivan and others \cite{simonssullivan2008diffcoh,bunkeschick2009smoothk}.
  \item Using this abstract approach we can construct examples of interest, such as Deligne cohomology, which refines integral homology.
\end{enumerate}

\section{Talk Topics}
Here is an overview of the possible talk topics we can cover this semester. There is some dependency, but many talks are independent of each other giving us significant flexibility, if we want to rearrange things.

\[
\begin{tabular}{|l|l|l|l|}
  \hline 
  \textbf{Topic} & \textbf{Section} & \textbf{Exp.~vs.~Research} & \# \textbf{Talks}\\ \hline
  Further Examples of Differential Cohomology Theories & \cref{subsec:examples} & Expository & 1 \\ \hline
  Twisted (differential) cohomology and applications & \cref{subsec:twisted} & Expository & 2\\ \hline
  Differential Cohomology in Cohesive $\infty$-Topoi & \cref{subsec:cohesion} & Expository & 2\\ \hline
  Differential Cohomology of Lie Groupoids & \cref{subsec:liegroupoids} & Research & 1\\ \hline
  Differential Refinements of Loop Group Representations & \cref{subsec:loopgroup} & Research & 2\\ \hline
  Cohomology Groups out of Differential Cohomology Theories & \cref{subsec:groups} & Research & 2\\ \hline
  Abstract Proof and further Applications of the Fracture Square & \cref{subsec:abstract} & Research & 2\\ \hline
\end{tabular}
\]
% \begin{enumerate}
%   \item Further Examples of Differential Cohomology Theories (\cref{subsec:examples})
%   \item Twisted (differential) cohomology and applications (\cref{subsec:twisted})
%   \item Differential Refinements of Loop Group Representations (\cref{subsec:loopgroup})
%   \item Differential Cohomology in Cohesive \texorpdfstring{$\infty$}{oo}-Topoi (\cref{subsec:cohesion})
%   \item Cohomology Groups out of Differential Cohomology Theories (\cref{subsec:groups})
%   \item Abstract Proof and further Applications of the Fracture Square (\cref{subsec:abstract})
% \end{enumerate}

Here are further details on each topic.

\subsection{Further Examples of Differential Cohomology Theories} \label{subsec:examples}
We already saw how to can use the methods of the fracture square to construct $\hat{\bZ}[l]$. However, the literature contains many other examples, constructed with different methods. How do these examples fit into our framework? Concrete examples are:
\begin{enumerate}
  \item Differential $K$-theory, as a differential refinement of $ku$  \cite{hopkinssinger2005diffcoh}.
  \item  \emph{differential algebraic $K$-theory} as a differential refinement of algebraic $K$-theory \cite{bunkegepner2021diffktheory}
  \item \emph{differential complex cobordism} as a differential refinement of complex cobordisms \cite{bunkeschickschroederwiethaup2009landweber}
\end{enumerate}

Beyond these already considered examples, it would be interesting if we could find some content about differential refinements of tmf.

\subsection{Twisted (differential) cohomology and applications} \label{subsec:twisted}
Twisted cohomology theories further generalize cohomology theories. They have been refined to a differential version. Given their myriad applications, both twisted cohomology and its differential version merit a careful analysis.

This topics would involve two parts (and hence probably two talks):
\begin{enumerate}
  \item The development of twisted cohomology theory, and its applications in geometry and physics \cite{rosenberg2024twistedcohomology}. 
  \item The differential refinement due to Bunke--Nikolaus, and its applications \cite{bunkenikolaus2019twisted}.
\end{enumerate}

\subsection{Differential Cohomology in Cohesive \texorpdfstring{$\infty$}{oo}-Topoi} \label{subsec:cohesion}
Schreiber has introduced a notion of differential cohomology in the abstract concept of a cohesive $\infty$-topos \cite{schreiber2013differentialcohomology}. This has then be applied to physical settings \cite{fiorenzasatischreiber2024charactermap}. This suggest a $2$-part talk with the following topics:

\begin{enumerate}
  \item Background on $\infty$-topos theory and cohesion, introducing differential cohomology in this setting, comparison to our approach \cite{lurie2009htt,schreiber2013differentialcohomology}
  \item Applications to physics \cite{fiorenzasatischreiber2024charactermap}
\end{enumerate}

\subsection{Differential Cohomology of Lie Groupoids} \label{subsec:liegroupoids}
While differential cohomology theories of manifolds are well-defined and studied extensively, the same cannot be said for Lie groupoids. Concretely, we can imagine several definitions in this context.

\begin{enumerate}
  \item We can embed Lie groupoids into sheaves of manifolds, and then define cohomology via left Kan extension. 
  \item We can think of a Lie groupoid as a special kind of simplicial manifold, and then define differential cohomology point-wise and then geometric realization.
  \item In certain cases, such as differential $K$-theory, we can use classical methods to define differential cohomology (via vector bundles with connection).
\end{enumerate}

Currently, it is unknown how these definitions relate to each other, and what properties they have, suggesting a talk analyzing these questions.

\subsection{Differential Refinements of Loop Group Representations} \label{subsec:loopgroup}
Building on work of Freed and Hopkins, twisted K-theory has known connections to the concept of \emph{loop group representations}.  As a follow up to the talks in \cref{subsec:twisted}, we can explore the following question:

\begin{question}
 Is there a differential analogue of loop group representations and their connection to twisted differential K-theory?
\end{question}


\subsection{Cohomology Groups out of Differential Cohomology Theories} \label{subsec:groups}
Recall that the Deligne cohomology groups do not arise as the homotopy groups of a single differential cohomology theory. Instead, there is a collection of differential cohomology theories, one for each degree, whose homotopy groups give the Deligne cohomology groups.

On the other hand, there are alternative methods to extract cohomology groups out of a single differential cohomology theory, as proposed by Bunke--Gepner \cite[Definition 2.23]{bunkegepner2021diffktheory}. This naturally results in the following question and possible talk topic:

\begin{question}
 Can we use the methods from \cite{bunkegepner2021diffktheory} to extract Deligne cohomology out of $\hat{\bZ}[1]$?
\end{question}

Understanding this approach can help us understand ways to extract cohomological data in a non-formal way that is more helpful in geometrically motivated applications.

\subsection{Abstract Proof and further Applications of the Fracture Square} \label{subsec:abstract}
One question is whether we want to explicitly go through the proof of the fracture square, and the fact that $\bR$-invariant sheaves are precisely spectra. This involves understanding advanced aspects of sheaf theory, such as \emph{recollements}.

The essential step in the proof is the following very technical result.

\begin{proposition}
  Assume we have the following data and assumptions:
  \begin{enumerate}
    \item A Grothendieck site $(\lC,J)$, such that $\lC$ has a terminal object.
    \item A stable $\infty$-category $\lT$.
    \item The inclusion functor $\Delta\colon \lT \to \Shv_\lT(\lC,J)$ is fully faithful and admits a left adjoint $L_{const}$. 
  \end{enumerate}
  Then the following holds:
  \begin{enumerate}
    % \item $\Delta$ admits a right adjoint $R_{const}$, given by evaluation at the terminal object.
    \item The full sub-category of $\Shv_\lT(\lC,J)$ consisting of sheaves $P$, such that $P(*)$ is the point admits a left adjoint $L^{\bot}\colon \Shv_\lT(\lC,J) \to \Shv_\lT(\lC,J)^\bot$.
    \item For every sheaf $P$ in $\Shv_\lT(\lC,J)$, there is a pullback square
    \[
    \begin{tikzcd}
      P \arrow[r] \arrow[d] & L^{\bot}P \arrow[d] \\
      L_{const}P \arrow[r] & L_{const}L^{\bot}P
    \end{tikzcd},
    \]
    inside $\Shv_\lT(\lC,J)$, where the inclusions are left implicit.
  \end{enumerate} 
\end{proposition}

This is a very general result. We use it for the following specific case:
\begin{lemma}
 Let $(\Euc,J)$ be the Euclidean site and $\Sp$ the stable $\infty$-category of spectra. Then the inclusion functor $\Delta\colon \Sp \to \Shv_{\Sp}(\Euc,J)$ is given by the constant presheaf functor.
\end{lemma}

In other words, the constant presheaf is already a sheaf. This directly implies that $\Delta$ is fully faithful and that it admits a left adjoint via colimit. So we can directly apply the result above to get the fracture square.

Beyond looking at the proof, we can also look at manifestations of this general result in the context of other sites, such as:
\begin{enumerate}
  \item The site of topological spaces.
  \item The site of coarse spaces.
  \item The site of diffeological spaces.
  \item The site of Lie groupoids.
\end{enumerate}

Concretely, we can pursue the following questions:
\begin{enumerate}
  \item Do these Grothendieck sites satify the assumptions of the proposition?
  \item If yes, what does the existence of a fracture square imply in these cases?
  \item If not, what are are the obstructions inherent to the theory?
\end{enumerate}


% \subsection{Applications to Physics}
% Twisted differential cohomology theories have found concrete applications in physics, which we should explore.
  
\section{Talk Dates}
Here are some first talk date plans
\[
 \begin{tabular}{|l|l|l|l|}
  \hline 
  & \textbf{Date} & \textbf{Speaker} & \textbf{Topic} \\ \hline 
  (0) & 14.10.2025 & Nima Rasekh & Review and Talk Distribution \\ \hline
  (1) & 29.10.2025 & Matthias Ludewig & Further Examples\\ \hline
  (2) & 05.11.2025 & Hannes Berkenhagen & Twisted Cohomology Theories\\ \hline
  (3) & 12.11.2025 & Alessandro Nanto & Twisted Diff. Cohomology Theories\\ \hline
  (4) & 19.11.2025 & Matthias Frerichs & $\infty$-Topoi and Cohesion\\ \hline
  (5) & 26.11.2025 & Nima Rasekh & Differential Cohomology and Cohesion\\ \hline
  (6) & 03.12.2025 & Christian Becker & Differential Cohomology of Lie Groupoids \\ \hline
  % (7) & 10.12.2025 & & \\ \hline
  % (8) & 17.12.2025 & & \\ \hline
  % (9) & 07.01.2026 & & \\ \hline
  % (10) & 14.01.2026 & & \\ \hline
  % (11) & 21.01.2026 & & \\ \hline
  % (12) & 28.01.2026 & & \\ \hline
 \end{tabular} 
\]

\chapter{Previously and Differential K-Theory \texorpdfstring{\\}{ } \normalsize \emph{Talk by Matthias Ludewig}}

\section{Recollection: Differential cohomology}
A differential cohomology theory is a sheaf $\hat{E}$ on manifolds taking values in some stable $\infty$-category $\lC$, usually chain complexes or spectra. It is called \emph{$\bR$-invariant} or \emph{homotopy invariant} if the projection map $M \times \bR \to M$ induces an equivalence. In this case, it is the sheafification of the presheaf constant equal to $\hat{E}(*)$, in formulas $\smash{\hat{E} \simeq \underline{\hat{E}(*)}}$. A differential cohomology theory is called \emph{pure} if $\hat{E}(*) \simeq 0$.

\subsection{The basic maps}
To begin with, the inclusion of homotopy invariant into all sheaves has both a left and right adjoint, 
\[
\begin{tikzcd}[row sep=0.2cm]
L_{hi} : \Shv(\Mfd, \lC)
\ar[r, shift left=1]
&
\Shv^{hi}(\Mfd, \lC)
\ar[l, shift left=1] : \mathrm{incl}
\\
\mathrm{incl} : \Shv^{hi}(\Mfd, \lC)
\ar[r, shift left=1]
&
\Shv(\Mfd, \lC)
\ar[l, shift left=1] : R_{hi}.
\end{tikzcd} 
\]
The right adjoint evaluates a differential cohomology theory at the point (which is an object $\hat{E}(*)$ of $\mathscr{C}$) and then sheafifying the constant presheaf given by $\hat{E}(*)$. We may then take the fiber, respectively cofiber of the canonical maps of these adjunctions,
\[
\begin{tikzcd}[row sep=0.2cm]
\mathrm{Def}(\hat{E}) \ar[r, "\mathrm{def}"] & \hat{E} \ar[r] & L_{hi}(\hat{E})
\\
R_{hi}(\hat{E}) \ar[r] & \hat{E} 
\ar[r, "\mathrm{curv}"] & \mathrm{Cyc}(\hat{E}),
\end{tikzcd}
\]
where we omitted the inclusion functor of homotopy invariant sheaves into all sheaves in notation. These are called the \emph{homotopification}, respectively \emph{geometric cycles} functors. Since point evaluation is a right adjoint (to the functor that sheafifies the constant sheaf given by an object of $\lC$), it preserves fiber/cofiber sequences, hence we see that $\mathrm{Cyc}(\hat{E})$ is a pure sheaf. It turns out that it is in fact the left adjoint of the inclusion functor from pure sheaves
\[
\begin{tikzcd}[row sep=0.2cm]
\mathrm{Cyc} : \Shv(\Mfd, \lC)
\ar[r, shift left=1]
&
\Shv^{pure}(\Mfd, \lC)
\ar[l, shift left=1] : \mathrm{incl}.
\end{tikzcd} 
\]
Moreover, the restriction of $\mathrm{Def}$ to pure sheaves turns out to be left adjoint to the $\mathrm{Cyc}$ functor, 
\[
\begin{tikzcd}
\mathrm{Def} : \Shv^{pure}(\Mfd, \lC)
\ar[r, shift left=1]
&
\Shv(\Mfd, \lC)
\ar[l, shift left=1] : \mathrm{Cyc}.
\end{tikzcd} 
\]
The functor $\mathrm{Def}$ and $\mathrm{Cyc}$ have the properties that
\[
L_{hi}\mathrm{Def}(\hat{E}) \simeq 0 \qquad \text{and} \qquad R_{hi}\mathrm{Cyc}(\hat{E}) \simeq 0.
\]
See \cite[\S3]{bunkenikolausvoelkl2016diffcoh} for the general discussion of these results.

\subsection{Differential cohomology diamond}
We may organize these data into a $3 \times 3$ grid, where every square is a pullback square,
\begin{equation*}
\begin{tikzcd}
\Sigma^{-1} L_{hi}\mathrm{Cyc}(\hat{E}) 
\ar[r] \ar[d] 
& 
R_{hi}(\hat{E})
\ar[r] \ar[d] 
&
0
 \ar[d] 
\\
\mathrm{Def}(\hat{E})
\ar[r] \ar[d] 
&
\hat{E}
\ar[r] \ar[d] 
&
\mathrm{Cyc}(\hat{E})
\ar[d] 
\\
0 \ar[r]
&
L_{hi}(\hat{E}) \ar[r, "\phi"] &
L_{hi}\mathrm{Cyc}(\hat{E}).
\end{tikzcd}
\end{equation*}
Note that the short sequences in the middle are \emph{not} fiber sequences; instead the composition yields interesting maps. We may twist the diagram in the middle and turn it on the side, to obtain the diamond shaped diagram
\begin{equation*}
\begin{tikzcd}[column sep=0.5cm]
& 
R_{hi}(\hat{E}) 
\ar[dr]\ar[rr]
& & 
L_{hi}(\hat{E}) 
\ar[dr, "\phi"] 
&
\\ 	
\Sigma^{-1} L_{hi}\mathrm{Cyc}(\hat{E})
\ar[ur]
\ar[dr]
& & 
\hat{E} 
\ar[ur]
\ar[dr, "\mathrm{curv}"]
& & 
L_{hi}\mathrm{Cyc}(\hat{E})
\\
& 
\mathrm{Def}(\hat{E})
\ar[ur, "\mathrm{def}"]
\ar[rr, "d"']
& & 
\mathrm{Cyc}(\hat{E}) 
\ar[ur]
&
\end{tikzcd}
\end{equation*}
which has the property that the two diagonal sequences are fiber sequences, and that the top and bottom row each are fiber sequences as well (by using the pasting law for Cartesian squares on the previous diagram). The naming of the diagram $d$ comes from the example of differential forms, where it corresponds to the exterior derivative.

\subsection{Fracture square}
Any differential cohomology theory $\hat{E}$ is determined by its cycles $\mathrm{Cyc}(\hat{E})$ (which is pure), its homotopification $L_{hi}(\hat{E})$ (which is homotopy invariant), and the map
\[
\phi : L_{hi}(\hat{E}) \longrightarrow L_{hi}\mathrm{Cyc}(\hat{E}),
\]
which by homotopy invariance is fully determined on the underlying map in $\lC$ over the point. \cite{bunkenikolausvoelkl2016diffcoh} call this map the \emph{characteristic map}.

Conversely, if we are given an object $E$ of $\lC$, a pure sheaf $P \in \Shv^{pure}(\Mfd, \lC)$ and a map $\phi: E \to L_{hi}(P)$ in $\lC$, then we may define a differential cohomology theory as the pullback
\[
\begin{tikzcd}
	\hat{E} \ar[r] \ar[d] & P
	\ar[d]
	\\
	E \ar[r, "\phi"] & L_{hi}(P).
\end{tikzcd}
\]
This theory satisfies
\[
L_{hi}(\hat{E}) = E, \qquad \mathrm{Cyc}(\hat{E}) = P, 
\]
and the map $\phi$ canonically extracted from $\hat{E}$ agrees with the previously given map $\phi$.


\section{Basic example: Differential forms with values in a chain complex}
Consider the sheaf
\[
\Omega : \Mfd^{\mathrm{op}} \longrightarrow \Ch_\bR
\]
of cochain complexes that assigns to a manifold $M$ the cochain complex $\Omega(M)$ of differential forms on $M$. If $C \in \Ch_\bR$ is any real cochain complex and $n$ is a natural number, we may define a new sheaf $\Omega \otimes_\bR C$, given by
\[
\Omega^n(M, C) := (\Omega \otimes_\bR C)^n(M) = \bigoplus_{p + q = n} \Omega^p(M) \otimes_{\bR} C^q,
\]
as well as its truncations $(\Omega \otimes_{\bR} C)^{\geq n}$, which are obtained by setting the complex to zero of $k < n$ over each manifold. Inverting quasi-isomorphisms, this remains a sheaf\footnote{Why again?} which we again denote by $(\Omega \otimes_{\bR} C)^{\geq n}$, valued in the derived category $\Ch_\bR[W^{-1}]$ of chain complexes.
 
\begin{theorem}\cite[Lemma~4.4]{bunkenikolausvoelkl2016diffcoh}
With $\hat{E} = (\Omega \otimes_\bR C)^{\geq n}$, we have
 \begin{align*}
R_{hi}(\hat{E})(*) &\simeq C^{\geq n}
\\
L_{hi}(\hat{E})(*) &\simeq C
\\
L_{hi}\mathrm{Cyc}(\hat{E})(*) &\simeq C^{\leq n-1}
\\
\mathrm{Cyc}(\hat{E}) &\simeq (\Omega \otimes_\bR C^{\leq n-1})^{\geq n})
\\
\mathrm{Def}(\hat{E}) &\simeq \Sigma \bigl(\Omega \otimes_\bR C)^{\leq n-1}\bigr).
 \end{align*}	
Moreover, the characteristic map $\phi$ is the truncation map $C \to C^{\leq n-1}$.
 \end{theorem}
 
We determine the differential cohomology diamond for this theory.
 
 \medskip

Given a chain complex $C \in \Ch_\bR$, there is a unique homotopy invariant differential cohomology theory $E_C$ valued in chain complexes. To calculate the differential cohomology diamond in this case, we have to understand what this differential cohomology theory assigns to an arbitrary manifold $M$. The answer\footnote{Maybe better: One answer, as the theory is only determined up to equivalence of chain complexes. But this is the answer that best fits our current considerations.} is
\[
E_C(M) = \Omega(M) \otimes_\bR C,
\]
the tensor product chain complex. One can see this from the fact that this is a differential cohomology theory which is homotopy invariant (by the Poincar\'e lemma) and clearly evaluates to $C$ on the point. The cohomology groups of this chain complex will be denoted by $H^n(M, C)$. We moreover observe that 
\[
H^n(R_{hi}(E)(M)) = H^0(M, C^n_{\mathrm{cl}}),
\]
where $C^n_{\mathrm{cl}} \subseteq C^n$ denotes the space of closed cochains of degree $n$ in the complex $C$. Indeed, we need to take the cohomology of the complex
\[
0 \longrightarrow \Omega^0(M, C^n) \longrightarrow \Omega^0(M, C^{n+1}) \oplus \Omega^1(M, C^n) \longrightarrow \cdots.
\]
The differential is $D = d\otimes 1 + 1 \otimes \delta$ (with $\delta$ the differential of $C$), so since the components in each of the two summands have to vanish separately, we obtain that $D(f \otimes c) = 0$ if and only if $df$ and $\delta c$ are individually closed.

To determine the non-homotopy invariant terms in the diamond, notice that the chain complex $\mathrm{Cyc}(\hat{E})$ starts with
\[
0 \longrightarrow \bigl(\Omega \otimes_\bR C^{\leq n-1}\bigr)^n \longrightarrow \bigl(\Omega \otimes_\bR C^{\leq n-1}\bigr)^{n+1},
\]
its degree $n$ cohomology is therefore
\[
\Omega_{\mathrm{cl}}^n(M, C^{\leq n-1}) = \left(\bigoplus_{p=2}^\infty \Omega^p(M, C^{n-p})\right)_{\mathrm{cl}} \oplus \Omega^1_{\mathrm{cl}}(M, C^{n-1}).
\]
Notice here that in the first sum, the de Rham differential mixes with the differential of $C$, and the bif sum of the left is generally \emph{not} equal to direct sum of the terms $\Omega^p_{\mathrm{cl}}(M, C^{n-p})$ or $\Omega^p_{\mathrm{cl}}(M, C_{\mathrm{cl}}^{n-p})$, which would mean closed de Rham forms valued in the vector space $C^{n-p}$, respectively $C^{n-p}_{\mathrm{cl}}$).
The last term above comes from the fact that we truncated the complex $C$ also.

The complex $\mathrm{Def}(\hat{E})$ is unbounded from the left and ends at
\[
\cdots \longrightarrow \Omega^{n-2}(M, C) \longrightarrow \Omega^{n-1}(M, C) \longrightarrow 0,
\]
the last non-zero term sitting in degree $n$. The result is therefore
\[
\mathrm{Def}(\hat{E})^n(M) = \Omega^{n-1}(M, C)/d\Omega^{n-2}(M, C).
\]
In total, the differential cohomology diamond in degree $n$ becomes
\begin{equation*}
\begin{tikzcd}[column sep=0.4cm]
&[-1.2cm] 
H^0(M, C^n_{\mathrm{cl}})
\ar[dr]\ar[rr]
& & 
H^n(M, C) 
\ar[dr, "\phi"] 
&
\\ 	
H^{n-1}(M, C^{\leq n-1})
\ar[ur]
\ar[dr]
& & 
\Omega^n_{\mathrm{cl}}(M, C)  
\ar[ur]
\ar[dr, "\mathrm{curv}"]
& & 
H^n(M, C^{\leq n-1})
\\
& 
\Omega^{n-1}(M, C)/\mathrm{im}(d)
\ar[ur, "\mathrm{def}"]
\ar[rr, "D"']
& & 
\Omega^n_{\mathrm{cl}}(M, C^{\leq n-1}) 
\ar[ur]
&
\end{tikzcd}
\end{equation*}

Above, we constructed a differential cohomology theory valued in chain complexes. We may turn this into a spectrum valued differential cohomology theory using the Eilenberg-McLane functor
\[
H : \Ch_\bR[W^{-1}] \longrightarrow \Mod_{H\bR}(\Sp) \subseteq \Sp,
\]
which preserves limits, hence also the sheaf condition. It has the property
\[
\pi_n(HC) = (HC)^{-n}(\mathrm{pt}) = H^{-n}(C),
\]
where the second term denotes the $n$-th $HC$-cohomology of the point and the third term denotes the cohomology of the chain complex $C$. This is enforced by the requirement that for a short exact sequences of cochain complexes, the resulting long exact sequence of cohomology groups is sent to the short exact sequence of homotopy groups.

\section{Differential \texorpdfstring{$K$}{K}-theory}
Following \cite[Def.~3.6]{bunkenikolausvoelkl2016diffcoh} a \emph{differential refinement} of a  cohomology theory $E$ is a differential cohomology theory $\hat{E}$ together with an equivalence $E \simeq L_{hi}(\hat{E})$. In general, there may be many differential refinements of a given cohomology theory $E$. The fracture square
\[
\begin{tikzcd}
	\hat{E} \ar[d]\ar[r] & P
	\ar[d]
	\\
	\underline{E} \ar[r, "\phi"] &L_{hi}\mathrm{Cyc}(P)
\end{tikzcd}
\] 
gives a characterization: A differential refinement of $E$ is ``the same'' as a pure sheaf $P$ together with a map of spectra $E \to L_{hi}\mathrm{Cyc}(P)(*)$. So there are generally plenty differential refinements of a theory and one is interested in ``geometrically interesting'' ones.

For complex $K$-theory, one is interested in differential refinements of $ku$ or $KU$ (the connective, respectively non-connective complex $K$-theory spectrum) which have something to do with vector bundles with connection. There are at least two different constructions in the literature, one given in \cite{hopkinssinger2005diffcoh}, one given in \cite{bunkenikolausvoelkl2016diffcoh}. Before discussing any of these, we give the case of line bundles, where the discussion is cleanest, to get a feeling for what to expect.

\subsection{The case of line bundles}
Consider the functor $\Line^\nabla : \Mfd^{\mathrm{op}} \to \Grpd$ that takes a manifold $M$ to the symmetric monoidal groupoid of line bundles with connection over $M$. This is a stack in symmetric monoidal groupoids on manifolds. Symmetric monoidal groupoids are the same thing as 1-truncated spectra, so we may view this directly as a sheaf of spectra\footnote{In \cite[Eq.~(39)]{bunkenikolausvoelkl2016diffcoh}, they sheafify, but it seems that this is not necessary on $\Mfd$ with values in $\lS$},
\[
\Line^\nabla \in \Shv(\Mfd, \Sp).
\]

\begin{remark}
\label{RemComMon}
This identification is somewhat non-trivial. The problem is that symmetric monoidal categories are \emph{not} monoid objects in the 1-category $\cat$ of categories, so the nerve functor does not take them to monoid objects in $\lS$.
\end{remark}

\begin{theorem}
We have
\begin{align*}
R_{hi}(\Line^\nabla)(*) 
&\simeq \Sigma H\bC^\times
\\
L_{hi}(\Line^\nabla)(*) &\simeq \Sigma H\bC^\times_{\mathrm{top}} 
\simeq \Sigma^2 H\bZ
\\
L_{hi}\mathrm{Cyc}(\Line^\nabla)(*) &\simeq \Sigma H\bC_\delta
\\
\mathrm{Cyc}(\Line^\nabla) &\simeq \Sigma H(\Omega^{\geq 2} \otimes_\bR \bC)
\\
\mathrm{Def}(\Line^\nabla) &\simeq \Sigma H(\Omega^{\leq 1} \otimes_\bR \bC ).
\end{align*}
The characteristic map is given by $\Sigma H\bC^\times \simeq \Sigma^2 H\bZ \to \Sigma^2 H \bC$, where the first is the boundary map of the exponential sequence and the second is induced by the inclusion map.
\end{theorem}

\begin{proof}
The first two statements are \cite[Lemma~5.2]{bunkenikolausvoelkl2016diffcoh}, observing that $\Line^\nabla \cong \Bdl^\nabla(\bC^\times)$. If $\Line^\nabla$ meant ``Hermitian line bundles with compatible connection'', then the last three statements would literally follow from \cite[Lemma 5.4]{bunkenikolausvoelkl2016diffcoh} combined with \cite[Lemma 4.5]{bunkenikolausvoelkl2016diffcoh} (where the right hand side does not involve tensoring with $\bC$). In our current setting, the results have to be slightly adapted.
\end{proof}

The differential cohomology diamond is most interesting in degree zero; in all other degrees, it is purely topological. There we get
\begin{equation*}
\begin{tikzcd}[column sep=0.4cm]
&[-0.2cm] 
H^1(M, \bC^\times)
\ar[dr]\ar[rr]
&[-1cm] & 
H^2(M, \bZ)
\ar[dr, "\phi"] 
&
\\ 	
H^1(M, \bC)
\ar[ur, "\exp_*"]
\ar[dr]
& & 
\pi_0\bigl(\Line^\nabla(M)\bigr)
\ar[ur, "\ch"]
\ar[dr, "\mathrm{curv}"]
& & 
H^2(M, \bC)
\\
& 
\Omega^{1}(M, \bC)/\mathrm{im}(d)
\ar[ur, "\mathrm{def}"]
\ar[rr, "d"']
& & 
\Omega^2_{\mathrm{cl}}(M, \bC) 
\ar[ur]
&
\end{tikzcd}
\end{equation*}
Here the diagonal map starting at $H^1(M, \bC^\times) \cong \Hom(\pi_1(M), \bC^\times)$ is the inclusion of flat line bundles, the map $\mathrm{def}$ sends a class $[\alpha]$ to the trivial line bundle with connection $d + \alpha$ (here changing $\alpha$ to $\alpha' = \alpha + df$ does not change the isomorphism class, as multiplication by $e^f$ is a gauge transformation from $d + \alpha$ to $d+ \alpha'$). The curvature map takes the curvature 2-form.

\subsection{Hopkins--Singer differential \texorpdfstring{$K$}{K}-theory.}
In \cite[\S{4.4}]{bunkenikolausvoelkl2016diffcoh}, Bunke--Nikolaus--V{\"o}lkl give the following construction, which they implicitly claim produces the same result as that of Hopkins-Singer \cite{hopkinssinger2005diffcoh}, although this is not clear to me.

The construction more generally produces differential refinements of arbitrary cohomology theories by mixing with differential forms. The input is a spectrum $E$, a number $n \in \bZ$, a chain complex $C \in \Ch_{\bR}$ and a map of spectra $c : E \to HC$. Then a differential cohomology theory $\hat{E}$ is defined as the pullback
\[
\begin{tikzcd}
\hat{E} 
\ar[r] \ar[d] & H((\Omega \otimes_\bR C)^{\geq n})\ar[d]
\\
\underline{E} \ar[r, "\underline{c}"]
& \underline{HC}.
\end{tikzcd}
\]

\begin{theorem}\cite[Thm.~4.7]{bunkenikolausvoelkl2016diffcoh}
\label{DiffRefinementThm}
We have
\begin{align*}
	L_{hi}(\hat{E})(*)
	&\simeq E
	\\
	L_{hi}\mathrm{Cyc}(\hat{E})(*) &\simeq H(C^{\leq n-1})
	\\
	\mathrm{Cyc}(\hat{E}) 
	&\simeq H((\Omega \otimes_\bR C^{\leq n-1})^{\geq n})
	\\
	\mathrm{Def}(\hat{E})
	&\simeq \Sigma H((\Omega \otimes_\bR C)^{\leq n-1})
\end{align*}
and $R_{hi}(\hat{E})(*)$ fits into the pullback square
\[
\begin{tikzcd}
	R_{hi}(\hat{E})(*)
	\ar[r]
	\ar[d] &
	H(C^{\geq n})
	\ar[d]
	\\
	E \ar[r, "c"] & HC
\end{tikzcd}
\]
Moreover, the characteristic map of the theory is the composition
\[
\begin{tikzcd}
\phi : E \ar[r, "c"] & HC \ar[r] & H(C^{\leq n-1}),
\end{tikzcd}
\]
where the last map is the truncation map. In particular, $\hat{E}$ is a differential refinement of $E$. 
\end{theorem}

To obtain a differential refinement of complex $K$-theory, we take $E = ku$ the complex connective $K$-theory functor, $C = \bC[x]$, where $x$ is a formal variable of degree $-2$ ($C$ is viewed as a chain complex with a trivial differential), $c$ is the Chern character map
\[
\ch : ku \longrightarrow H(\bC[x])
\]
and the integer $n \in \bZ$ is taken to be zero. In other words, we make the following definition, which specializes the Hopkins--Singer construction to this case.

\begin{definition}
The differential cohomology theory 	$\widehat{ku}_{\mathrm{HS}}$ is defined as the pullback
\[
\begin{tikzcd}
\widehat{ku}_{\mathrm{HS}}
\ar[r, dashed] 
\ar[d, dashed]
&
H\bigl((\Omega \otimes_\bR \bC[x])^{\geq 0}\bigr)
\ar[d]
\\
\underline{ku} \ar[r, "\ch"] 
& 
\underline{H(\bC[x])}.
\end{tikzcd}
\]
\end{definition}

To describe $R_{hi}(\widehat{ku}_{\mathrm{HS}})$, we need $\bC^\times$-valued $K$-theory, which is defined as the Chern cofiber of the Chern character map $\ch : KU \to H(\bC[x, x^{-1}])$. Its connective version is defined as the cofiber of the connective Chern character map,
\[
ku_{\bC^\times} := \mathrm{cofib}\bigl(ku \stackrel{\ch}{\longrightarrow} H(\bC[x])\bigr), \quad \text{equivalently}
\quad
\Sigma^{-1} ku_{\bC^\times} = \mathrm{fib}\bigl(ku \stackrel{\ch}{\longrightarrow} H(\bC[x])\bigr).
\]
We may also define a spectrum $\widetilde{ku}_{\bC^\times}$ in the same way, but using the reduced Chern character, 
\[
\widetilde{ku}_{\bC^\times} := \mathrm{cofib}\bigl(ku \stackrel{\widetilde{\ch}}{\longrightarrow} H(x\bC[x])\bigr), \quad \text{equivalently}
\quad
\Sigma^{-1} \widetilde{ku}_{\bC^\times} = \mathrm{fib}\bigl(ku \stackrel{\widetilde{\ch}}{\longrightarrow} H(x\bC[x])\bigr).
\]
As the cofiber of connective spectra, both these spectra are again connective. In non-positive degrees, the $ku$-cohomology groups coincide with the classical $K$-theory groups, and similarly, the non-positive $ku_{\bC^\times}$-cohomology groups coincide with the $\bC^\times$-valued cohomology groups considered in Karoubi \cite{karoubi1987cyclichomology} or Lott \cite{lott1994rzindextheory}. The long exact sequence of homotopy groups yields that the homotopy groups are given by
\[
\pi_n(ku_{\bC^\times}) = \begin{cases}
 \bC^\times	& n \geq 0 ~\text{even}
 \\
 0 & \text{else}.
 \end{cases}
 \qquad
 \pi_n(\widetilde{ku}_{\bC^\times}) = \begin{cases}
 \bC^\times	& n \geq 2 ~\text{even}
 \\
 \bZ & n=1
 \\
 0 & \text{else}.
 \end{cases}
\]
In other words, the the comparison map $ku_{\bC^\times} \to \widetilde{ku}_{\bC^\times}$ induces an isomorphism on homotopy groups. We denote by $K^n(M, \bC^\times)$, $k^n(M, \bC^\times)$, respectively $\tilde{k}^n(M, \bC^\times)$ the cohomology groups of the non-connective, connective, respectively reduced connective $\bC^\times$-valued $K$-theory. We have
\begin{align*}
K^n(M, \bC^\times) \cong k^n(M, \bC^\times)\rlap{\qquad$n\leq -1$}
\\ k^n(M, \bC^\times) \cong \tilde{k}^n(M, \bC^\times) \rlap{\qquad$n\leq -2$}
\end{align*}
The Atiyah-Hirzebruch spectral sequence should then yield that there is a short exact sequence
\[
\begin{tikzcd}
0
\ar[r]
&
k^{-1}(M, \bC^\times)	
\ar[r]
&
\tilde{k}^{-1}(M, \bC^\times) \ar[r] 
& 
H^0(M, \bZ)
\ar[r]
&
0.
\end{tikzcd}
\]

\begin{remark}
The reason for this terminology is that due to the fact that the Chern character is an isomorphism after tensoring with $\bC$, $H(\bC[x)$ may be viewed as ``$\bC$-valued $K$-theory'', and the cofiber of the map from $\bZ$-valued $K$-theory to $\bR$-valued $K$-theory deserves to be named ``$\bC^\times$-valued $K$-theory''.
\end{remark}

\begin{theorem}
We have
\begin{align*}
R_{hi}(\widehat{ku}_{\mathrm{HS}})(*) &\simeq \Sigma^{-1} ku_{\bC^\times}
\\
L_{hi}(\widehat{ku}_{\mathrm{HS}})(*)
&\simeq ku
\\
L_{hi}\mathrm{Cyc}(\widehat{ku}_{\mathrm{HS}})(*)
&\simeq H(x\bC[x])
\\
\mathrm{Cyc}(\widehat{ku}_{\mathrm{HS}})
&\simeq H\bigl((\Omega \otimes_\bR x\bC[x])^{\geq 0}\bigr)
\\
\mathrm{Def}(\widehat{ku}_{\mathrm{HS}})
&\simeq \Sigma H\bigl((\Omega \otimes_\bR \bC[x])^{\leq -1}\bigr)
\end{align*}
The characteristic map is given by the reduced connective Chern character,
\[
\widetilde{\ch} : ku \stackrel{\ch}{\longrightarrow} H(\bC[x]) \longrightarrow H(x\bC[x]).
\]
\end{theorem}

\begin{proof}
All claims follow directly from \cref{DiffRefinementThm} except for the calculation of $R_{hi}(\widehat{ku}_{\mathrm{HS}})(*)$, which is, also by \cref{DiffRefinementThm}, given as a pullback; concretely, it fits into the left square of the following diagram
\[
\begin{tikzcd}
	R_{hi}(\widehat{ku}_{\mathrm{HS}})(*)
	\ar[r]
	\ar[d] &
	H\bC
	\ar[d]
	\ar[r] & * \ar[d]
	\\
	ku \ar[r, "\ch"] & H(\bC[x]) \ar[r] & H(x\bC[x]).
\end{tikzcd}
\]
The second diagram comes from the short exact sequence $\bC \to \bC[x] \to x\bC[x]$ of $\bC$-vector spaces. Since $H$ preserves limits, it is also a pullback diagram. By the pasting law, also the exterior diagram is a pullback diagram, which implies that
\[
R_{hi}(\widehat{ku}_{\mathrm{HS}}) = \mathrm{fib}\bigl(ku \stackrel{\widetilde{\ch}}{\longrightarrow} H(x\bC[x])\bigr).
\]
Notice here that the bottom composition of the above diagram is the reduced Chern character. By definition, this fiber is the de-suspension of $ku_{\bC^\times}$.
\end{proof}

The negative even degree cohomology groups for $H(\bC[x])$ may be identified with the direct sum of all even degree ordinary cohomology groups. The differential cohomology diamond for Hopkins-Singer differential $K$-theory therefore becomes
\begin{equation*}
\begin{tikzcd}[column sep=0cm]
&[-0.2cm] 
\widetilde{k}{}^{-1}(M, \bC^\times)
\ar[dr]\ar[rr]
& &[0.9cm]
K^0(M) 
\ar[dr, "\widetilde{\ch}"] 
&
\\ 	
H^{\mathrm{odd}}(M, \bC)
\ar[ur]
\ar[dr]
& & 
\widehat{ku}{}_{\mathrm{HS}}^0(M) 
\ar[ur]
\ar[dr, "\mathrm{curv}"]
& & 
H^{\mathrm{ev} \geq 2}(M, \bC)
\\
& 
\Omega^{\mathrm{odd}}(M, \bC)/\mathrm{im}(d)
\ar[ur, "\mathrm{def}"]
\ar[rr, "d"']
& & 
\Omega_{\mathrm{cl}}^{\mathrm{ev} \geq 2}(M, \bC)
\ar[ur]
&
\end{tikzcd}
\end{equation*}
%Another diamond, which looks somewhat more reasonable, is
%\begin{equation*}
%\begin{tikzcd}[column sep=0cm]
%&[-0.2cm] 
%K^{-1}(M, \bC^\times)
%\ar[dr]\ar[rr]
%& &[0.9cm]
%K^0(M) 
%\ar[dr, "{\ch}"] 
%&
%\\ 	
%H^{-1}(M, \bC[x])
%\ar[ur]
%\ar[dr]
%& & 
%\widehat{ku}{}_{\mathrm{HS}}^0(M) 
%\ar[ur]
%\ar[dr, "\mathrm{curv}"]
%& & 
%H^0(M, \bC[x]))
%\\
%& 
%\Omega^{-1}(M, \bC[x])/\mathrm{im}(d)
%\ar[ur, "\mathrm{def}"]
%\ar[rr, "d"']
%& & 
%\Omega_{\mathrm{cl}}^0(M, \bC[x])
%\ar[ur]
%&
%\end{tikzcd}
%\end{equation*}

%The Atiyah-Hirzebruch spectral sequence should then yield that there is a short exact sequence
%\[
%\begin{tikzcd}
%0
%\ar[r]
%&
%K^{-1}(M, \bC^\times)	
%\ar[r]
%&
%P^0(M) \ar[r] 
%& 
%H^0(M, \bZ)
%\ar[r]
%&
%0
%\end{tikzcd}
%\]

\subsection{Universal differential \texorpdfstring{$K$}{K}-theory}
In \cite{bunkenikolausvoelkl2016diffcoh}, a differential cohomology theory is defined directly from vector bundles with connection. For each manifold $M$, we obtain a symmetric monoidal groupoid $\Vect^\nabla(M)$ of vector bundles with connection. The functor
\[
M \longmapsto \Vect^\nabla(M)
\]
is a sheaf valued in symmetric monoidal groups on the site $\Mfd$ of manifolds. Taking nerves, we view groupoids as elements of the $\infty$-category $\lS$ of spaces. This promotes $\Vect^\nabla$ to a sheaf valued in the category $\CMon(\lS)$ of commutative monoid objects in spaces.

We may now apply the group completion functor $\texttt{Grp}$ which is the left adjoint in the adjunction
\[
\begin{tikzcd}
\mathrm{grp} : \CMon(\lS)
\ar[r, shift left=1]
&
\CMon^{\mathrm{grp}}(\lS)
\ar[l, shift left=1] : \mathrm{incl},
\end{tikzcd} 
\]
where $\lS$ is the $\infty$-category of spaces and the map from right to left is the inclusion of group-like commutative monoids in $\lS$ into all commutative monoids. Recall further that $\CMon^{\mathrm{grp}}(\mathcal{S})
\simeq \Sp_{\geq 0},$
the $\infty$-category of connective spectra. However, group completion, as a left adjoint functor, only preserves colimits and not generally limits, hence applying it to the sheaf $M \mapsto \Vect^\nabla(M)$ generally destroys the sheaf condition.

\begin{definition}
We define $\widehat{ku}{}_{\mathrm{uni}}$ as the sheafification of the functor 
\[
\Mfd^{\mathrm{op}} \longrightarrow \CMon^{\mathrm{grp}}(\lS) \simeq \Sp_{\geq 0} \subseteq \Sp, \qquad
M \longmapsto \mathrm{grp}(\Vect^\nabla(M))
\]
\end{definition}



\begin{theorem}\cite[Lemma~6.3]{bunkenikolausvoelkl2016diffcoh}
We have
\begin{align*}
R_{hi}(\widehat{ku}_{\mathrm{uni}})(*)
&\simeq K\bC
\\
L_{hi}(\widehat{ku}_{\mathrm{uni}})(*) 
&\simeq ku,
\end{align*}
where $ku$ is the connective complex $K$-theory spectrum and $K\bC$ is the algebraic $K$-theory spectrum of $\bC$.\footnote{There is no non-connective algebraic $K$-theory spectrum of $\bC$ as the algebraic $K$-theory of fields vanishes in negative degrees.} Moreover, the canonical map $\smash{R_{hi}(\widehat{ku}_{\mathrm{uni}}) \to L_{hi}(\widehat{ku}_{\mathrm{uni}})}$ is the comparison map between algebraic and topological $K$-theory.
\end{theorem}

\begin{remark}
In \cite{bunkenikolausvoelkl2016diffcoh}, Lemma~6.3, it is also proved that the same statement also holds when 	$\widehat{ku}_{\mathrm{uni}}$ is replaced by the differential cohomology theory defined using vector bundles without Hermitian metrics and connections.
\end{remark}

\begin{remark}
Both $\mathrm{Def}(\widehat{ku}{}^\nabla_{\mathrm{uni}})$ and $\mathrm{Cyc}(\widehat{ku}{}^\nabla_{\mathrm{uni}})$ seem generally hard to understand. Since $\mathrm{Cyc}(\widehat{ku}{}^\nabla_{\mathrm{uni}})$ is pure, evaluating the right square at the point yields a fiber sequence
\[
K\bC \longrightarrow ku \longrightarrow L_{hi}\mathrm{Cyc}(\widehat{ku}{}^\nabla_{\mathrm{uni}})(*).
\]
We obtain that
\[
L_{hi}\mathrm{Cyc}(\widehat{ku}{}^\nabla_{\mathrm{uni}})(*) \simeq \mathrm{cofib}(K\bC \to ku) \simeq \Sigma\mathrm{fib}(K\bC \to ku) =: \Sigma k^{\mathrm{rel}}\bC,
\]
the relative connective $K$-theory spectrum of $\bC$. Hence $\mathrm{Cyc}(\widehat{ku}{}^\nabla_{\mathrm{uni}})$ is a pure differential refinement of a suspension of the relative $K$-theory spectrum $ k^{\mathrm{rel}}\bC$. 
\end{remark}

The diamond is therefore
\begin{equation*}
\begin{tikzcd}[column sep=0.5cm]
& 
K\bC 
\ar[dr]\ar[rr]
& & 
ku 
\ar[dr, "\phi"] 
&
\\ 	
\Sigma^{-1} k^{\mathrm{rel}}\bC
\ar[ur]
\ar[dr]
& & 
\widehat{ku}_{\mathrm{uni}}
\ar[ur]
\ar[dr, "\mathrm{curv}"]
& & 
k^{\mathrm{rel}}\bC
\\
& 
\mathrm{Def}(\widehat{ku}{}_{\mathrm{uni}})
\ar[ur, "\mathrm{def}"]
\ar[rr, "d"']
& & 
\mathrm{Cyc}(\widehat{ku}{}_{\mathrm{uni}}) 
\ar[ur]
&
\end{tikzcd}
\end{equation*}



\begin{definition}\cite[Remark 6.6]{bunkenikolausvoelkl2016diffcoh}
Let $T$ be a differential cohomology theory. An \emph{additive $T$-valued differential characteristic class for complex vector bundles} is a map
\[
\Vect^\nabla \longrightarrow T
\]
of $\CMon(\lS)$-valued sheaves on $\Mfd$.
\end{definition}

\begin{remark}
\label{RemarkUniqueFactorization}
Any additive $T$-valued differential characteristic class for complex vector bundles factors uniquely through $\widehat{ku}_{\mathrm{uni}}$, as
\begin{align*}
\mathrm{Map}(\Vect^\nabla,  T)
&\cong 
\mathrm{Map}(\mathrm{grp}(\Vect^\nabla),  T)
\\
&\cong 
\mathrm{Map}(\mathrm{shv}(\mathrm{grp}(\Vect^\nabla)),  T)
\\
&\cong 
\mathrm{Map}(\widehat{ku}{}_{\mathrm{uni}},  T),
\end{align*}
using that $T$ is group complete and a sheaf.
\end{remark}


\begin{example}
If $T$ is the $\bR$-invariant differential cohomology theory defined by the spectrum $\Sigma^n HA$ for some abelian group $A$, then by \cref{RemarkUniqueFactorization}, we have
\begin{align*}
\mathrm{Map}(\Vect^\nabla,  \Sigma^n \underline{HA})
&\cong 
\mathrm{Map}(\widehat{ku}{}_{\mathrm{uni}},  \Sigma^n \underline{HA})
\\
&\cong 
\mathrm{Map}(L_{hi}(\widehat{ku}{}_{\mathrm{uni}}),  \Sigma^n \underline{HA})
\\
&\cong 
\mathrm{Map}(L_{hi}(\widehat{ku}{}_{\mathrm{uni}})(*),  \Sigma^n HA)
\\
&\cong 
\mathrm{Map}(ku,  \Sigma^n HA).
\end{align*}
Hence a homotopy class of such maps assigns, for each manifold $M$, a map
\[
\varphi : K^0(M) \longrightarrow H^n(M, A)
\]
with the property that $\varphi(V \oplus W) = \varphi(V) + \varphi(W)$; in other words, an $A$-valued additive characteristic class of degree $n$ in the classical sense.
\end{example}





We construct a transformation of differential cohomology theories
\[
\widehat{ku}{}_{\mathrm{uni}} \longrightarrow \widehat{ku}_{\mathrm{HS}}.
\]
To this end, we observe that we have a commutative diagram
\[
\begin{tikzcd}
	\widehat{ku}_{\mathrm{uni}} \ar[r] \ar[d] & H((\Omega \otimes_\bR \bC[x])^{\geq 0})
	\ar[d]
	\\
	\underline{ku} \ar[r, "\ch"] 
	& \underline{H(\bC[x])}
\end{tikzcd}
\]
where the top horizontal map is given as follows. First, we have a map of monoids
\[
\widehat{\ch} : \Vect^\nabla(M) \longrightarrow \Omega^0(M, \bC[x]), \qquad (V, \nabla) \longmapsto \ch(\nabla)
\]
sending a vector bundle with connection to its Chern character form. This yields a map of functors $\Mfd^{\mathrm{op}} \to \CMon(\lS)$, when the right hand side is considered as a discrete abelian category. Generally, under the identification of $\CMon(\lS) \simeq \Sp_{\geq 0}$, an abelian group $A$, viewed as a discrete space is sent to $HA$. Applying this to the map above, we get a map
\[
\mathrm{grp}(\Vect^\nabla(M)) \longrightarrow H\bigl(\Omega^{0}(M, \bC[x])\bigr) \hookrightarrow H\bigl(\Omega^{\geq 0}(M, \bC[x])\bigr).
\]
Varying the manifold $M$, the right hand side is a sheaf, hence we get the desired map of sheaves $\widehat{ku}_{\mathrm{uni}} \to H((\Omega \otimes_\bR \bC[x])^{\geq 0})$. The left vertical map is the canonical map that forgets the connection data. Commutativity of the diagram on homotopy groups follows from Chern-Weil theory. To get a fully coherent filler of the diagram, one has to work a little bit, see \cite[\S6.1]{bunkenikolausvoelkl2016diffcoh}.

\begin{remark}
There is a non-trivial consequence of this. The transformation constructed above yields a map of spectra
\[
K\bC \simeq R_{hi}(\widehat{ku}_{\mathrm{uni}})(*) \longrightarrow R_{hi}(\widehat{ku}_{\mathrm{HS}})(*) = \Sigma^{-1}\widetilde{ku}_{\bC^\times}.
\]
In particular, we have $\pi_{2n+1}(R_{hi}(\widehat{ku}_{\mathrm{HS}})(*)) \cong \bC^\times$ for $n \geq 0$, hence we obtain maps
\[
K_{2n+1}(\bC) \longrightarrow \bC^\times.
\]
These are so-called \emph{regulator} maps. They can be verified to be non-trivial.
\end{remark}

\section{\texorpdfstring{$\bC^\times$}{C}-valued \texorpdfstring{$K$}{K}-theory}
Let $KU_{\bC^\times}$ be the spectrum which is the cofiber of the non-connective Chern character $\ch : KU \longrightarrow H\bC[x, x^{-1}]$. The corresponding $K$-theory groups $K^n(M, \bC^\times)$ fit into a long exact sequence
\[
\begin{tikzcd}
K^{\mathrm{ev}}(M) \ar[r]
& 
H^{\mathrm{ev}}(M, \bC[x, x^{-1}])
\ar[r]
&
K^{\mathrm{ev}}(M, \bC^\times)
\ar[d]
\\
K^{\mathrm{odd}}(M, \bC^\times) \ar[u] & H^{\mathrm{odd}}(M, \bC[x, x^{-1}]) \ar[l]& K^{\mathrm{odd}}(M) 
\ar[l]
\end{tikzcd}
\]
All these groups are 2-periodic. We aim to describe this long exact sequence and all its maps (although this is mainly guesswork).

\subsection{Chern-Simons forms.}

Let $\nabla$ and $\nabla'$ be two connections on a manifold. Then there is an odd differential form $\mathrm{CS}(\nabla, \nabla')$ with the property that
\begin{equation}
\label{DefiningPropertyChernCharacter}
d \mathrm{CS}(\nabla, \nabla') = \ch(\nabla') - \ch(\nabla),
\end{equation}
where $\ch(\nabla')$ and $\ch(\nabla)$ are the Chern character forms of the connections. It may be defined as follows: On the manifold $M \times [0, 1]$, take the connection $\tilde{\nabla} = (1-t)\mathrm{pr}_M^*\nabla + t\mathrm{pr}_M^*\nabla'$ and let
\[
\mathrm{CS}(\nabla, \nabla') := \int_{[0, 1]}\ch(\tilde{\nabla}) \in \Omega^{\mathrm{odd}}(M, \bC) 
\]
be defined by integration over the fiber. More generally, we may replace the convex interpolation of $\nabla$ and $\nabla'$ by a general connection on $M \times [0, 1]$ that restricts to $\nabla$, respectively $\nabla'$ at $\partial [0, 1]$. In this case, the Chern-Simons form $\mathrm{CS}(\tilde{\nabla})$ is defined in a similar fashion and it also satisfies \cref{DefiningPropertyChernCharacter}. Notice that if two connections $\tilde{\nabla}$ and $\tilde{\nabla}'$ on $M \times [0, 1]$ are gauge equivalent, then $\ch(\tilde{\nabla}) = \ch(\tilde{\nabla}')$, hence $\mathrm{CS}(\tilde{\nabla}) = \mathrm{CS}(\tilde{\nabla})$. Moreover, if the connection $\tilde{\nabla}$ is constant in $t$, then $\mathrm{CS}(\tilde{\nabla}) = 0$.

%More generally, an arbitrary connection $\tilde{\nabla}$ on $M \times \Delta_2$ induces three paths of connections by restricting to $\partial \Delta_2$, which we denote by $\tilde{\nabla}^{01}$, $\tilde{\nabla}^{12}$ and $\tilde{\nabla}^{02}$.
%We may then define
%\[
%\CS_2(\tilde{\nabla}) := \int_{\Delta_2} \ch(\tilde{\nabla}) \in \Omega^{\mathrm{ev}}(M, \bC), 
%\]
%satisfying
%\[
%d\CS_2(\tilde{\nabla}) = \mathrm{CS}(\tilde{\nabla}^{01}) + \mathrm{CS}(\tilde{\nabla}^{12}) - \mathrm{CS}(\tilde{\nabla}^{02}).
%\]
%In particular, if $\tilde{\nabla}^{12}$ is constant so that $\tilde{\nabla}^{01}$ and $\tilde{\nabla}^{02}$ restrict to the same connections at the end point, then $\mathrm{CS}(\tilde{\nabla}^{12}) = 0$.
%Hence the above formula shows that if we interpolate between a pair of connections using two different paths, then the corresponding Chern-Simons forms differ by an exact term.

%Given three connections $\nabla, \nabla'$ and $\nabla''$ on $M$, convex interpolation defines a connection $\tilde{\nabla}$ on $M$ and we may define
%\[
%\CS_2(\nabla, \nabla', \nabla'') := \int_{\Delta_2} \ch(\tilde{\nabla}) \in \Omega^{\mathrm{ev}}(M).
%\]
%Then
%\[
%d\CS_2(\nabla, \nabla', \nabla'') = \mathrm{CS}(\nabla, \nabla') + \mathrm{CS}(\nabla', \nabla'') - \mathrm{CS}(\nabla, \nabla'').
%\]

\medskip



If $g : M \to U(\infty)$ is a smooth map (in the sense that it factors through some $U(m)$ and is smooth there), it has an odd Chern character, which is the differential form
\[
\ch(g) := -\mathrm{CS}(d, d+g^{-1}dg) \in \Omega^{\mathrm{odd}}_{\mathrm{cl}}(M, \bC).
\]
Here we view $d$ and $d+ g^{-1}dg$ as connections on the trivial vector bundle $\underline{\bC}^m$. Notice here that the connection $d+ g^{-1}dg$ is gauge equivalent to $d$, hence has the same Chern character as $d$, namely zero. From \cref{DefiningPropertyChernCharacter}, we therefore see that $\ch(g)$ is a closed form. Including the minus sign in the definition will be important later.

If $g, h : M \to U(\infty)$ are homotopic via a homotopy $(g_t)_{t \in [0, 1]}$ (which we require to factor through some finite $U(m)$), this yields a map $\tilde{g} : M \times [0, 1] \to U(m)$. The \emph{odd Chern/Simons form} is then defined as
\[
\mathrm{CS}((g_t)_{t \in [0, 1]}) = \int_{[0, 1]} \ch(\tilde{g}) \in \Omega^{\mathrm{ev}}(M),
\]
which satisfies
\[
d \mathrm{CS}(g, h) = \ch(h) - \ch(g).
\]
One may check that the Chern-Simons forms for two different homotopies differ by an exact form.
%
% then we may consider the connection $\tilde{\nabla}$ on $M \times \Delta_2$ which on $\partial_{01}\Delta_2$ is given by the convex interpolation between $d$ and $d + g^{-1}dg$, on $\partial_{02}\Delta_2$ by the convex interpolation between $d$ and $d + h^{-1} dh$ and on $\partial_{12} \Delta_2$ the path of connections $d + g_t^{-1} dg_t$.
%This last connection is gauge equivalent constant one, so $\mathrm{CS}(\tilde{\nabla}^{12}) = 0$.
%We therefore obtain that
%\[
%\mathrm{CS}(g, h) := \CS_2(\tilde{\nabla}) \in \Omega^{\mathrm{ev}}(M, \bC)
%\]
%satisfies
%\[
%d \mathrm{CS}(g, h) = \ch(h) - \ch(g).
%\]
%We call it the \emph{odd Chern-Simons form}.



\subsection{The \texorpdfstring{$\bC^\times$}{C}-valued \texorpdfstring{$K$}{K}-theory groups.}
We first describe the group $\tilde{K}^{\mathrm{odd}}(M, \bC^\times)$ for $n$ odd. Consider equivalence classes $(V, \nabla, \eta)$, where $V$ is a complex vector bundle, $\nabla$ is a connection on $V$ and $\eta$ is an odd differential form such that $d\eta = \smash{\widetilde{\ch}}(V) = \ch(V) - \mathrm{rk}(V)$. In particular, the reduced Chern character of $V$ vanishes. Two such triples $(V, \nabla, \eta)$ and $(V', \nabla', \eta')$ are equivalent if there exists a vector bundle isomorphism $\phi : V \to V'$ (not necessarily connection preserving) such that
\[
\mathrm{CS}(\nabla, \phi^*\nabla') = \eta' - \eta \mod \text{exact forms}.
\]
The set of equivalence classes of these triples form a monoid under direct sum and addition of differential forms, and we denote by $\tilde{K}^{\mathrm{odd}}(M, \bC^\times)$ the Grothendieck group of this monoid. Then
\[
K^{\mathrm{odd}}(M, \bC^\times) \subset \tilde{K}^{\mathrm{odd}}(M, \bC^\times)
\] 
is the subgroup consisting of those elements with virtual dimension zero.

\begin{remark}
Of course, a vector bundle with vanishing Chern character is not necessarily flat. For example, the Chern character already vanishes when all Chern classes are torsion (because the classifying map lifts along $BU_\delta \to BU$). On the other hand, if the manifold is simply connected, any flat vector bundle must be trivial. 
\end{remark}

The group $K^{\mathrm{ev}}(M, \bC^\times)$ consists of pairs $(g, \eta)$, where $g : M \to U(\infty)$ is a smooth, and $\eta$ is an even differential form satisfying $d\eta = \ch(g)$, the odd Chern character of $g$. Two pairs $(g, \eta)$ and $(g', \eta')$ are equivalent if $g$ and $g'$ are homotopic and 
\[
\mathrm{CS}(g, g') = \eta' - \eta \mod \text{exact forms}.
\]
The set of such triples is a monoid under direct sum and $K^{\mathrm{ev}}(M, \bC^\times)$ is the Grothendieck group of this monoid.

The maps $K^{\mathrm{odd}}(M, \bC^\times) \to K^{\mathrm{ev}}(M)$ forgets the connection and differential form data and the map $K^{\mathrm{ev}}(M, \bC^\times) \to K^{\mathrm{odd}}(M)$ forgets the differential form data. The map $H^{\mathrm{odd}}(M, \bC) \to K^{\mathrm{odd}}(M, \bC^\times)$ takes a class $[\rho]$ and sends it to the triple $(0, 0, \rho)$. This is well-defined by the equivalence relation in $K^{\mathrm{ev}}(M, \bC^\times)$ (we need to specify $\eta$ only up to exact forms). Similarly, $H^{\mathrm{ev}}(M, \bC) \to K^{\mathrm{ev}}(M, \bC^\times)$ sends $[\rho]$ to $(1, \rho)$.

\subsection{We have a complex}
Then $[g] \in K^{\mathrm{odd}}(M)$ is sent to $[\ch(g)]$ in $H^{\mathrm{odd}}(M, \bC)$. Its image in $K^{\mathrm{odd}}(M, \bC^\times)$ is $(0, 0, \ch(g))$. Stabilizing, this is equivalent to $(\underline{\bC}^m, d, \ch(g)) - (\underline{\bC}^m, d, 0)$. But we have
\[
(\underline{\bC}^m, d, \ch(g)) \sim (\underline{\bC}^m, d, 0)
\]
Indeed, the vector bundle automorphism $g: \underline{\bC}^m \to \underline{\bC}^m$ yields
\[
- \ch(g) = \mathrm{CS}(d, d + g^{-1}dg) = \mathrm{CS}(d, g^*d).
\]

$[V] \in K^{\mathrm{ev}}(M)$ is sent to $\ch([V]) \in H^{\mathrm{ev}}(M, \bC)$.
This is sent to $(1, \ch(V))\in K^{\mathrm{odd}}(M, \bC^\times)$, where $\ch(V)$ is an arbitrary differential form representative of the Chern class of $V$.{\color{red}[Why is this zero?]}

\subsection{Exactness.}
We show exactness at $K^{\mathrm{odd}}(M, \bC^\times)$. Consider the kernel of the map $K^{\mathrm{odd}}(M, \bC^\times) \to K^{\mathrm{ev}}(M)$ in the case that $n$ is odd. A formal difference $[V, \nabla, \eta] - [V', \nabla', \eta']$ is in the kernel if and only if $[V] - [V'] = 0$ in $K^{\mathrm{ev}}(M)$. By changing representatives, we may assume that $V'$ is trivial, which implies that $V$ is stably trivial. By changing representatives again, we may assume that $V$ is trivial as well. Hence elements in the kernel are represented by triples of the form $(\underline{\bC}^m, d+\alpha, \eta)$, where $\alpha \in \Omega^1(M, M_m(\bC))$ and $\eta$ is an odd differential form (minus a trivial bundle with trivial connection data that sets the virtual rank to zero). Now, 
\[
d\bigl(\mathrm{CS}(d, d+\alpha) - \eta\bigr) =  d\mathrm{CS}(d, d+\alpha) - d\eta = \ch(d+\alpha) - \ch(d+\alpha) = 0
\]
and $\mathrm{CS}(d, d+\alpha) - \eta$ is an odd closed form, hence in the image of the map from de Rham cohomology.

Now consider the map $K^{\mathrm{ev}}(M, \bC^\times) \to K^{\mathrm{odd}}(M)$.
Suppose that $(g, \eta)$ is in the kernel of this map. Then $g$ is null homotopic, hence after choosing representatives, we may assume that $g$ is constant, so $\ch(g) = 0$. That means that $\eta$ is closed, so it is again in the map from de Rham cohomology.

\section{The Hopkins-Singer differential \texorpdfstring{$K$}{K}-theory groups}

The following should now be a valid description of the group $\widehat{ku}{}^0_{\mathrm{HS}}(M)$: Elements are represented by triples $(V, \nabla, \eta)$, where $\eta \in \Omega^{\mathrm{odd}}(M, \bC)$ (with no further condition on $\eta$). The geometric equivalence relation is $(V, \nabla, \eta) \sim (V', \nabla', \eta')$ iff there exists a vector bundle isomorphism $\phi : V \to V'$ such that $ \mathrm{CS}(\nabla, \phi^*\nabla') = \eta - \eta'$. Again, taking the Grothendieck completion of the resulting monoid yields an abelian group.

In this description, the arrow from $\tilde{k}^{-1}(M, \bC^\times)$ is the obvious inclusion, the map to $K^0(M)$ is the projection onto the $K$-theory class and
\[
\mathrm{def}([\eta]) = (0, 0, -\eta), \qquad \mathrm{curv}(V, \nabla, \eta) = \ch(\nabla) - d\eta.
\]

\begin{remark}
Hopkins--Singer \cite{hopkinssinger2005diffcoh} give the following cocycle description of the group $\widehat{ku}_{\mathrm{HS}}(M)$. Let $BU$ the classifying space of the infinite unitary group $U = \colim_n U(n)$. Then $\lF = \bZ \times BU$ is a classifying space for the functor $K^0$. We denote by $C^*(\lF, \bR[x])$ the singular cochains on $\lF$ with values in $\bR[x]$, where $x$ has degree $-2$. Let $\ch_{\mathrm{univ}} \in C^0(\lF, \bR[x])$ be a singular cochain representative of the universal Chern class (for some reason denoted by $\iota$ by Hopkins--Singer). The elements of the group $\smash{\widehat{ku}}{}^0_{\mathrm{HS}}(M)$ are then represented by triples $(h, c, \omega)$, where   $c: M \to \lF$ is a continuous map, $\omega \in \Omega^0(M, \bR[x])$ is a de Rham form and $h \in C^{-1}(M, \bR[x])$ is such that
\[
dh = c^*\ch_{\mathrm{univ}} - \omega.
\]
Here $\omega$ is viewed as a singular cochain using the de Rham map. Two elements $(c, h, \omega)$ and $(c', h', \omega')$ represent the same element if there exists a triple $(\tilde{h}, \tilde{c}, \tilde{\omega})$ over $M \times [0, 1]$ such that $\tilde{\omega} = \mathrm{pr}_M^*\omega$ and that restricts to the given cocycles at zero and one.

This description should correspond to ours by choosing a universal Hermitian connection on $BU$ and taking $\ch_{\mathrm{uni}}$ to be its Chern character form, see \cite{narasimhanramanan1961universalconnections}.
\end{remark}

\chapter{What (is) a Twist(ed Cohomology Theory?)! \texorpdfstring{\\}{ } \normalsize \emph{Talk by Hannes Berkenhagen}}

The aim of this talk is to review twisted cohomology theory with the aim of later discussing twisted differential cohomology theories \cite{bunkegepner2021differential}. For these talks the main source is \cite{abghr2014infty,abg2018thom}.

\section{Twisted Cohomology}
Let $R$ be a ring spectrum, meaning a monoid object in the $\infty$-category of spectra $\Sp$. From this we get a presentable stable $\infty$-category $\Mod_R$ of left $R$-module spectra. Objects therein are morphisms of the form $R \wedge M \to M$ satisfying the usual associativity and unit conditions up to coherent homotopies.
\begin{remark}
	Every stable category $\lC$ is enriched over spectra. Given two objects $C,D$, we will denote by $\Hom_\lC(C,D)$ the hom-spectrum and by $\hom_\lC(C,D)=\Omega^\infty\Hom_\lC(C,D)$ the underlying hom-space.  
\end{remark}
Note we have an adjunction 
\[
\begin{tikzcd}[column sep=2cm]
\Sp \arrow[r, shift left, "R \wedge -"] & \Mod_R \arrow[l, shift left, "{\Hom_{\Mod_R}(R,-)}"]
\end{tikzcd}, 
\]
where the right adjoint is equivalent to the forgetful functor.
\begin{definition}
	Let $R$ be a ring spectrum. An \emph{$R$-line} is an $R$-module $L$ such that $L \simeq R$.
\end{definition}
\begin{definition}
	Let $\Line_R$ be the full sub-$\infty$-groupoid of $\Mod_R$ spanned by the $R$-lines.
\end{definition}
By construction, $\Line_R$ is equivalent to the category with a single object $R$ and hom-space $GL_1R$, the $\infty$-group of $R$-linear automorphisms of $R$. Notice that $GL_1(R)\subseteq\hom_{\Mod_R}(R,R)\simeq\hom_{\Sp}(\bS,R)=\Omega^\infty R$. 
\begin{lemma}\label{lem:units}
	$GL_1(R)$ fits into the pullback square
	\begin{center}
		\begin{tikzcd}
			GL_1(R)\arrow[r]\arrow[d] & \Omega^\infty R\arrow[d]\\
			\pi_0(R)^\times \arrow[r] & \pi_0(R)
		\end{tikzcd}
	\end{center}In particular, the inclusion $GL_1(R)\to\Omega^\infty R$ induces an isomorphism on $n$-homotopy groups, for all $n\geq1$. 
\end{lemma}
\begin{proof}
	$\pi_0(R)\simeq\hom_{\Ho\Mod_R}(R,R)$, where $\Ho\Mod_R$ is the homotopy category of $R$-modules, the right-vertical arrow corresponds to $\hom_{\Mod_R}(R,R)\to\hom_{\Ho\Mod_R}(R,R)$ mapping a morphism to its homotopy class, and $\pi_0(R)^\times\subseteq\hom_{\Ho\Mod_R}(R,R)$ is the set of isomorphisms. Finally, a morphism in a $\infty$-category $\lC$ is an equivalence if and only if its homotopy class is an isomorphism in $\Ho\lC$. 
\end{proof}
\begin{definition}
	Let $X$ be a space. Denote by $\Mod_R(X)$ the $\infty$-category of $R$-module spectra parametrized over $X$, i.e. the functor category $\Fun(X^{\op},\Mod_R)$.
\end{definition}

\begin{definition}
	Let $X$ be a space. Denote by $\Line_R(X)$ the $\infty$-groupoid of $R$-line spectra parametrized over $X$, i.e. the functor category $\Fun(X^{\op},\Line_R)$. 
\end{definition}
Since $\Line_R\simeq BGL_1(R)$, functors $X^{\op}\to\Line_R$ classify $GL_1(R)$-principal bundles over $X$. In this sense, elements of $\Line_R(X)$ are $\infty$-local systems. 
\begin{example}
	Let $R_X\colon X^{\op} \to*\to \Line_R$ be the constant functor with value $R$. 
\end{example}
We now proceed to the Thom construction. We introduce it to full generality. 
\begin{definition}[{Thom spectrum}] \label{def:thom}
	The \emph{Thom $R$-module spectrum} is the functor
	\[
	R \colon \lS/\Mod_R \to \Mod_R,
	\]
	which sends $f\colon X^{\op} \to \Mod_R$ to the $R$-module spectrum $R^f$ given by colimit of $f$. 
\end{definition}
\begin{remark}The colimit of a functor $F:\lC\to\lX$ is the left Kan extension of $F$ along the terminal functor $p:\lC\to*$. If $\lX$ is cocomplete, a functor $\pi:\lC\to\lD$ induces a pullback functor $\pi^*:\Fun(\lD,\lX)\to\Fun(\lC,\lX)$, which is cocontinuous, therefore using left Kan extension along $\pi$ we can construct a left adjoint $\pi_!$.  
\end{remark}
Let us have a look at some examples of Thom spectra. 
\begin{example}[Trivial twist] A functor $f:\lX\to\lY$ induces a post-composition functor $f_*:\lS/\lX\to\lS/\lY$ that is left adjoint to $f^*$ sending $F:\lC\to\lY$ to the canonical projection $\lX\times_{\lY}\lC\to\lX$. If $\lX=*$, we can identify $\lS/\lX$ with $\lS$, and $f$ with an object $E\in\lY$. If $\lY$ is cocomplete, there is a colimit functor $\pi_!:\lS/\lY\to\lY$ and the composition \begin{equation}
	\begin{tikzcd}
		\lS \arrow[r,"f_*"] & \lS/\lY \arrow[r,"\pi_!"] & \lY
	\end{tikzcd}
\end{equation}preserves colimits, therefore it is identified by the image of the one-point space, which is $E$. If $\lY=\Mod_R$, then $f$ corresponds to a $R$-module $E$. The composition \begin{equation}
	\begin{tikzcd}
		\lS \arrow[r,"\Sigma^\infty_+"] & \Sp \arrow[r,"E\wedge-"] & \Mod_R
	\end{tikzcd}
\end{equation}also preserves colimits and maps the one-point space to $E$, therefore the Thom spectrum of $E_X$ is $E\wedge\Sigma^\infty_+X$, for every space $X$. 
\end{example}
% In Diff_Coh_2, include some background on homotopy colimits vs oo-colimits. In particular, that a functor BG->S is modelled by a G-simplicial set X, and its colimit by the homotopy quotient X/G

%--------------TO BE CONTINUED-----------------


\begin{definition}[{Twisted cohomology}] \label{def:twistedcohomology}
	Let $R$ be a ring spectrum, $X$ be a space, $p\colon X \to *$ the terminal functor, and $f\colon X^{\op} \to \Line_R$ a $R$-line bundle over $X$.  The \emph{$f$-twisted $R$-cohomology of $X$} is defined as the mapping spectrum 
	\[ R_f(X)	\coloneqq \Hom_{\Mod_R}(Mf, R) \simeq \Hom_{\Mod_R(X)}(f,p^*R) \simeq \Hom_{\Mod_R(X)}(f,R_X) .\]
	Similarly, the \emph{$f$-twisted $R$-homology of $X$} is defined as
	\[ R^f(X) \coloneqq \Hom_{\Mod_R}(R,Mf) \simeq Mf. \]
	The $f$-twisted $R$-cohomology groups of $X$ are defined as the homotopy groups of $R_f(X)$, i.e.
	\[ R^n_f(X) \coloneqq \pi_0(\Hom_{\Mod_R}(Mf, \Sigma^n R)) \cong \pi_{-n}(\Hom_{\Mod_R(X)}(f,R_X)). \]Similarly, the $f$-twisted $R$-homology groups of $X$ are defined as the homotopy groups of $R^f(X)$, i.e. 
	\[ R_n^f(X) \coloneqq \pi_0(\Hom_{\Mod_R}(\Sigma^n R, Mf)) \cong \pi_n(Mf). \]
\end{definition}

\begin{example}[Trivial twist] If $f:X\to \Line_R$ factors through $*$, then $f$ factors as $X^{\op}\to*\to\lS$, the constant factor with value $*$, and $R\wedge\Sigma^\infty_+(-):\lS\to\Mod_R$. The latter functor commutes with colimits, being a left adjoint, while the colimit of the latter is $X$ itself, then $Mf\simeq R\wedge\Sigma^\infty_+X$. In particular, $f$-twisted $R$-cohomology and $R$-homology of $X$ reduce to ordinary (untwisted) $R$-cohomology and $R$-homology of $X$. 
\end{example}


\section{Examples of Twisted Cohomology}
We now proceed to analyze several examples of twisted cohomology	theories. This requires some preliminary lemmas.

\begin{lemma}
	Consider a space $X$ and the $\infty$-categorical Yoneda's embedding $y:X\to\Fun(X^{\op},\lS)$. The colimit of $y$ is the terminal pre-sheaf on $X$, i.e. the pre-sheaf with constant value the one-point space. 
\end{lemma}

\begin{proof}
	Let $S$ be a pre-sheaf on $X$, consider then the slice category $X_{/S}$ of pairs $(x,\phi)$, where $x$ is an object of $X$ and $\phi:y(x)\to S$. The density theorem for $\infty$-categories states that $S$ is equivalent to the colimit of $X_{/S}\to X\xrightarrow{y}\Fun(X^{\op},\lS)$, the first map being the canonical projection. Take $S=*$, then $X_{/*}\to X$ is an equivalence, hence the claim. 
\end{proof}
Let $G$ be a topological group and $BG$ the $\infty$-groupoid with a single object $*$ and hom-space $G$. The category $\lS_G\coloneqq\Fun(BG,\lS)$ is equivalent to the category of $G$-spaces.
\begin{lemma}
	Consider $X=BG$, a $G$-space $f:X\to\lS$ and its left Kan extension $f_!:\lS_G\to\lS$, then $f_!\simeq(-\times E)/G$, where $E=f(*)$. 
\end{lemma}
\begin{proof}
	Evaluate at $*$, then $f_!(y(*))=E$, by definition, and $y(*)\simeq G$, as $G$-spaces, hence $(y(*)\times E)/G\simeq(G\times E)/G\simeq E$. Since $f_!$ and $(-\times E)/G$ agree on representables and are colimit-preserving, they are equivalent. 
\end{proof}
\begin{example}\label{ex:bundle}Take the space $BO(n)$ and $f_n:BO(n)\to\lS_*$ the $n$-sphere $S^n$ with $O(n)$-action coming from the one-point compactification of the regular action on $\bR^n$. Let $\alpha_n=\Sigma^{\infty-n}f_n:BO(n)\to\Sp$, then $\alpha_n(*)=\Sigma^{\infty-n}f_n(*)=\Sigma^{\infty-n}S^n\simeq\bS$, so $\alpha_n$ factors through $\Line_\bS$. Let $X=BO(n)^{\op}$ and $p:X\to*$ the terminal functor, then \[M\alpha_n=p_!\Sigma^{\infty-n}f_n\simeq\Sigma^{\infty-n}p_!(f_n)_!y\simeq\Sigma^{\infty-n}(f_n)_!\underbrace{p_!(y)}_{\simeq*}\simeq\Sigma^{\infty-n}(*\times S^n)/O(n)\]
Let $P=EO(n)$ be the universal $O(n)$-bundle and $M$ a $O(n)$-space, then $*\times_{O(n)} M$ is modelled by the \emph{strict} quotient $(P\times M)/O(n)$, then \[S^n/O(n)=\mathrm{cofib}(\bR^n_0\subseteq\bR^n)/O(n)\simeq\mathrm{cofib}(\underbrace{*\times_{O(n)}\bR^n_0}_{\simeq E^n_0}\subseteq\underbrace{ *\times_{O(n)}\bR^n}_{\simeq E^n})=\mathrm{Th}(E^n) \]
where $E^n=P\times_{O(n)}\bR^n\to BO(n)$ is the universal $n$-dimensional vector bundle, hence $M\alpha_n\simeq\Sigma^{\infty-n}\mathrm{Th}(E^n)$. 
\end{example}
The functor $BO(n)\to\Line_\bS$ induces a $\infty$-group homomorphism $j_n:O(n)\to GL_1(\bS)$, mapping $\phi$ to $\Sigma^{\infty-n}\mathrm{Th}(\phi)$. Consider the suspension morphism $s_n=\bR\oplus-:O(n)\to O(1+n)$, then $$j_n(\bR\oplus\phi)=\Sigma^{\infty-n-1}\mathrm{Th}(\bR\oplus\phi)\simeq\Sigma^{\infty-n-1}\underbrace{\mathrm{Th}(\bR)}_{\simeq S^1}\wedge\mathrm{Th}(\phi)\simeq\Sigma^{\infty-n}\mathrm{Th}(\phi)=j_n$$Recall that the colimit over the suspension morphisms $s_n$ is the stable orthogonal group $O$.
\begin{definition}
	Denote by $j$ the induced group homomorphism $O\to GL_1(\bS)$, called the \emph{$J$-homomorphism}.  
\end{definition} 
\begin{example}
	Let $X=O^{\op}$ and take $Bj:BO\to\Line_\bS$, then $Mj$ is denoted $MO$ and called the \emph{real bordism spectrum}. 
\end{example}
Denote by $M$ the extended Thom spectrum functor $\Grp^{\op}_\infty/\Mod_R\to\Mod_R$, this is a left adjoint to the functor sending a $R$-module to the corresponding functor $*\to\Mod_R$. In particular, $M$ preserves colimits and  $Bj\simeq\colim_nBj_n$, therefore we have the following:
\begin{theorem}
	$MO\simeq\colim_nMO(n)=\colim_n\Sigma^{\infty-n}\mathrm{Th}(E^n)$. 
\end{theorem}
\begin{example}\label{ex:bordismspectra}
	A group homomorphism $\xi:G\to O$ induces a functor $f:BG\to\Line_\bS$. The Thom spectrum $Mf$ is denoted $MG$ or $M\xi$, and called \emph{$G$-bordism spectrum}. For $G=U,SO,Spin$, and $String$, we obtain the \emph{complex, oriented, spin}, and \emph{string bordism spectra}. 
\end{example}
\begin{remark}
	In \cref{ex:bordismspectra} we might take $G=\{*\}$, the one-point group, then $MG\simeq\bS$, which, if it didn't have a name, might be called \emph{framed bordism spectrum}, following the naming convention in \cref{ex:bordismspectra} and in line with the theorem that $\pi_*(\bS)\simeq\Omega^\mathrm{fr}_*$, the bordism ring of framed (trivialized tangent bundle) smooth manifolds.  
\end{remark}
Let $R$ be a ring in sets, then $R$ is a $A_\infty$-ring spectrum (actually, $E_\infty$), $\Omega^\infty R$ is equivalent to $R$ with discrete topology ($\pi_0(\Omega^\infty R)\simeq R$, as sets, and every other homotopy group vanish). In particular, $GL_1(R)$ is simply $R^\times$ with discrete topology. Consider then the fiber sequence $SO\to O\to \bZ^\times\simeq GL_1(\bZ)$.
\begin{example}
	$X=BO^{\op}$ and $\alpha=w_1:BO\to\Line_\bZ$, the 1st Stiefel-Whitney class (delooping of the determinant $O\to GL_1(\bZ)$), then $Mw_1$ is a $\bZ$-module spectra. Let $i:SO\subseteq O$, then $w_1i$ factors through the point, so $M(w_1i)\simeq\bZ\wedge\Sigma^\infty_+SO$. 
\end{example}
Given $f:X^{\op}\to\Line_R$ and a sequence $F\xrightarrow{i} Y\xrightarrow{\pi} X$, there is an induced sequence of Thom $R$-module spectra $MF\to MY\to MX$. If $\pi i$ factors through the point, $MF\simeq R\wedge\Sigma^\infty_+F$. 
\begin{lemma}\label{lem:adj}
	Let $R$ be a ring spectrum and $X$ a connected monoidal $\infty$-groupoid, then \[\hom_{\Mon(\Sp)}(\Sigma^{\infty}_+X,R)\simeq\hom_{\Mon(\lS)}(X,GL_1(R))\] 
\end{lemma}
\begin{proof}
	Since $X$ is connected, the space of homomorphisms $X\to GL_1(R)$ is equivalent to the space of homomorphisms $X\to\Omega^\infty R$, then use that $(\Sigma^\infty,\Omega^\infty)$ is a monoidal adjunction (The monoidal structure on spectra is such that $\Sigma^\infty$ is strong monoidal). 
\end{proof}
\begin{remark}
	Notice that we can weaken the result. Namely, if $X$ is 1-connected (pointed and connected), then the space of functors (of $\infty$-groupoids) $X\to GL_1(R)$ is equivalent to the space of functors $X\to\Omega^\infty R$ such that $*\to X\to\Omega^\infty R$ is an equivalence. This last space is equivalent, via the $(\Sigma^\infty_+,\Omega^\infty)$ adjunction, to the space of morphisms of spectra $\Sigma^\infty_+X\to R$, such that $\bS\to\Sigma^\infty_+X\to R$ represents a unit in $\pi_0(R)$.
\end{remark}
\begin{remark}\label{rmk:strengthen}
	Notice that we can also strengthen the result. Namely, if $X$ is a connected, commutative monoid object and $R$ is a commutative ring spectrum, then $\Omega^\infty R$ and $GL_1(R)$ are also commutative monoid objects. Using the same argument, together with the fact that $(\Sigma^\infty,\Omega^\infty)$ is actually a \emph{symmetric} monoidal adjunction, we conclude that \[\hom_{\CMon(\Sp)}(\Sigma^\infty_+X,R)\simeq\hom_{\CMon(\lS)}(X,GL_1(R))\] 
\end{remark}
\begin{remark}\label{rmk:moncat}
	Let $\lD$ be a monoidal $\infty$-category. Consider $\cat_\infty/\lD$, the $\infty$-category of functors into $\lD$, with monoidal structure given by \begin{center}
		\begin{tikzcd}
			(F:\lA\to\lD,G:\lB\to\lD) \arrow[r,|->] & (\lA\times\lB\xrightarrow{F\times G}\lD\times\lD\xrightarrow{\otimes}\lD)
		\end{tikzcd}
	\end{center} The monoidal unit is the functor $*\to\lD$ picking out the monoidal unit of $\lD$. If $\lD$ is symmetric monoidal, then so is $\cat_\infty/\lD$. A (commutative) monoid object in $\cat_\infty/\lD$ is given by a (symmetric) monoidal category $\lC$ and a (symmetric) monoidal functor $F:\lC\to\lD$. 
\end{remark}
In view of \cref{rmk:moncat}, let $R$ be commutative ring spectrum, then $\Mod_R$ is a symmetric monoidal $\infty$-category and $\Line_R$ is a symmetric monoidal $\infty$-groupoid. The category $\Grpd^{\op}_\infty/\Line_R$ is then symmetric monoidal and (commutative) monoid objects are given by (symmetric) monoidal $\infty$-groupoids $X^{\op}$ a (symmetric) monoidal functors $X^{\op}\to\Line_R$. One can then check that $M$ is a symmetric monoidal functor, so that (commutative) monoid objects are sent to (commutative) monoid objects in $\Mod_R$, i.e. (commutative) $R$-algebras. 
\begin{example}\label{ex:tmf}
	Let $\mathrm{Tmf}$ be the commutative ring spectrum of topological modular forms (see \cref{rmk:names}) and $\sigma:MString\to\mathrm{tmf}$ the $String$-orientation of $\mathrm{tmf}$. In the sequence \begin{center}
		\begin{tikzcd}
			BString \arrow[r] & BO \arrow[r,"Bj"] & \Line_\bS
		\end{tikzcd}
	\end{center}all functors are symmetric monoidal, so that $MString$ is a commutative $\bS$-algebra, i.e. a commutative ring. The $String$-orientation of $\mathrm{tmf}$ is also a commutative ring homomorphism. In the fiber sequence $K(\bZ,3)\to BString\to BSO$, the fiber map $i:K(\bZ,3)\to BString$ is also a symmetric monoidal, so the composition 
	\begin{center}
		\begin{tikzcd}
			\Sigma^\infty_+K(\bZ,3)\arrow[r,"Mi"] & MString \arrow[r,"\sigma"] & \mathrm{tmf}
		\end{tikzcd}
	\end{center}is a commutative ring homomorphism. Using \cref{lem:adj}, we conclude that the induced homomorphism $K(\bZ,3)\to\Omega^\infty\mathrm{tmf}$ (which is a homomorphism of commutative monoid objects, given \cref{rmk:strengthen}) factors through $GL_1(\mathrm{tmf})$, and so it induces a (symmetric monoidal) functor $K(\bZ,4)\to\Line_\mathrm{tmf}$, i.e. \emph{2-bundle gerbes twist $\mathrm{tmf}$}.
\end{example}
\begin{remark}\label{rmk:names}The spectrum of topological modular forms comes in three main flavors, namely:
\begin{enumerate}
	\item $\mathrm{TMF}$, i.e. the global sections of the spectral structure sheaf $\rO^{top}:(\mathrm{Aff}/\lM_{ell})^{\op}\to\CMon(\Sp)$ on the (\' etale site of the) moduli stack of elliptic curves.
	\item $\mathrm{Tmf}$, i.e. the global sections of the spectral structure sheaf $\bar{\rO}^{top}:(\mathrm{Aff}/\bar{\lM}_{ell})^{\op}\to\CMon(\Sp)$ on the (\' etale site of the) \emph{compactified} moduli stack of elliptic curves. The inclusion $\lM_{ell}\hookrightarrow\bar{\lM}_{ell}$ induces a commutative ring homomorphism $\mathrm{Tmf}\to\mathrm{TMF}$.
	\item $\mathrm{tmf}$, i.e. the connective cover of $\mathrm{Tmf}$. By definition, there is a commutative ring homomorphism $\mathrm{tmf}\to\mathrm{Tmf}$. 
\end{enumerate}In \cite{ando2010multiplicative}, $\mathrm{tmf}$ is used to denote our $\mathrm{TMF}$, in \cite{goerss2009topologicalmodularformsaftern}, $\mathrm{tmf}$ is used to denoted our $\mathrm{Tmf}$, and \cite{douglas2014topological} has the same notation as us. In \cref{ex:tmf}, we use $\mathrm{tmf}$ to mean the connective cover of $\mathrm{Tmf}$. 
\end{remark}
Let us go down one step in the chromatic ladder. 
\begin{example}\label{ex:ktheory}
	Recall the fiber sequences for $Spin$ and $Spin^c$: \[\bZ_2\to Spin\to SO,\qquad S^1\to Spin^c\to SO\] All the spaces involved are commutative groups. Applying the Thom spectrum functor to the delooped sequences, we get\[\Sigma^\infty_+K(\bZ_2,1)\to MSpin\to MSO,\qquad\Sigma^\infty_+K(\bZ,2)\to MSpin^c\to MSO \]
	Let $\sigma:MSpin\to KO$ and $\sigma^c:MSpin^c\to KU$ be the Atiyah-Bott-Shapiro orientation of real and complex $K$-theory (see \cite{atiyahbottshapiro1964clifford}).
	Similar to \cref{ex:tmf}, we get homomorphisms 
	\begin{center}
		\begin{tikzcd}
			\Sigma^\infty_+K(\bZ_2,1)\arrow[r] & MSpin \arrow[r,"\sigma"] & KO, & \Sigma^\infty_+K(\bZ,2)\arrow[r] & MSpin^c \arrow[r,"\sigma^c"] & KU
		\end{tikzcd}
	\end{center}
	Using \cref{lem:adj} again and delooping, we obtain functors $K(\bZ_2,2)\to\Line_{KO}$ and $K(\bZ,3)\to \Line_{KU}$, i.e. \emph{real, resp. complex, bundle gerbes twist real, resp. complex, $K$-theory}. 
\end{example}

\section{Twists via Picard Groupoids and Grading}
This section requires some some further details. Up until now we defined everything via $\Line_R$, however for many applications we need to work with $\Pic_R$ instead.
\begin{definition}
	Given a monoidal $\infty$-category $(\lC,\otimes,1)$, an object $M$ is \emph{invertible} if there is an object $N$ such that $N\otimes M\simeq M\otimes N\simeq 1$. The \emph{Picard $\infty$-groupoid} of $\lC$ is the sub-$\infty$-groupoid generated by invertible modules. 
\end{definition}
\begin{remark}\label{rmk:dualinver}
	A monoidal category $(\lC,\otimes,1)$ is \emph{closed} if, for every object $M$, the \emph{left tensoring with $M$} functor $M^\otimes:\lC\to\lC$ admits a right adjoint $F(M,-)$. If $M$ is invertible, with inverse $N$, then $M^\otimes$ is an equivalence with $N^\otimes$ as inverse. In particular, we can promote $(M^\otimes,N^\otimes)$ to an adjoint equivalence. By uniqueness of adjoint functors $F(M,-)\simeq N^\otimes$, and so \[N\simeq N\otimes1= N^\otimes(1)\simeq F(M,1)=:DM\] 
\end{remark}
\begin{definition}
	Given a closed monoidal category $(\lC,\otimes,1)$, the functor $D:=F(-,1):\lC^{\op}\to\lC$ will be called \emph{duality}. Given an object $M$, the object $DM$ is called the \emph{dual of $M$}. 
\end{definition}
\begin{definition}
If $R$ is a ring spectrum, $\Mod_R$ is closed monoidal. Denote by $\Pic_R$ the Picard $\infty$-groupoid of $\Mod_R$. 
\end{definition}
\begin{remark}
	$\Sigma^nR$ is invertible, with inverse $\Sigma^{-n}R$. In particular, there is a map $\bZ\times\Line_R\to\Pic_R$. However, this map need not be neither injective (if $R$ is $n$-periodic), nor surjective (see \cite{hill2011tmf}). 
\end{remark}
As mentioned in, the Thom spectrum functor makes sense for every functor $f:X^{\op}\to\Mod_R$. However, all examples of twists encountered so far came from functors into $\Line_R$. An example of twist that is not the result of a $R$-line bundle is the \emph{degree shift}. 
\begin{definition}
	Denote by $M$ the \emph{Thom $R$-module spectrum} functor \[\Grpd^{\op}_\infty/\Mod_R\to\Mod_R\]sending a functor $f:X^{\op}\to\Mod_R$ to its colimit.  
\end{definition}
\begin{example}
	Let $f:X^{\op}\to\Line_R$ be a twist. Denote by $\Sigma^nf$ the composition of $f$ with the shift functor $\Sigma^n:\Line_R\to\Pic_R$. Since $\Sigma^n$ is an equivalence, it commutes with colimits, so $$M\Sigma^nf\simeq \Sigma^nMf$$If $f=R_X$, then $M\Sigma^nf\simeq \Sigma^nR\wedge\Sigma^\infty_+X$, so  $\Sigma^nf$-twisted $R$-cohomology and $R$-homology correspond to normal $R$-cohomology and $R$-homology with a degree shift by $n$. 
\end{example}

\section{Umkehr Map}
We now proceed to the construction of the umkehr map in twisted cohomology theories. Here we follow \cite{abg2018thom}. Let $R$ be a ring spectrum, denote by $D_R$ the duality of $\Mod_R$. Given a invertible $R$-twist $\alpha:X^{\op}\to\Pic_R$, denote by $D_R\alpha$ the post-composition of $\alpha^{\op}$ with $D_R$.

\begin{remark} Depending on which cohomological degrees we want, we sometimes use the following alternative definition of	twisted cohomology:
	\[ R^{\alpha}(X) =\pi_0\hom_{\Mod_R}(M(D_R\alpha), R) \]
 meaning we use the dual/inverse of $\alpha$. To justify the use of $D_R\alpha$, consider the following: Let $\beta=D_R\alpha$, then
 \[\hom_{\Mod_R}(M\beta,R)\simeq\hom_{\Mod_{R}(X)}(D_R\alpha,R_X)\simeq\hom_{\Mod_R(X)}(R_X,\alpha\otimes_{R_X}R_X)\simeq\hom_{\Mod_R(X)}(R_X,\alpha)\]The second equivalence is a consequence of the fact that, if $M$ be an invertible $R$-module, left tensoring with $M$ is left adjoint to tensoring with $D_RM$, and $D_RD_RM\simeq M$, see \cref{rmk:dualinver}. If we think of $\alpha$ as a bundle of invertible $R$-modules over $X$, then $\hom_{\Mod_R(X)}(R_X,\alpha)$ is the spectrum of global sections of $\alpha$, which aligns with the idea that {twisted cohomology is the homotopy groups of the global sections of a bundle of spectra}. 
\end{remark}
If $R=\bS$, we denote the duality $D_R$ by simply $D$. 
\begin{definition}
	The \emph{Spanier-Whitehead dual} of a space $X$ is the dual spectrum of $\Sigma^\infty_+X$ with respect to the sphere spectrum.
\end{definition}
\begin{remark}\label{rmk:monunit}
	Recall that $\Sp$ is a closed symmetric monoidal category, with $\Hom_{\Sp}(E,-)$ being the right adjoint to the functor given by smash product with $E$. Let $(\lC,\otimes,1)$ a closed monoidal category, where $F(X,-)$ is the right adjoint to the functor given by left tensoring with $X$, then \[\\hom_\lC(-,F(1,E))\simeq\hom_\lC(1\otimes-,E)\simeq\hom_\lC(-,E)\]Using Yoneda's lemma, we conclude that $E\simeq F(1,E)$, for all $E$. 
\end{remark}
\begin{example}
	By \cref{rmk:monunit}, the Spanier-Whitehead dual to $*$ is the sphere spectrum, since \[D*=\Hom(\Sigma^\infty_+*,\bS)=\Hom(\Sigma^\infty S^0,\bS)=\Hom(\bS,\bS)\simeq\bS\]
\end{example}
\begin{definition}
	Denote by $\phi$ the functor $\lS^{\op}\to\bS/\Sp$ from spaces to the category of spectra under $\bS$, mapping $X$ to the map of spectra \begin{center}
		\begin{tikzcd}
			\phi(X):\bS\simeq D*\arrow[r,"Dp"] & DX
		\end{tikzcd}
	\end{center}where $p$ is the terminal map $X\to *$.
\end{definition} 
\begin{definition}
	Denote by $\widehat{\Mfd}_\infty$ the topological groupoid of closed smooth manifolds and diffeomorphisms. The set of diffeomorphisms are topologized by the weak $C^\infty$ topology, see \cite{hirsch1994differentialtopology}. Let $\Mfd_\infty$ be homotopy coherent nerve of $\widehat{\Mfd}_\infty$. 
\end{definition}
Let $B$ be a connected compact space and $\pi\colon X \to B$ a continuous function such that, for every $b\in B$, the fiber $\pi^{-1}(b)=:X_b$ is equipped with the structure of a closed smooth manifold, which varies continuously in $B$, in the sense of being classified by a functor
\[f:B \to \Mfd_\infty\]
Given such a classifying map $f$, consider the following composition \begin{center}
	\begin{tikzcd}
		B^{\op}\arrow[r] & \Mfd_\infty^{\op} \arrow[r] & \lS^{\op} \arrow[r,"\phi"] & \bS/\Sp
	\end{tikzcd}
\end{center}where the middle functor is the one forgetting the smooth structure. The above composition is then an object of
\[ \Fun(B^{\op},\bS/\Sp) \simeq \bS_B/\Fun(B^{\op}, \Sp) \]i.e. it is a natural transformation \[\phi_{X/B}:\bS_B\to D_B(f)\]where $D_B$ denotes the Spanier-Whithead duality applied to (the opposite of) $f:B\to\lS$ point-wise. Let $p:B\to *$ be the terminal functor and $p_!:\Fun(B^{\op},\Sp)\to\Sp$ the left Kan extension along $p$, i.e. the colimit functor. Applying $p_!$ to $\phi_{X/B}$, we obtain a morphism of spectra \begin{equation}\label{eq:thom}
	\begin{tikzcd}
		\Sigma_+^\infty B\simeq p_!p^*\bS=p_!\bS_B \arrow[rr,"p_!(\phi_{X/B})"] & & p_!D_B(f)
	\end{tikzcd}
\end{equation}
Let $T_{X/B}\to X$ be the vector bundle of fiber-wise tangent vectors, classified by $X^{\op}\to BO(n)$, where $n$ is the manifold dimension of the fibers\footnote{$B$ is connected, so the fibers have constant dimension.}. Denote by $\alpha$ the composition \begin{equation}\label{eq:alpha}
	\begin{tikzcd}
		X^{\op}\arrow[r,"T_{X/B}"] & BO(n) \arrow[r] & \lS_* \arrow[r,"\Sigma^\infty"] & \Sp
	\end{tikzcd}
\end{equation}the middle arrow being the $n$-sphere $S^n$ with $O(n)$-action coming from the one-point compactification of the regular action on $\bR^n$. The above diagram takes values in the Picard $\infty$-groupoid of the sphere spectrum. 
\begin{theorem}\label{thm:pontry}
	$p_!D_B(f)\simeq M(-\alpha)=:X^{-T_{X/B}}$.
\end{theorem}
\begin{definition}
	The morphism 
	\begin{center}
		\begin{tikzcd}
		\PT(f):\Sigma^{\infty}_+B\arrow[r] & X^{-T_{X/B}}	
		\end{tikzcd}
	\end{center}
	is called \emph{Pontryagin-Thom transfer map}. 
\end{definition}

Consider now a general situation. Let $R$ be a ring spectrum and $\alpha:X^{\op}\to\Line_R$ a $R$-line bundle. 
\begin{definition}
	An \emph{orientation} of $M\alpha$ is lift 
	\begin{center}
		\begin{tikzcd}
			& *\arrow[d]\\
			X^{\op}\arrow[r,"\alpha"']\arrow[ru,dashed] & \Line_R
		\end{tikzcd}
	\end{center}Explicitly, it is an equivalence of $R_X$-modules $\alpha\to R_X$. 
\end{definition}
Applying $p_!$, we see that an equivalence $t:\alpha\to R_X$ induces an equivalence
\begin{center}
	\begin{tikzcd}
		M\alpha\arrow[r,"p_!(t)"] & R\wedge\Sigma^{\infty}_+X
	\end{tikzcd}
\end{center}In our case, $\alpha$ in \cref{eq:alpha} is not valued in $\Line_\bS$, but $\Sigma^{-n}\alpha$ is. An orientation of $\Sigma^{-n}\alpha$ induces an equivalence \begin{equation}\label{eq:thomiso}
	\begin{tikzcd}
		\Sigma^{-n}M\alpha\arrow[r,"\simeq"] & R\wedge\Sigma^{\infty}_+X
	\end{tikzcd}
\end{equation}
\begin{definition}
	The isomorphism in cohomology induced by \cref{eq:thomiso} is called \emph{Thom isomorphism}. 
\end{definition}
Finally, we have the following definition:
\begin{definition}
	Assuming $\Sigma^{-n}\alpha:X^{\op}\to\Line_R$ is orientable, the \emph{Umkehr map} is the map: 
	\begin{center}
		\begin{tikzcd}
			R^*(\Sigma^\infty_+X)\arrow[rr,"\text{Thom iso.}"] & & R^{*-n}(X^{-T_{X/B}}) \arrow[rr,"\PT(f)^{*-n}"] & & R^{*-n}(\Sigma^{\infty}_+B)
		\end{tikzcd}
	\end{center}
\end{definition}
\begin{remark}
	The name \emph{Umkehr} comes from the map going in the opposite direction to $R^*(\Sigma^\infty_+X)\to R^*(\Sigma^\infty_+B)$, given by pulling back along $\pi$. 
\end{remark}
We now introduce the \emph{twisted Umkehr map}. Given a twist $\beta\colon B^{\op} \to \Pic_R$, we can smash to get a map 
\begin{center}
	\begin{tikzcd}
		\beta\simeq\bS_B\wedge_B\beta \arrow[r] & D_B(f)\wedge_B\beta
	\end{tikzcd}
\end{center}
Applying the functor $p_!$ once again, we obtain a morphism \begin{center}
	\begin{tikzcd}
		M\beta \arrow[r] & p_!(D_B(f)\wedge_B\beta)\simeq X^{-T_f+\beta\pi}
	\end{tikzcd}
\end{center}
where $\beta\pi:X^{\op}\to\Pic_R$ is the composition of $\pi$ and the twist $\beta$.
\begin{definition}
	The \emph{twisted Pontryagin-Thom transfer map} is the morphism:
	\begin{center}
		\begin{tikzcd}
			\PT(f,\beta):M\beta\arrow[r] & X^{-T_f+\beta\pi} 
		\end{tikzcd}
	\end{center}The \emph{twisted Umkehr map} is the map of twisted cohomology groups induced by $\PT(f,\beta)$. 
\end{definition}
\begin{example}
	Assume that $X^{-T_{X/B}+\beta\pi}$ is orientable, i.e. the following diagram commutes \begin{center} 
		\begin{tikzcd}
			X \arrow[r, "-T_f"] \arrow[d] & \Pic_\bS \arrow[d] \\ 	B \arrow[r, "\alpha"] & \Pic_R 
		\end{tikzcd}
	\end{center}then $X^{-T_{X/B}+\beta\pi}\simeq R\wedge\Sigma^\infty_+X$ and the twisted Umkehr map becomes:
	\begin{center}
		\begin{tikzcd}
			R^*(\Sigma^\infty_+X)\arrow[r] & R^{*-\beta}(B)
		\end{tikzcd}
	\end{center}
\end{example}

\chapter{Classical VS Modern Twisted K-Theory \texorpdfstring{\\}{ } \normalsize \emph{Talk by Alessandro Nanto}}

The goal of this talk is to compare different models of twisted $K$-theory. Concretely, we learn about the comparison between the original approach to twisted $K$-theory by Atiyah--Segal \cite{atiyahsegal2004twistedktheory} and the modern approach in \cite{abghr2014infty} that builds on the abstract approach to twisted cohomology via $\infty$-categorical methods. There is a third approach by Freed--Hopkins--Teleman \cite{freedhopkinsteleman2011loopgroupsi}, that we will not consider here.

\begin{remark}
	Throughout this talk we adopt the following notational conventions:
	\begin{enumerate}
		\item For an $\infty$-category $\lC$ and objects $X,Y \in \lC$, we denote by $\hom_{\lC}(X,Y)$ the mapping \emph{space} from $X$ to $Y$ in $\lC$. 
		\item If an $\infty$-category $\lC$ is enriched over spectra (e.g. if $\lC$ is stable), denote the \emph{mapping	spectrum} from $X$ to $Y$ by $\Hom_{\lC}(X,Y)$. 
  \item Note, in this case we have the relationship $\Omega^{\infty} \Hom_{\lC}(X,Y) \simeq \hom_{\lC}(X,Y)$, i.e., the underlying space of this mapping spectrum is then equivalent to the mapping space $\hom_{\lC}(X,Y)$.
	\end{enumerate}
\end{remark}

\section{Reviewing the \texorpdfstring{$\infty$}{oo}-categorical Approach to twisted Cohomology}
Let us recall the modern approach again. Given a ring spectrum	$R$, one can consider its space of invertible $R$-modules $\Line_R$. Now, given a space $X$, a twist of $R$ on $X$ is a map $\alpha \colon X \to \Line_R$. Then we define twisted homology as
\[X^\alpha \coloneq M \alpha = \colim (X \xrightarrow{ \ \alpha \ } \Line_R \to \Mod_R) \]
where the last map is the inclusion of $\Line_R$ into $\Mod_R$. The $R$-cohomology groups twisted by $\alpha$ are then defined as the homotopy groups of the mapping spectrum $\Hom_R(X^\alpha,  R)$, which is the spectrum given by the mapping spaces of maps from $X^\alpha$ to $\Sigma^n R$	in $\Mod_R$.

\section{The \texorpdfstring{$\infty$}{oo}-categorical Approach to twisted Cohomology via Sections}
Eventually we want to compare this definition with the classical approach. In this section, as a first step, we translate the definition into a more suitable form via sections, which will then allow us to compare it to the classical approach. 

Recall that the objects in $\Line_R$ are invertible $R$-modules. Let $-\alpha$ denote the map 
\[ - \alpha\coloneq X \xrightarrow{\alpha} \Line_R \xrightarrow{\mathrm{inv}} \Line_R, \]

We now have the following simple lemma. 

\begin{lemma}
	Let $M$ be an invertible	$R$-module. Then there is a natural equivalence of mapping spectra
	\[ \Hom(-M,R) \simeq \Hom(R,M) \] 
 This in particular means that the inverse of $M$ is given by $\Hom(M,R)$.
\end{lemma}

\begin{proof}
 We know that $\Hom(M,-)$ is the right adjoint to $M \otimes_R -$. But $M^{-1} \otimes_R -$ is also a right adjoint to $M \otimes_R -$, so they must be equivalent, meaning $M^{-1} \otimes R \simeq \Hom(M,R)$.  
\end{proof}

Based on this lemma we have the following chain of equivalences 
\[ \Hom_R(X^{-\alpha},R) \simeq \Hom(-\alpha,R_X) \simeq \Hom(R_X,\alpha) \]
Here in the last step we used the previous lemma point-wise. This proves the following proposition.

\begin{proposition}
	 The twisted $R$-cohomology spectrum twisted by $-\alpha$ is given by the mapping spectrum $\Hom(R_X,\alpha)$.
\end{proposition}

Let us now recall that every ring spectrum $R$ comes with a unique ring map $\bS \to R$, which induces an adjunction between $\bS$-modules (i.e. spectra) and $R$-modules:
\[
\begin{tikzcd}[column sep=2cm]
	\Mod_{\bS} \arrow[r, shift left=1.8, "R \otimes_{\bS} -", "\bot"'] & \Mod_R \arrow[l, shift left=1.8, "\text{Forget}"]
\end{tikzcd}.
\]
If we apply this adjunction point-wise to $\alpha$ we immediately get the following lemma.
\begin{lemma}
	There is an equivalence of mapping spectra \[ \Hom_R(R_X,\alpha) \simeq \Hom(\bS_X,\alpha). \]
\end{lemma}

Finally, let $\Omega^{\infty-n} \alpha \colon X \to \Line_R \to \lS$ be the composition of $\alpha$ with $\Omega^\infty\colon \Line_R \to \Sp \to \lS$. As $\bS \simeq \Sigma^\infty_+ *$, the adjunction between $\Sigma^\infty_+$	and $\Omega^\infty$ gives us the following equivalence of spaces:
\[\hom_{\Sp}(\bS_X,\alpha) \simeq \hom_{\lS}(*_X,\Omega^\infty\alpha).\]
Here $*_X$ is the constant functor at the point. The data of the right hand side is exactly a section. Hence combining these equivalences we get the following theorem.

\begin{theorem}
 A point in the underlying space of the twisted $R$-cohomology spectrum $\Omega^\infty\Hom_R(X^{-\alpha},R)$ is equivalent to a section of the bundle over $X$ classified by the map $\Omega^\infty \alpha \colon X \to \lS$.
\end{theorem}

Moreover, by similar analysis we have the following more general theorem.

\begin{theorem} \label{thm:twisted_R_cohomology_section}
 A point in the $n$-th underlying space of the twisted $R$-cohomology spectrum $\Omega^{\infty-n}\Hom_R(X^{-\alpha},R)$ is equivalent to a section of the bundle over $X$ classified by the map $\Omega^{\infty-n} \alpha \colon X \to \lS$.
\end{theorem}

\section{\texorpdfstring{$K$}{K}-Theory \texorpdfstring{{\` a}}{a} la Atiyah}
Before we can comprehend how Atiyah--Segal defined twisted $K$-theory, we first need to review how Atiyah defined classical $K$-theory via Fredholm operators in \cite{atiyah1967ktheory} (also in \cite{janich1965fredholm}). 

\begin{remark}
	Recall that every infinite-dimensional separable Hilbert space is isomorphic to each other. We will hence fix one such choice $\lH$ throughout these two sections.
\end{remark}

% , and fix a space $X$. Let $PU(\lH)$ be the projective unitary group of $\lH$, i.e. the quotient of the unitary group $U(\lH)$ by its center $S^1$. Let $P \to X$ be a principal $PU(\lH)$-bundle over $X$. 

We now recall the notion of a Fredholm operator.

\begin{definition}
	A Fredholm operator on $\lH$ is a bounded linear operator $F \colon \lH \to \lH$ such that both its kernel and cokernel are finite-dimensional. The space of Fredholm operators is denoted $\Fred(\lH)$ and is topologized as a subspace of the space of bounded linear operators with the norm topology.
\end{definition}

\begin{remark}
Fredholm operators admit many alternative characterizations. For example, it is equivalent to being invertible modulo compact operators i.e. 
\[\Fred(\lH) = \{A \in \lB(\lH) : AB - 1, BA - 1 \in \mathscr{K} \text{ for some } B \in \lB(\lH)\} \]
\end{remark}

We now define the spectrum $K_n$, given level-wise by the space of Fredholm operators on $\lH$:
\[ K_n = \Fred^{(n)}(\lH)\]

Here $\Fred^{(0)}(\lH)$ is simply $\Fred(\lH)$, and the other can be characterized similarly.

We will claim the following result regarding these spaces that we shall not prove.

\begin{lemma}
 The spaces $K_n$ assemble into a spectrum, meaning there are maps
	\[ K^{(n)}(\lH) \to \Omega \Fred^{(n+1)}(\lH). \]
\end{lemma}

How does this construction relate to $K$-theory? We have the following remark that can give us some intuition.

\begin{remark} \label{rem:Fredholm_Ktheory_map}
 For a compact space $X$, and given map $X \to \Fred^{(0)}(\lH)$, we can associate the formal difference of vector bundles given by the kernel and cokernel of the associated family of Fredholm operators (which are by assumption finite-dimensional). This associated to every element in $[X,\Fred^{(0)}(\lH)]$ an element in $K^0(X)$. 
\end{remark}

This map is not just coincidental and in fact we have the following major result.

\begin{theorem}[Atiyah--{J\"anich}]
 The map from \cref{rem:Fredholm_Ktheory_map} is an isomorphism.
\end{theorem}

This relates classical $K$-theory to Fredholm operators, meaning the spectrum above is indeed coincides with $K$-theory.

\section{Atiyah--Segal's Approach}
We now proceed to review the original approach to twisted $K$-theory by Atiyah--Segal \cite{atiyahsegal2004twistedktheory}. This approach generalizes the Fredholm approach to $K$-theory from the regular to the twisted setting. 

\begin{definition} \label{def:twisted_K_theory_atiyah}
 Fix a topological space $X$. Let $PU(\lH)$ be the projective unitary group of $\lH$, i.e. the quotient of the unitary group $U(\lH)$ by its center $S^1$. Let $P \to X$ be a principal $PU(\lH)$-bundle over $X$. Then we define the $\tau$-twisted $K$-theory group as the space of sections
	\[K^0_\tau \coloneq \pi_0(\Gamma(X,P \times_{PU(\lH)} K)). \]
\end{definition}

\begin{remark}
	Here, $P \times_{PU(\lH)} KU_0(\lH)$ denotes the associated bundle of spectra over $X$ with fibre $KU_0(\lH)$. The associated bundle can be described more explicitly as the bundle whose fibre over $x \in X$ is given by the $P_x \times_{PU(\lH)} K_0$, meaning we take the products of the fibers and quotient out the simultaneous action of $PU(\lH)$ on the fibers $P_x$ and $K_0$.
\end{remark}

\begin{remark} \label{rem:classifying_map}
	As we have $BPU(\lH) \simeq K(\bZ,3)$, the bundle $P \to X$ is classified via a map $\tau \colon X \to BPU(\lH) \simeq K(\bZ,3)$.
\end{remark}
 

\section{Comparing Classical and Modern}
We are now ready to compare the classical and modern approaches to twisted $K$-theory. Let $X$ be a topological space and $P \to X$ be a principal $PU(\lH)$-bundle over $X$. Following \cref{rem:classifying_map} this is classified by a map $\tau \colon X \to K(\bZ,3)$. 

For such a map we now have the following explicit commutative diagram 
\[ 
\begin{tikzcd}
 X \arrow[dr, "\tau"'] \arrow[r, "\tau"] & {K(\bZ,3)} \arrow[r, "?"] \arrow[d, "\mathrm{id}"] & BGL_1(KU) \arrow[r] \arrow[d] & \Line_{KU} \arrow[d, "\Omega^{\infty - n}"] \\
 & {K(\bZ,3)} \arrow[r] & BAut(K_n) \arrow[r] & \lS
\end{tikzcd},
\]
which we now explain in more detail. 
\begin{enumerate}
	\item The top row corresponds to the modern approach, as it takes a map into $K(\bZ,3)$ to a $KU$-line bundle. Following \cref{thm:twisted_R_cohomology_section}, taking $\Omega^{\infty - n}$ and then sections of the resulting bundle correspond precisely to classes of $\tau$-twisted $KU$-cohomology.
	\item The bottom row corresponds to the classical approach. Indeed, post-composing $\tau$ with the map $K(\bZ,3) \to BAut(K_n)$ corresponds to the constructing the associated bundle. In this case, it is by definition (\cref{def:twisted_K_theory_atiyah}) the case that sections of this bundle give us elements in twisted $K$-theory. 
	\item Hence compatibility between these two approaches corresponds to these diagrams commuting, which is a direct observation.
\end{enumerate}

% Here we use the associated bundle construction, because we need to transition from $BPU(\lH)$-bundles to $BGL_1(KU)$-bundles, and the associated bundle construction allows us to do this via the action of $PU(\lH)$ on $K$.

% This relates back to the modern approach, as both are now given via sections of bundles of spectra over $X$.

% We can summarize this via the following diagram 

% Here the top row corresponds to the modern approach, while the bottom row corresponds to the Atiyah--Segal approach.

\section{Some Comments regarding Freed--Hopkins--Teleman}
Ideally we could add some comments regarding the Freed--Hopkins--Teleman approach here, but that might turn out to be beyond existing capabilities.

\chapter{Wrong Way! Push-Forward in Twisted K-Theory \texorpdfstring{\\}{ } \normalsize \emph{Talk by Konrad Waldorf}}

Today we look at the construction of pushforward maps in twisted K-theory due to Carey and Wang \cite{careywang2008thomisomorphism}.

\section{The Classical Case}
Let us do a quick review of the classical case. Given a map $f\colon X \to Y$ and induced map on K-theory $f^*\colon K(Y) \to K(X)$, we would like to construct a pushforward map $f_!\colon K(X) \to K(Y)$. In the classical case, this requires $w_2(X)	= f^* w_2(Y)$. This precisely corresponds to choosing a $\spin^c$-structure on the virtual bundle $TX - f^* TY$. This is the necessary condition to define this map.

\section{Reviewing twisted K-Theory}
Before we proceed to generalize the pushforward map to twisted K-theory, we first review the definition of twisted K-theory. Recall that 
\[ H^3(M,\bZ) \cong PU(\lH) - \Bdl(M)/ \sim \] 
where $\lH$ is a separable Hilbert space and $PU(\lH)$ is the projective unitary group. Choosing $P$ a principal $PU(\lH)$-bundle over $M$ representing a class in $H^3(M,\bZ)$, we can form the associated bundle of Fredholm operators $\Fred(P) := P \times_{PU(\lH)} \Fred(\lH)$. Then the $0$-th twisted K-theory is defined as
\[K^0(M,P) := \pi_0(\mathrm{Map}(M, P\times_{PU(\lH)} \Fred(\lH))) \cong \pi_0(\mathrm{Map}(P, \Fred(\lH)))_{PU(\lH)}. \]
More generally 
\[ K^n(M,P) := \pi_0(\mathrm{Map}(P, \Omega^n \Fred(\lH)))_{PU(\lH)}. \]

In the coming sections we consider the pushforward map in two cases: the torsion case and the non-torsion case.

\section{Understanding the Torsion Case}
Let us consider the case when the bundle is torsion. Then $P$ is a principal $PU(n)$-bundle, using $U(n) \hookrightarrow	U(\lH)$. In this case we can obtain a bundle gerbe $\lL_P$, which has the following form 
 \[ 
	\begin{tikzcd}
	U(n) \arrow[d] & \Gamma_P \arrow[d] \arrow[l] & E \arrow[d] \\ 
	PU(n) & P \times_M P \arrow[l] \arrow[r, shift left = 1.8] \arrow[r, shift right = 1.8] & P \arrow[d] \\ 
	 & & M  
	\end{tikzcd}
	\]  
	This structure gives us a map 
	\[ 
	\phi\colon (\Gamma_P|_{p_1,p_2} \times_{U(1)} \bC) \otimes E_{p_2} \to E_{p_1}
	\]
	meaning $(E,\phi)$	is a $\lL_P$-twisted vector bundle.	We now can prove	the following theorem.

	\begin{theorem}
		Let $P$ be a principal $PU(n)$-bundle. Then, we have an isomorphism $K(\lL-\Mod) \cong K_{U(n), \mathrm{scal}}(P)$.
	\end{theorem}

	Here $K_{U(n), \mathrm{scal}}(P)$ is the $U(n)$-equivariant K-theory of $P$ with scalar action of the center $U(1) \subset U(n)$. The proof proceeds by mapping a pair $(E,\phi)$ to the data $(U(n) \times E \to E)$ that maps $(A,v) \mapsto \phi(A,v)$.

	We now have the following second theorem, relating twisted K-theory to twisted vector bundles.

	\begin{theorem}
	 Let $P$ be a principal $PU(n)$-bundle. Then we have an isomorphism
	 \[ K^0(M,P) \cong K(\lL_P-\Mod). \]
	\end{theorem}

	Here we are implicitly using the $PU(\lH)$-bundle $\tilde{P} = P \times_{PU(n)} PU(\lH)$, meaning 
	\[K^0(M,P) = K^0(M,\tilde{P}) = \pi_0(\mathrm{Map}(\tilde{P}, \Fred(\lH)))_{PU(\lH)}.\] 
Now we also have $\Fred(\lH)	\cong \Fred_{U(n)-cts}(\lH)$, meaning it is the Fredholm operators with continuous $U(n)$-action.

Now before we proceed, recall that we have a an equivariant version of Atiyah-J\"anich theorem (which holds if certain conditions are satisfied), meaning 
\[ [M,\Fred_{G-cts}(\lH)]_G \cong K^0(\Vect_G(M)). \]

We can now restrict this bijection to a bijection of subsets 
\[[M,\Fred_{U(n)-cts}(\lH^{\mathrm{scal}})]_G \cong K^0(\Vect_{U(n),\mathrm{scal}}(M)), \]
where $\lH^{\mathrm{scal}}$ is the Hilbert space with scalar action of the center $U(1) \subset U(n)$. 

The construction proceeds now by replacing $\lH$ with $\lH^{\mathrm{scal}}$ in the definition of twisted K-theory, meaning we go back to the first step and do 
	\[K^0(M,P) = K^0(M,\tilde{P}) = \pi_0(\mathrm{Map}(\tilde{P}, \Fred(\lH^{\mathrm{scal}})))_{PU(\lH)}.\]
	which via the bijection above is isomorphic to $K_{U(n),\mathrm{scal}}(P)$, which by the previous theorem is isomorphic to $K(\lL_P-\Mod)$. 

\begin{remark}
	Due to recent work, we know that if $P$ is torsion, then $\lL_P \cong \lA_P = P \times_{PU(n)}\bC^{n \times n}$, which is an isomorphism of $2$-vector bundles, and induces an isomorphism $\lL_P-\Mod \cong \lA_P-\Mod$.
\end{remark}

\section{Thom Isomorphism in the Torsion Case}
Let $V \to M$ be a $\bR$-vector bundle. Then we get $Tr(V) \to M$,	the associated $SO(n)$-bundle, which comes with a sub-bundle with a $\spin^c(n)$-structure. This gives us a bijection $\lL_{Tr(V)} \cong CL(V)$, where $CL(V)$ is the complex Clifford bundle of $V$.

Here we can use the result by Karoubi.

\begin{theorem}
	We have an isomorphism $K(CL(V)-\Mod) \cong K(\mathrm{Th}(V))$.
\end{theorem}

We now want to generalize this to more general	twists, which should be the following result.

\begin{theorem}
	There is an isomorphism
	\[K^0(M,P+W_3(V)) \cong K^0(\mathrm{Th}(V),\pi^*P). \]
\end{theorem}

\section{Non-Torsion Case}
We now aim to generalize these results to the non-torsion case. Concretely, we want prove an analogue of the following theorem:
\[K(\lL_P-\Mod) \cong K_{U(n),\mathrm{scal}}(P) \]
We can try to reproduce the same diagram 
\[ 
	\begin{tikzcd}
	 \Gamma_P \arrow[d] & E \arrow[d] \\ 
	 P \times_M P \arrow[r, shift left = 1.8] \arrow[r, shift right = 1.8] & P \arrow[d] \\ 
	 & M  
	\end{tikzcd},
	\]
	but we cannot compare to $U(n)$ anymore, since $P$ is not a $PU(n)$-bundle. So, we do not get the original result, but rather a restricted version.

	\begin{theorem}
	 $K^0(M,P) \cong K(\lL_P-\Mod^{U_2})$.
	\end{theorem}
	Here $U_2 \subseteq U(\lH)$	is the group of unitary objects that differ from the identity by a Hilbert-Schmidt operator, and $\lL_P-\Mod^{U_2}$ is the category of $\lL_P$-twisted vector bundles whose transition functions take values in $U_2$.

	Here $U^2$ can also be described via colimits, as $U^2 \cong \colim_{n \to \infty} U(n)$. So this result can be interpreted as the fact that filtered colimits preserve the original theorem.

Now we analogously have an isomorphism 
\[\lL_P-\Mod^{U_2} \cong \lA_p- \Mod^{U_2}. \]

\begin{remark}
	Similar to above, we anticipate an isomorphism	of $2$-vector bundles $\lL_P \cong P \times_{PU(\lH)} \mathrm{HS}(\lH)$, that induces the isomorphism above as a special case. This would require a suitable Morita category for these Hilbert-Schmidt operators, which is currently unknown.
\end{remark}

\section{Twisted Pushforward}
As a last step we can deduce the Thom	isomorphism in this non-torsion case. 
Let $f\colon X \to Y$. Then we have $W_3(f) \coloneq W_3(TX \oplus f^* TY)$. This gives us a pushforward map 
\[f_!\colon K^*(X,W_3(f) + f^*P) \to K^{(*  + (\dim X - \dim Y) \mod 2)}(Y,P), \]
for $P \to Y$ some $PU(\lH)$-bundle.

The construction of this map proceeds as follows: We can factor $X \to Y$ into $X \hookrightarrow Y \times S^n \to Y$. The first step uses the Thom isomorphism, and the second step uses $C^*$-algebra techniques.

\begin{remark}
 If $Y$ is trivial, then this construction gives us a version of an index. 
\end{remark}

\section{D-Branes}
	There is an application of these methods to physics. Given a $U(n)$-bundle $\lG$ over $X$ with sub objects $Q$, such that the restriction along $Q$ has a  $\lG|_Q$-module $\lE$.

	It is claimed by Witten that these d-branes are classified by some version of K-theory. This can be made precise as follows (with some adjustments). Given $\lE$ and $\lG|_Q \otimes \lL_f$-module we get an element in  $K^0(Q, \lG|_Q \otimes \lL_f)$. Using the pushforward map, we have a map
	\[K^0(Q, \lG|_Q \otimes \lL_f) \to K^*(X,\lG) \] 
giving us an element	in $K^*(X,\lG)$, as desired.

We can in particular consider the case when $X = \lG$, which is known as the WZW-model in physics. In this case, $Q$ are conjugacy classes $C_g$ of some element $g$ in the group and for $f\colon C_g \to \lG$, we have 
\[W_3(f) = (f^* \lG_{bas})^{\check{c}}\]  
	
This is a non-trivial observation. Indeed comparing with the Freed-Hopkins-Teleman approach \cite{freedhopkinsteleman2011loopgroupsi}, we see that 
\[K_{\lG}(\lG, \lG^k_{bas}) \cong \mathrm{Rep}^{k + \check{c}}(L\lG), \]
meaning this $\check{c}$ Coxeter shift appears naturally	in the representation theory of the loop group. 

Hence this approach suggests a possible geometric interpretation of these constructions.

\chapter{Non-Abelian Differential Cohomology \texorpdfstring{\\}{ } \normalsize \emph{Talk by Matthias Frerichs}}

Today we look at differential cohomology in the cohesive setting and in particular the generalization to non-abelian differential cohomology in the sense of Schreiber \cite{schreiber2013cohesive,fiorenzasatischreiber2024charactermap}.

We saw previously that differential cohomology lifts order cohomology theories represented by spectra. However, spectra can only classify phenomena that are inherently stable, whereas certain applications require unstable structures (see \cref{rem:motivation} for a more detailed discussion).

We hence discuss a ``non-abelian'' generalization of differential cohomology, that can incorporate such unstable phenomena. Concretely, instead of working with $E_\infty$-groups (which correspond to connective spectra), we want to work with $E_2$-groups, meaning homotopy commutative group objects in spaces.

\section{Non-Abelian Cohomology}
Let us first define non-abelian cohomology in the classical setting. HEre the idea is to generalize from spectra to spaces.

\begin{definition}
	Let $X, A$ be $\infty$-groupoids. The \emph{non-abelian cohomology} of $X$ with coefficients in $A$ is defined as
	\[	H^0(X; A) := \pi_0 \mathrm{Map}(X, A). \]
\end{definition}

We can generalize to higher cohomology if $A$ admits a delooping.

\begin{definition}
	Let $X$ be an $\infty$-groupoid and $A$ an	$E_2$-group with delooping $\mathrm{B}A$. The \emph{non-abelian cohomology} of $X$ with coefficients in $A$ is defined as
	\[	H^1(X; A) := \pi_0 \mathrm{Map}(X, \mathrm{B} A). \]
\end{definition}	

Of course, if $A$	admits further deloopings we can continue and obtain higher non-abelian cohomology groups.

\begin{definition}
	Let $X$ be an $\infty$-groupoid and $A$ admits deloopings $\mathrm{B}^n A$. The \emph{non-abelian cohomology} of $X$ with coefficients in $A$ is defined as
	\[	H^n(X; A) := \pi_0 \mathrm{Map}(X, \mathrm{B}^n A). \]
\end{definition}	

If $A$ admits infinite deloopings, then $A$ is precisely an $E_\infty$-group and hence corresponds to a connective spectrum. In that case we recover the classical definition.


\begin{remark}
	If $A$ is a spectrum, then we have $\mathrm{B}^n A = \Omega^{\infty - n} A$ and we recover the ordinary cohomology groups
	\[	H^n(X; A) = \pi_0 \mathrm{Map}(X, \Omega^{\infty - n} A) \cong [X, \Omega^{\infty - n} A] \cong A^n(X). \]
\end{remark}

Let us motivate this story.

\begin{remark} \label{rem:motivation}
	\begin{enumerate}
		\item The string bundle of a manifold $X$ is classified by homotopy classes of maps $X \to B \mathrm{String}$, which coincides with $H^1(X; \mathrm{String})$, originally defined by Giraud via non-abelian Cech cohomology \cite{giraud1971cohnonabelian}. Here $\mathrm{String}$ is an $E_1$-group (it is not abelian), so $B\mathrm{String}$ is just $E_2$.
		\item Another motivation comes from electro magnetics. Let $P$ be a $U(1)$-principal fiber bundle on $X$, then $[X,BU(1)] \cong [X, K(\bZ,2)] \cong [X,\bC\bP^\infty]$. We know that $\bC\bP^\infty$ is a colimit of $\bC\bP^n$, none of the intermediate stages are $\infty$-loop spaces. Generalizing the definition permits looking at $H(X,\bC\bP^n)$, which should have relevance in physics.
	\end{enumerate}
\end{remark}
	
\section{Differential Non-Abelian Cohomology}
We now want to generalize this approach to differential cohomology. Concretely, we want to look at differential refinements of non-abelian cohomology. We will do so by looking at differential refinements of intrinsic cohomology of a cohesive $\infty$-topos.

Naively, given an $\infty$-topos $\lG$ with cohesive structure and $X,A$ objects in $\lG$, we would like to define the non-abelian cohomology as 
\[ H^n(X; A) := \pi_0 \mathrm{Map}_{\lG}(X, \mathrm{B}^n A). \]
assuming $A$ admits a delooping $\mathrm{B}^n A$ in $\lG$.

Our goal is now to extract the fracture square, analogous to \cite{bunkenikolausvoelkl2016diffcoh}. 

\[
\begin{tikzcd}[row sep=1cm, column sep=2cm]
 \hat{E} \arrow[r] \arrow[d] & Z \arrow[d, "\text{homotopification}"] \\ 
	E \arrow[r, "\text{characteristic map}"] & \lH(Z)
\end{tikzcd} 
\]

How can we do this in the non-abelian case? We first review relevant definitions.

\begin{definition}
 A cohesive	$\infty$-topos $\lG$ is an $\infty$-topos together with an adjoint quadruple
	\[	\Pi \dashv \mathrm{Disc} \dashv \Gamma \dashv \mathrm{Codisc},	\]
	where $\Gamma\colon \lG \to \lS$ is the global sections functor,	such that
	\begin{enumerate}
		\item $\mathrm{Disc}$ and $\mathrm{Codisc}$ are fully faithful,
		\item $\Pi$ preserves finite products.
	\end{enumerate}
\end{definition}

\begin{example}
	Let $\Mfd$ be the category of smooth manifolds with the usual open cover topology. Then the $\infty$-topos of sheaves of $\infty$-groupoids on $\Mfd$ is cohesive, where
	\begin{enumerate}
		\item $\Pi$ is given by taking the fundamental $\infty$-groupoid,
		\item $\mathrm{Disc}$ is given by regarding an $\infty$-groupoid as a constant sheaf,
		\item $\Gamma$ is given by taking global sections,
		\item $\mathrm{Codisc}$ is given by regarding an $\infty$-groupoid as a codiscrete sheaf.
	\end{enumerate}
\end{example}

We can now define an analogue of the de Rham complex in the setting of a cohesive $\infty$-topos.

\begin{definition}
	Let $\lG$ be a cohesive $\infty$-topos and $A	\in \lG$. The \emph{de Rham complex} of $A$ is defined as
	\[
	\begin{tikzcd}
	 A \arrow[d] \arrow[r] & * \arrow[d] \\
		\mathrm{Disc} \circ \Pi(A) \arrow[r] &  \Pi_{dR}A
	\end{tikzcd}	
	\]
	and for an object $A$ the \emph{coefficient object} of $A$
	\[
	\begin{tikzcd}
		\flat_{dR}A \arrow[d] \arrow[r]	& * \arrow[d] \\
		\mathrm{Disc} \circ \Gamma(A) \arrow[r] & A 
	\end{tikzcd}
	\]
\end{definition}

Note here, there is an adjunction 
\[ 
\begin{tikzcd}
	\lG \arrow[r, shift left=2, "\Pi_{dR}", "\bot"'] & \lG \arrow[l, shift left=2, "\flat_{dR}"]
\end{tikzcd}
\]

With that at hand, we can formulate the following definition.

\begin{definition}
	For objects $X,A$ in a cohesive $\infty$-topos $\lG$ the \emph{differential non-abelian cohomology} of $X$ with coefficients in $A$ is defined as the set
	\[	H_{dR}(X, A) := H(\Pi_{dR}X,A) \simeq H(X, \flat_{dR} A). \]
\end{definition}

Notice this abstract definition recovers the classical case. 

\begin{proposition} \label{prop:cohomologycomputation}
 Let $X$ be a sheaf of manifolds. Then 
	\[ \mathrm{Map}(X,\flat_{dR}	\mathbf{B}^n U(1)) \simeq 
	\begin{cases}
	 H^{n}_{dR} & \text{for } n \geq 2, \\
		\Omega^1_{cl} & \text{for } n = 1, \\
		0 & \text{for } n = 0.
	\end{cases} 	
	\]
\end{proposition}

\begin{remark}
	Recall that in our original approach to differential cohomology we often used the pure sheaf as $\Omega^{\geq n}$. So, this definition does diverge slightly. However, it appears this difference does not play an adverse role.
\end{remark}

We will not prove this result, however, we	sketch the idea of the proof.
\begin{proof}[Idea]
	 Fundamentally the result follows from the following equivalence:
	\[\mathrm{Map}(X, \flat_{dR} B^n	U(1)) \simeq \mathrm{Map}(C(\{U_i\}), \Phi(\Omega^1(-) \to \Omega^2(-) \to ... \Omega^n_{cl})), \]
	where $C(\{U_i\})$ is the Cech nerve of a good open cover of $X$ and $\Phi$ is the Dold-Kan correspondence.

	This observation fundamentally relies on the computation that the sheaf associated to the chain complex
\[	\Phi(\Omega^1(-) \to \Omega^2(-) \to ...	\to \Omega^n_{cl}). \] 
 is given by $\flat_{dR} B^n	U(1)$.

	This allows unwinding the definition of cohomology in \cref{prop:cohomologycomputation} via $\flat_{dR}$ into a computation with differential forms, which recovers the classical de Rham cohomology groups.
\end{proof}

% Fundamentally this result relies on the computation, which shall remain without proof, of $\flat_{dR} B^n	U(1)$ as the sheaf associated to the chain complex
% \[	\Phi(\Omega^1(-) \to \Omega^2(-) \to ...	\to \Omega^n_{cl}). \]

We can now use this generalized definition to define curvature maps.

\begin{definition} \label{def:curvaturemap}
	Let $A$ be an object in $\lG$, such that $B^{n+1} A$ exists. The \emph{curvature map} is	defined as the map pullbacks
	\[ 
	\begin{tikzcd}
		B^nA \arrow[r] \arrow[d, "\mathrm{curv}"] & * \arrow[d]  \\
		\flat{dR} B^{n+1} A \arrow[r]  \arrow[d] & \mathrm{Disc} \circ \Gamma  B^{n+1} A \arrow[d] \\ 
		* \arrow[r] & B^{n+1} 
	\end{tikzcd}
	\]
\end{definition}

With this curvature map we can finally state the fracture square.
\begin{theorem}
	Let $A$ admit deloopings $B^n A$. Then we have a pullback square
	\[ 
		\begin{tikzcd}
			H_{diff}(X,B^nA) \arrow[r] \arrow[d] & H_{dR}(X,A) \arrow[d]  \\ 
			\mathrm{Map}_{sG}(X,B^nA) \arrow[r, "\mathrm{curv}_*"] & \mathrm{Map}_{sG}(X, \flat_{dR} B^{n+1} A)
		\end{tikzcd},
		\]
		where the bottom map is induced by the curvature map from \cref{def:curvaturemap}.
	\end{theorem}

The claim, which shall remain unproven, is that in the case of $A = U(1)$ this recovers the ordinary differential cohomology fracture square.

\chapter{(A Peek into) Twisted Differential Cohomology \texorpdfstring{\\}{ } \normalsize \emph{Talk by Alessandro Nanto}}

Today we want discuss twisted differential cohomology theory following \cite{bunkegepner2021diffktheory}.

\section{Reviewing Twisted Cohomology}
Let us first recall the definition of twisted cohomology theories, as discussed in a previous talk \cite{berkenhagen2025diffcoh8}. Let $X$ be a topological spaces, $R$ be a commutative ring spectrum. We denote by $\Mod_R$ the $\infty$-category of $R$-module spectra. The Picard $\infty$-groupoid of $R$ is defined as 
\[\Pic_R := \mathrm{Pic}(\Mod_R) \subset \Mod_R\]
the full sub- $\infty$-groupoid of invertible $R$-modules and equivalences.

\begin{definition}
A \emph{twist} of $R$ over $X$ is a map of spaces ($\infty$-groupoids)
\[\alpha : X \to \Pic_R.\]
\end{definition}

\begin{definition}
Given a twist $\alpha : X \to \Pic_R$, the \emph{$\alpha$-twisted $R$-cohomology of $X$} is defined as
\[R^\alpha(X) := \mathrm{Map}_{\Mod_R}(R, \alpha(X)).\]
\end{definition}

We can now compute the twisted cohomology groups as homotopy groups of the mapping spectrum.
\[R^{k + \alpha} = \pi_k(R^\alpha(X)) = \pi_k\left(\mathrm{Map}_{\Mod_R}(R, \alpha(X))\right).\]
This coincides with the $n$-homotopy group of $\Gamma(X, \alpha)$ where $\Gamma(X, \alpha)$ is the spectrum of sections of the bundle of spectra over $X$ associated to the twist $\alpha$.

\section{Reviewing Differential Cohomology}
We now want to recall the definition of differential cohomology theories. Here we don't follow the more modern approach via pure and $\bR$-invariant sheaves of spectra, as studied in \cite{bunkenikolausvoelkl2016diffcoh}, but rather the more classical approach \cite{hopkinssinger2005diffcoh}. This approach is also the one used in previous talks to study differential $K$-theory \cite{ludewig2025diffcoh7}.

We first review some notations and definition. Let $\Ch(\bR)$ be the $1$-category of chain complexes of real vector spaces. We denote by $\lD(\bR)$ the derived $\infty$-category obtained from $\Ch(\bR)$ by inverting quasi-isomorphisms. Finally, let $C$ be an object in an $\infty$-category $\lC$. Then $\underline{C}$ denotes the constant sheaf with value $C$.

\begin{definition}
  Let $X$ be a spectrum. A \emph{differential refinement of $X$} is a triple $(V, c\colon R \wedge X \to HV)$, where 
  \begin{enumerate}
    \item $V$ is a chain complex of real vector spaces,
    \item $HV$ is the Eilenberg-MacLane spectrum associated to $V$, via the stable Dold-Kan correspondence,
    \item $\alpha$ is an equivalence of spectra.
  \end{enumerate}
\end{definition}

Before we can proceed, we need to introduce further notation. Let $V$ be a chain complex of real vector spaces. Let $\Omega^*V$ be the sheaf of differential forms with values in $V$. 

\begin{definition}
 Let $\lM$ be a sheaf of chain complexes on the site of manifolds. The \emph{naive truncation} $\lM^{\geq n}$ is defined as $\lM$ if $k > n$ and $0$ otherwise.
\end{definition}

We can now get higher categorical sheaves out of sheaves of truncated chain complexes via localizations.

\begin{lemma}
  Let $\lM$ be a sheaf of $C^\infty$-modules, then post-composition with the localization functor $\Ch(\bR) \to \lD(\bR)$ preserves limits, and hence preserves the sheaf condition.
\end{lemma}

\begin{definition}
  Let $X$ be a $(X,V,c)$ be a differential refinement of a spectrum $X$, and $n \in \bZ$. Let $F^n(X,V,c)$, the \emph{differential function spectrum}, be defined as the pullback
  \[
  \begin{tikzcd}
    F^n(X,V,c) \arrow[rrr] \arrow[d] & & & H(\Omega^\bullet V^{\geq n}) \arrow[d, "c"] \\
    \underline{X} \arrow[r, "1 \wedge X"] & \underline{H\bR \wedge X} \arrow[r, "c"] &  \underline{HV} \arrow[r] & H(\Omega^\bullet V)
  \end{tikzcd}
  \]   
\end{definition}

Finally we can use the differential function spectrum to define differential cohomology groups.

\begin{definition}
  Let $X$ be a spectrum, $(X,V,c)$ be a differential refinement of $X$, and $n \in \bZ$. Then 
  \[ \hat{X}^n(M) \coloneq \pi_{-n}(F^n(X,V,c)(M))\]
  is the \emph{$n$-th differential $X$-cohomology group} of the manifold $M$.
\end{definition}

\section{Towards Twisted Differential Cohomology}
We can now combine the two previous definitions to define twisted differential cohomology theories. Let $R$ be a commutative ring spectrum, $M$ a manifold. We now want to define sheaf theoretic analogue of a twist, generalizing our previous definitions. Here we start using non-trivial definitions and results of \cite{bunkegepner2021diffktheory}.

\begin{definition}
  Let $\Mod_{\underline{R}}(M)$ be the $\infty$-category of sheaves of $\underline{R}$-module spectra on $M$ (i.e.~sheaves valued in $\Mod_R$). Define $\Pic^{loc}_{\underline{R}}(M)$ as the full sub-$\infty$-groupoid of $\Mod_{\underline{R}}(M)$ spanned by the locally constant $\underline{R}$-modules. The objects of $\Pic^{loc}_{\underline{R}}(M)$ are called \emph{topological $R$-twists} of $M$.
\end{definition}

This name is motivated by the following observation. There is an equivalence 
\[ \Fun(M^{top},\Pic_R) \to \Pic^{loc}_{\underline{R}}(M) \hookrightarrow \Shv((\Mfd_{/M})^{op},\Mod_R),\]
which sends $f$ to the sheaf that sends $g\colon N \to M$ to the $R$-module of global sections of $f\circ g$. Here $M^{top}$ is the underlying topological space of the manifold $M$. Here we also used the fact that $\Shv((\Mfd_{/M})^{op},\Mod_R)$ is equivalent to the $\Shv((\Mfd)^{op},\Mod_R)$ objects over the locally constant sheaf $\underline{M}$. So, the topological $R$-twists of $M$ are precisely the twists that come from a functor $M \to \Pic_R$, which is what a twist should really mean.


We now use this notion to define twisted differential cohomology. 

\begin{remark}
  Recall that the localization functor $\Ch(\bR) \Mod_{H\bR}$ is lax monoidal, meaning it preserves commutative algebra objects. 
\end{remark}

The previous remark implies that the composition of the stable Dold-Kan equivalence $\Ch(\bR) \to \lD(\bR)$ with the localization $\Ch(\bR) \Mod_{H\bR}$ is lax monoidal and hence preserves algebra objects. This means we get a functor
\[\CMon(\Ch(\bR)) \to \CMon(\Mod_{H\bR}) \]
Here the left hand side is the category of differential graded commutative algebras (CDGAs) over $\bR$, while the right hand side is the category of commutative $H\bR$-algebras.

\begin{definition}
  Let $R$ be a commutative ring spectrum, A \emph{differential (ring) refinement} of $R$ is a triple $(A, c, R)$ where
  \begin{enumerate}
    \item $A$ is a CDGA over $\bR$,
    \item $c\colon R \wedge H\bR \to HA$ is an equivalence of commutative $H\bR$-algebras.
  \end{enumerate}
\end{definition}

We can in fact make this definition more categorical.

\begin{definition}
  We define the $\infty$-category $\widehat{\mathrm{Ring}}$ of differential ring refinements as the pullback
  \[
  \begin{tikzcd}
    \widehat{\mathrm{Ring}} \arrow[r] \arrow[d] & \mathrm{DGCAlg} arrow[d, "H"] \\
    \CMon(\Sp) \arrow[r, " - \wedge H\bR"] & \CMon(\Mod_{H\bR})
  \end{tikzcd}
  \]
\end{definition}

Notice we can use the above to definition to again define a new sheaf of ring spectra via pullback.

\begin{definition}
  Let $(R,A,c)$ be a differential refinement of $R$. Then the \emph{differential function ring spectrum} $\hat{R}$ is defined as the pullback
  \[
  \begin{tikzcd}
    \hat{R} \arrow[d] \arrow[rrr] & & & H(\Omega^\bullet A)^{\geq n} \arrow[d] \\ 
    \underline{R} \arrow[r] & \underline{R \wedge H\bR} \arrow[r, "c"] & \underline{HA} \arrow[r] & H(\Omega^\bullet A)
  \end{tikzcd}
  \]
\end{definition}

We now want to define a notion of ``differential twists'' and ``differential module refinements'' of topological twists. For this we need to introduce some more notation and elaborate properties.

\begin{definition}
  Let $A$ be a CDGA over $\bR$. Let $\Pic^{wloc, fl}_A(M)$ be the full sub-$\infty$-groupoid of $\Mod_{\Omega^\bullet A}(M)$ spanned by the sheaves $\lM$ of $\Omega^\bullet A$-modules which are
  \begin{enumerate}
    \item \emph{weakly locally constant}, i.e. locally equivalent to a constant sheaf of $\Omega^\bullet A$-modules, as a sheaf valued in $\lD(\bR)$,
    \item \emph{$K$-flat}, i.e. such that the functor $\lM \otimes_{\Omega^\bullet A} -$ preserves quasi-isomorphisms,
    \item \emph{invertible}, i.e. there exists $\lN$ such that $\lM \otimes_{\Omega^\bullet A} \lN \cong \Omega^\bullet A$ as sheaves valued in $\Ch(\bR)$.
  \end{enumerate} 
\end{definition}

We are now ready to give the main definition of twisted differential cohomology.

\begin{definition}
  Let $E$ be a topological twist in $\Pic^{loc}_{\underline{R}}(M)$, a \emph{differential module refinement of $E$} is a triple $(E, \lM, c)$ where
  \begin{enumerate}
    \item $\lM \in \Pic^{wloc, fl}_A(M)$,
    \item $c\colon E \wedge H\bR \to H\lM$ is an equivalence in $\Pic^{loc}_{\underline{HA}}(M)$.
  \end{enumerate}
\end{definition}

\begin{definition}
  Let $(E, \lM, c)$ be a differential module refinement of a topological twist $E \in \Pic^{loc}_{\underline{R}}(M)$. Let $F(E,\lM,c)$ be the pullback
  \[
  \begin{tikzcd}
    F(E,\lM,c) \arrow[rr] \arrow[d] & & H(\lM^{\geq 0}) \arrow[d] \\
    E \arrow[r] & \underline{H\bR} \wedge E \arrow[r, "c"] &  \underline{H\lM} 
  \end{tikzcd}
  \]
  Denote by $\hat{R}^{(E,\lM,c)}(M) := \pi_0(F(E,\lM,c)(M))$ the \emph{twisted differential $R$-cohomology group} of $M$ associated to the differential refinement $(E,\lM,c)$ of the topological twist $E$.
\end{definition}

\begin{definition}
  Let $\mathrm{Tw}_{\hat{R}}(M)$, the \emph{$\infty$-category of differential twists}, be the $\infty$-category defined as the pullback
  \[
  \begin{tikzcd}
    \mathrm{Tw}_{\hat{R}}(M) \arrow[r] \arrow[d] & \Pic^{wloc, fl}_{\Omega^\bullet A}(M) \arrow[d, "H"] \\
    \Pic^{loc}_{\underline{R}}(M) \arrow[r, " - \wedge H\bR"] & \Pic^{loc}_{\underline{HA}}(M)
  \end{tikzcd}
  \]
\end{definition}

As a next step we can explore applications and examples of this abstract setup, in particular in the context of differential twisted $K$-theory.

\section{Differential Twisted Refinements of \texorpdfstring{$K$}{K}-Theory}
% We now want to apply the previous definitions to the case of complex $K$-theory
This section is some experimental effort towards understanding twisted differential $K$-theory and how we get classes via bundle gerbes.

Recall that for $K$-theory, $GL_1K = BS^1 \times \bZ/2\bZ \times BSU$, as by definition we have the pullback diagram 
\[ 
\begin{tikzcd}
 GL_1K \arrow[r] \arrow[d] & \Omega^\infty KU \arrow[d] \\
 \pi_0(KU)^\times \arrow[r] & \pi_0(KU)
\end{tikzcd},
\]
where $\bZ \times BU \simeq \Omega^\infty K \simeq \Hom_{\lS}(*, K) \simeq \Hom_{\Mod K}(K,K)$.

So for a given bundle gerbe, meaning a map $X \to K(\bZ,3)$, we get via post-composition with $K(\bZ,3) \to BGL_1K$, a twist of $K$-theory. So we have bijection between homotopy classes of maps from $X$ into $\tau_3 BGL_1 K$ (the $3$-truncation of $BGL_1K$) and homotopy classes of super bundle gerbes on $X$
\[[X, \tau_3 BGL_1 K] \cong \pi_0 sBldGrb(X).\]

More generally, we also have a bijection between homotopy classes of maps from $X$ into $\tau_3\Pic_K$ (the $3$-truncation of $\Pic_K$) and homotopy classes of $2$-line bundles on $X$
\[[X, \tau_3\Pic_K] \cong \pi_0 2LineBdl(X).\]

We now have differential refinements of $K(\bZ,3)$ and $BGL_1K$ given by Deligne cohomology and the Hopkins-Singer model of differential $K$-theory, respectively. So we can ask whether the map $K(\bZ,3) \to BGL_1K$ admits a differential refinement. Recall that this model of differential $K$-theory is given by the pullback
\[ 
\begin{tikzcd}
 \hat{K} \arrow[r] \arrow[d] & H\Omega^*(-, \bR[u,u^{-1}]) \arrow[d] \\
 \underline{K} \arrow[r] & \underline{H\bR[u,u^{-1}]} 
\end{tikzcd},
\]
where $|u| = -2$.

Similarly, Deligne cohomology is given by the pullback
\[
\begin{tikzcd}
 \bZ(3) \arrow[r] \arrow[d] & H\Omega^*(-)^{\geq 3} \arrow[d] \\
 \underline{H\bZ} \arrow[r] & \underline{H\bR} 
\end{tikzcd}.
\]

Now given a commutative ring spectrum $R$, we can consider the corresponding differential refinements, as a DGCA $A$ and an equivalence $c\colon R \wedge H\bR \to HA$. This gives us a diagram 
\[ 
\begin{tikzcd}
 \widehat{R(n)} \arrow[r] \arrow[d] & H\Omega^*(-,A)^{\geq n} \arrow[d] \\ 
 \underline{R} \arrow[r] & \underline{HA} 
\end{tikzcd}
\]
and we have 
\[
\hat{R}^n(M) = \pi_{-n}(\widehat{R(n)}(M)).
\]

A twist of $\hat{R} = (R,A,c)$ over $X$ consists of
\begin{enumerate}
  \item a topological twist $E\Sing X \to  \Pic_R$,
  \item a sheaf $\lM$ of $\Pic^{wloc,flat, ...}_{\Omega^*A}|_X \Omega^*(-,A)|_X$-modules on $X$,
  \item an equivalence $\tau\colon\underline{E} \wedge H A \to Hi\lM$ in $\Pic^{loc}_{\underline{HA}}(X)$. Here $i\colon \Ch(\bR) \to \lD(\bR)$ is the localization functor. Here $\underline{E}$ is the locally constant sheaf of $R$-modules associated to $E$ (i.e., the sheafification of the constant sheaf).
\end{enumerate}

Let us now see how we can go from twisted Deligne cohomology to twisted differential $K$-theory. 
\begin{example}
  Let $\omega \in \Omega^3(M)$, and take $(\Omega^*(-,\bR[u,u^{-1}]), d + u\omega)$. Now given a bundle gerbe with connection, we get a degree $3$-form $\omega$ as curvature, and the bundle gerbe gives us the map $E\colon M \to BGL_1K$. From the perspective of the definition 
  \begin{enumerate}
    \item $\hat{R} = \hat{K} = (K,\bR[u,u^{-1}], \text{chern character})$
    \item $E$ is obtained from the bundle gerbe, as explained above,
    \item $\lM$ is the sheaf of $\Omega^*(-,\bR[u,u^{-1}])$-modules given by $(\Omega^*(-,\bR[u,u^{-1}]), d + u\omega)$, and $\tau$ is the equivalence given by the Chern character, which is an equivalence of sheaves of $H\bR$-modules, as the Chern character is an equivalence after tensoring with $\bR$.
  \end{enumerate}
 Note here we have the computation on $(\Omega^*A, d + \lambda)$
 \[(d + u \omega)(a x) = da \cdot x \pm a \cdot dx + u\omega a x = d a \cdot x \pm a (d + u\omega) x \]
\end{example}

{\footnotesize
\bibliographystyle{alpha}
\bibliography{main}
}

\end{document}